\chapter{Composite Theta Lifts and Theorem A}
\label{chap:upperbound}

\begin{chapterguide}
\statementlink{thm:short-forms}{Theorem~A(a)} gives a spanning family of short integral forms for the ramified source space. It bounds the squared norms of the untwisted lifts of their normalized images under every complex automorphism. \statementlink{thm:index-bound}{Theorem~A(b)} bounds the local lattice indices of the normalized images of integral source forms. Both estimates use the local gains in~\eqref{eq:a-local-gain} and the rational normalizing factor $c_B$ in~\eqref{eq:normalized-poincare}. The seesaw identity of \secref{sec:a-seesaw} transfers the upper-chain estimates for the composite lift to the lower chain.
\end{chapterguide}

\section{Short integral forms and archimedean norms}
\label{sec:a-short-basis}

Let $U$ be the space of \emph{ramified source forms} and $W$ the space of \emph{good target forms}. Put $r=\dim_{\mathbb C}W$. Let $G$ be the product of general linear groups $\mathrm{GL}_{r_0}(\mathbf k_w)$ at the selected finite places $w\mid\pi_M$ above $p_0,p_1$.

For $g\in G$ write $P_g\colon U\to W$ for the \emph{Poincaré map} of~\eqref{eq:normalized-poincare}. Its construction applies the Poincaré series map, right translation by $g$, and projection to the good target space. Write $\Theta_1$ for the lower-chain theta lift for the trivial twist and $\|z\|^2_\sigma$ for the squared norm of $z$ under the complex automorphism $\sigma$. Write $I\subseteq U$ for the source forms integral at every $p$-adic place.

Write $\chi_K=\chi_d$ for the quadratic character of $K=\mathbb Q(\sqrt{-d})$. Let $I_p$ consist of the positive odd integers if $\chi_d(p)=-1$, the positive integers if $p=d$, and no indices if $\chi_d(p)=1$. Take $v_p^*=1$ in case $p\ne d$ and $1/2$ in case $p=d$. The \emph{local gain} at the prime $p$ is
\begin{equation}\label{eq:a-local-gain}
\begin{aligned}
 s_p(y,\tau)&=v_p^*\sum_{e\in I_p}\bigl(y/p^e-\kappa_{\nu_{p^e}\tau}\bigr)_+,\\
 \mu_p&=\sum_{\tau\in\Phi}s_p(\kappa_\tau,\tau),\qquad
 \gamma_p=m\left(\sum_{\tau\in\Phi_a}s_p(a,\tau)+k\mu_p\right).
\end{aligned}
\end{equation}
For $S'\subseteq\Phi$ write $C_p(S')$ for the \emph{collision sum} among unmarked places, defined over the ordered pairs $(\tau,j)$ by
\[
 C_p(S')=\sum_{\substack{\tau\in\Phi_a\setminus S',\ j\in I_p\\
                         \nu_{p^j}\tau\in\Phi_a\setminus S'}}\frac a{p^j}.
\]

Retain the notation of \statementref{Definition}{def:twisted-lattice}. For a finite family $y$, the condition $\mathcal I_\delta(y)$ requires nonnegative local losses $j_p$ such that $\sum_{p\in J}j_p\log p\le\delta$ for every finite set $J$ of primes. At each $p$-adic place above $p$, an integral matrix with determinant valuation at most $j_p$ must send $y$ into the lattice defined by all admissible twists simultaneously. The normalization is $v_p(p)=1$. The covolume upper bound $\mathcal V_V(y)$ requires the Gram determinant of the conjugated coordinates of $\Theta_1(y)$, using the original rational pairing matrix, to have absolute value at most $e^V$ under every complex automorphism.

\thmsource{Lean source: \leansourcelink{src/Lib/Adelic/Reduction/FirstMinimumShortBases.lean}{887}.}
\begin{lemma}[A short basis for a definite unitary lattice]
\label{lem:a-unitary-lattice-basis}
\emitthmsource
Let $E$ be a CM field with $f_E$ infinite places, and let $g$ be a finite adelic point of the standard definite unitary group of rank $n\ge1$. Let the positive integer $\Delta_0$ belong to the different of $E/\mathbb Q$. Then there exists $Y\in\mathrm{GL}_n(E)$ such that $g^{-1}Y$ is a finite adelic integral matrix and every entry of $Y$ at each complex embedding has absolute value at most $6\Delta_0(f_E n)^2$.
\end{lemma}
\begin{proof}
Write $L=E^n\cap g\widehat{\mathcal O}_E^{\,n}$ and let $\iota_\infty$ be the real coordinate map obtained by choosing one complex embedding at each infinite place. Then $\Lambda=\iota_\infty(L)$ is a full lattice in real dimension $d_{\mathbb R}=2f_E n$. Write $h$ for the standard \emph{Hermitian form}. Unitarity of $g$ makes $g\widehat{\mathcal O}_E^{\,n}$ self-dual for the finite adelic $h$-pairing, where the integral dual consists of the vectors whose pairings against the whole lattice lie in $\widehat{\mathcal O}_E$. Density of $L$ in the finite completion then expresses self-duality as
\[
 \{v\in E^n:h(v,L)\subseteq\mathcal O_E\}=L.
\]
For any nonzero vector $v\in L$, the value $h(v,v)$ is a totally positive algebraic integer in the maximal real subfield of $E$. If $\|\iota_\infty(v)\|<1$, the images of $h(v,v)$ under all real embeddings of that subfield would be less than $1$, though their product is a positive integer. Thus each nonzero vector of $\Lambda$ has length at least $1$.\leanfn{See \href{https://github.com/lsernau/grh-lean/blob/72fead102f663d5c7a407964a13dc0cf9d977207/grh/src/Lib/Adelic/Reduction/FirstMinimumShortBases.lean\#L504}{The integral Hermitian pairing on the lattice}, \href{https://github.com/lsernau/grh-lean/blob/72fead102f663d5c7a407964a13dc0cf9d977207/grh/src/Lib/Adelic/Reduction/FirstMinimumShortBases.lean\#L560}{Lower bound for the length of a nonzero lattice vector}.}

Put $L^\vee=\{v\in E^n:\operatorname{Tr}_{E/\mathbb Q}h(v,L)\subseteq\mathbb Z\}$. Since $L$ is an $\mathcal O_E$-module, integrality of the trace pairing against $L$ implies $h(v,L)$ belongs to the inverse different $\mathfrak D_{E/\mathbb Q}^{-1}$. Multiplication by $\Delta_0$ then makes all pairings integral, and the preceding identity gives $\Delta_0 L^\vee\subseteq L$. At infinity we have $\operatorname{Tr}_{E/\mathbb Q}h(v,l)=2\langle\iota_\infty(v),\iota_\infty(l)\rangle_{\mathbb R}$. This trace pairing is nondegenerate on $\mathbb Q$, so the trace-dual basis of any $\mathbb Z$-basis of $L$ still lies in $E^n$. Hence the Euclidean dual of the scaled lattice $\Lambda'=(2/\Delta_0)\Lambda$ is
\[
 (\Lambda')^*=\iota_\infty(\Delta_0 L^\vee)\subseteq \Lambda.
\]
Its shortest nonzero vector also has length at least $1$.\leanfn{See \href{https://github.com/lsernau/grh-lean/blob/72fead102f663d5c7a407964a13dc0cf9d977207/grh/src/Lib/Adelic/Reduction/FirstMinimumShortBases.lean\#L674}{Trace duality and the scaled lattice}.}

The transference bound~\cite{LLS1990} between the covering radius and the first minimum of the dual lattice gives covering radius at most $r=d_{\mathbb R}^2$ for $\Lambda'$. The vectors in $\Lambda'$ of length at most $3r$ then span the real coordinate space. If their span were proper, take a vector $u$ of length $2r$ orthogonal to that span. By the covering radius there exists $\lambda\in\Lambda'$ with $\|\lambda-u\|\le r$. Thus $\|\lambda\|\le3r$, so $\lambda$ would lie in the span, contrary to the distance $2r$ of $u$ from it. Multiplying the short vectors by $\Delta_0/2$ gives a family in $L$ of lengths at most $3r\Delta_0/2=6\Delta_0(f_E n)^2$. Since their coordinates at infinity span the real space, they span $E^n$. Choose an $E$-basis from the family as the columns of $Y$. Then $g^{-1}Y$ is integral, and at each complex embedding the absolute values of the entries are at most the lengths of their columns.\leanfn{See \href{https://github.com/lsernau/grh-lean/blob/72fead102f663d5c7a407964a13dc0cf9d977207/grh/src/Lib/Adelic/Reduction/FirstMinimumShortBases.lean\#L150}{A weak lattice transference theorem}, \href{https://github.com/lsernau/grh-lean/blob/72fead102f663d5c7a407964a13dc0cf9d977207/grh/src/Lib/Adelic/Reduction/FirstMinimumShortBases.lean\#L809}{Constructing a short basis from the covering radius}.}
\end{proof}

\thmsource{Lean source: \leansourcelink{src/Lib/Shimura/HeckeTower/HeckeData.lean}{979}.}
\begin{lemma}[A stabilizer criterion for principal congruence levels]
\label{lem:a-free-congruence-level}
\emitthmsource
Let $E$ be a CM field with at least two infinite places, with $n\ge1$, and let $\Gamma$ be the rational point group of the standard definite unitary group of rank $n$. Let $K$ be a subgroup of its finite adelic integral points reducing to the identity at the distinct finite places in $\mathcal W$. If
\[
 \prod_{w\in\mathcal W}\operatorname{Nm}(w)>(2n)^{[E:\mathbb Q]},
\]
then for every finite adelic point $g$ we have $\Gamma\cap gKg^{-1}=\{1\}$.
\end{lemma}
\begin{proof}
If $\gamma\in\Gamma\cap gKg^{-1}$, then $k=g^{-1}\gamma g\in K$. Put $t=\operatorname{tr}(\gamma)-n\in E$. By conjugation invariance of the trace, $t$ is integral at every finite place and reduces to zero at $w\in\mathcal W$. Thus $t\in\mathcal O_E$ and the trace difference belongs to $\prod_{w\in\mathcal W} w$. If $t\ne0$, divisibility of principal ideals gives
\[
 |\operatorname{Nm}_{E/\mathbb Q}(t)|
 \ge \prod_{w\in\mathcal W}\operatorname{Nm}(w).
\]
On the other hand, $\gamma$ is unitary at every complex embedding, so its eigenvalues have modulus $1$. Hence $|\tau(t)|\le2n$ and $|\operatorname{Nm}_{E/\mathbb Q}(t)|\le(2n)^{[E:\mathbb Q]}$, contradicting the hypotheses. Thus $t=0$. Consequently, at each complex embedding the real parts of the eigenvalues of $\gamma$ sum to $n$, and each is at most $1$. All eigenvalues equal $1$. Since a unitary matrix is diagonalizable, $\gamma=1$.\leanfn{See \href{https://github.com/lsernau/grh-lean/blob/72fead102f663d5c7a407964a13dc0cf9d977207/grh/src/Lib/Shimura/HeckeTower/HeckeData.lean\#L923}{A trace argument for stabilizers at principal congruence level}.}
\end{proof}

\thmsource{Lean source: \leansourcelink{src/New/UpperBound/Size/ShortColumns.lean}{362}.}
\begin{lemma}[Short integral forms at a torsion-free level]
\label{lem:a-free-level-short-forms}
\emitthmsource
Let $E$ be a CM field with $f_E$ infinite places, $G_f$ the finite adelic point group of the standard definite unitary group of rank $n\ge1$, $\Gamma$ its rational point group, and $K$ an integral level satisfying $\Gamma\cap gKg^{-1}=\{1\}$ for every $g\in G_f$. Fix Weyl weight $(\kappa,e)$ and write $s$ for the total number of boxes in columns of height less than $n$ at the infinite places. Let the positive integer $\Delta_0$ lie in the different of $E/\mathbb Q$ and put $C=6\Delta_0(f_E n)^2$. The algebraic forms of level $K$ are spanned by $C^{6(nf_E+s)}$-short forms. A form is $M$-short if it is integral at every $p$-adic place and the \emph{Weyl norm} at every $\operatorname{Aut}(\mathbb C)$-conjugate point value is at most $M$.
\end{lemma}
\begin{proof}
Fix representatives of the double-coset space $\Gamma\backslash G_f/K$. For a representative $g$ and a Weyl coordinate vector $v$, define the single-class form $\delta_{g,v}$ to be zero off $\Gamma gK$, to take value $v$ at $g$, and to extend by the rational action on the left and $K$-invariance on the right. Triviality of the stabilizers makes this well defined. Every form is the sum of these single-class forms for its values at the representatives. If $v$ has \emph{Weyl gauge} at most $1$ at the representative $g$, then moving $\delta_{g,v}(\gamma gk)$ back to the identity cancels the left $\gamma$ action with $\gamma^{-1}$. The remaining $k^{-1}$ preserves the integral lattice. The single-class form is therefore integral at that place. At infinity, the rational action preserves the Weyl norm under every complex embedding, so the pointwise bound follows from the bound for $v$ at the representative.\leanfn{See \href{https://github.com/lsernau/grh-lean/blob/72fead102f663d5c7a407964a13dc0cf9d977207/grh/src/Lib/Adelic/AutomorphicForms/HeckeEigenforms.lean\#L221}{Evaluation isomorphism at a torsion-free level}, \href{https://github.com/lsernau/grh-lean/blob/72fead102f663d5c7a407964a13dc0cf9d977207/grh/src/New/UpperBound/Size/ShortColumns.lean\#L330}{Integrality and norms of forms supported on a single class}.}

For each $h\in G_f$, \statementref{Lemma}{lem:a-unitary-lattice-basis} supplies $Y(h)\in\mathrm{GL}_n(E)$ such that $h^{-1}Y(h)$ is a finite adelic integral matrix and the entries of $Y(h)$ have absolute value at most $C$ under each complex embedding. Fix a representative $g$ and write $Y_0=Y(g)$ and $Y_\tau=Y(g^{e_\tau})$. Let $\rho_{\kappa,0}$ be the Weyl representation without the determinant powers, and use $\tau$ also for the chosen complex embedding. Put
\[
 A_g=\left(\prod_\tau\det\tau(Y_\tau)\right)
      \rho_{\kappa,0}\bigl((\tau(Y_0))_\tau\bigr).
\]
The operator $A_g$ is invertible. Take the coordinates of the array $y$ to be roots of unity, and let $\operatorname{ev}_\kappa(y)$ be the vector of the values at $y$ of the integral Weyl basis polynomials. Set $v_y=A_g\operatorname{ev}_\kappa(y)$. We show that the vectors $v_y$ span the Weyl coordinate space. If a linear functional vanished on all $\operatorname{ev}_\kappa(y)$, the corresponding finite linear combination of the basis polynomials would vanish whenever all variables are roots of unity. There are infinitely many roots of unity, so the univariate root bound applied to each variable shows the polynomial vanishes identically. Independence of the basis polynomials then makes the functional zero. Invertibility of $A_g$ gives the claim.\leanfn{See \href{https://github.com/lsernau/grh-lean/blob/72fead102f663d5c7a407964a13dc0cf9d977207/grh/src/Lib/Adelic/Reduction/FirstMinimumShortBases.lean\#L888}{A short basis for the finite adelic lattice}, \href{https://github.com/lsernau/grh-lean/blob/72fead102f663d5c7a407964a13dc0cf9d977207/grh/src/Lib/GLn/SchurModule/WeylKernel.lean\#L188}{Spanning by the value vectors}.}

At a $p$-adic place, move $v_y$ to the identity by $g^{-1}$ at weight $(\kappa,e)$. Its Weyl part becomes the Schur matrix given by the images of $g^{-1}Y_0$ at the embeddings, and the determinant part is the product of the determinants of the images of $(g^{e_\tau})^{-1}Y_\tau$. Both kinds of matrices are integral, and the entries of the Schur matrices are integral polynomials. Roots of unity have absolute value $1$ at every finite place, so each coordinate of $\operatorname{ev}_\kappa(y)$ is integral. Thus all coordinates of the transported vector are integral and the gauge of $v_y$ at $g$ is at most $1$.\leanfn{See \href{https://github.com/lsernau/grh-lean/blob/72fead102f663d5c7a407964a13dc0cf9d977207/grh/src/New/UpperBound/Size/ShortColumns.lean\#L186}{Local Weyl integrality}.}

Fix $\sigma\in\operatorname{Aut}(\mathbb C)$. The entries of $\sigma\tau(Y_0)$ and $\sigma\tau(Y_\tau)$ still have absolute value at most $C$, and the entries of $\sigma(y)$ still have absolute value $1$. Let $T_\tau$ be the square matrix of the first $n$ columns of $\sigma(y)$ at $\tau$. Using the value vector $u$ of the standard tableau array gives
\[
 \rho_{\kappa,0}\bigl((\sigma\tau(Y_0))_\tau\bigr)
       \operatorname{ev}_\kappa(\sigma(y))
 =\operatorname{Schur}_\kappa\bigl((\sigma\tau(Y_0)T_\tau)_\tau\bigr)u,
 \qquad \|u\|_{\mathrm{Weyl}}^2=1.
\]
The entries of each product matrix have absolute value at most $nC$, so its squared singular values are at most $n^4C^2\le C^6$. The Schur operator bound for the squared Weyl norm uses these as bases with the row lengths of the Young diagram as exponents, which sum to $s$. Hence the squared norm above is at most $C^{6s}$. Hadamard's inequality also gives
$|\det(\sigma\tau(Y_\tau))|\le(\sqrt n\,C)^n\le C^{3n}$, so the squared product of the $f_E$ determinants is at most $C^{6nf_E}$. Combining the two bounds gives $\|\sigma(v_y)\|^2_{\mathrm{Weyl}}\le C^{6(nf_E+s)}$.\leanfn{See \href{https://github.com/lsernau/grh-lean/blob/72fead102f663d5c7a407964a13dc0cf9d977207/grh/src/New/UpperBound/Size/ShortColumns.lean\#L275}{Archimedean bounds for Weyl vectors}, \href{https://github.com/lsernau/grh-lean/blob/72fead102f663d5c7a407964a13dc0cf9d977207/grh/src/Lib/GLn/SchurModule/SchurNorm.lean\#L197}{A bound for squared singular values of the Schur operator}.}

Thus each $\delta_{g,v_y}$ satisfies the requirements. Since the vectors $v_y$ span all coordinates at each representative, the short integral forms span the algebraic forms at that level.
\end{proof}

\thmsource{Lean source: \leansourcelink{src/New/UpperBound/Size/ShortColumns.lean}{976}.}
\begin{lemma}[Short integral forms in the full source space]
\label{lem:a-short-source-basis}
\emitthmsource
For admissible data $D$, the full source space $U_0$ is spanned by $e^B/4^f$-short forms.
\end{lemma}
\begin{proof}
The initial source space has rank $k$ and its level lies inside the chosen principal congruence level. This level reduces to the identity at $f$ places of residue degree $1$ above the chosen prime $p_0$, whose norms have product $p_0^f$. The scales give $p_0>(2k)^2$ and hence $p_0^f>(2k)^{[E:\mathbb Q]}$. \statementref{Lemma}{lem:a-free-congruence-level} gives trivial stabilizers at every point for the source level.\leanfn{See \href{https://github.com/lsernau/grh-lean/blob/72fead102f663d5c7a407964a13dc0cf9d977207/grh/src/New/UpperBound/Objects/LiftChains.lean\#L514}{Stabilizer bounds for the principal congruence level}.}

Take $\Delta_0=q_1d$. A generator of the relative different $\mathfrak D_{E/K}$ divides $q_1$, and the generator $\sqrt{-d}$ of $\mathfrak D_{K/\mathbb Q}$ divides $d$. Transitivity of the different gives $\Delta_0\in\mathfrak D_{E/\mathbb Q}$. In a Young diagram for the source space, every nonempty column of height less than $k$ has height $1$, with at most $a$ such columns per infinite place. Hence the total number $s$ of boxes is at most $fa$.\leanfn{See \href{https://github.com/lsernau/grh-lean/blob/72fead102f663d5c7a407964a13dc0cf9d977207/grh/src/New/Derivatives/Field/PrimeDecomposition.lean\#L253}{The different and the conductor}, \href{https://github.com/lsernau/grh-lean/blob/72fead102f663d5c7a407964a13dc0cf9d977207/grh/src/New/UpperBound/Size/ShortColumns.lean\#L806}{Counting boxes in the source weights}.}

Put $C=6(q_1d)(fk)^2$. \statementref{Lemma}{lem:a-free-level-short-forms} gives the squared-norm threshold $C^{6(kf+s)}$, and the scale estimates give
\[
 f\log 4+6(kf+fa)\log C\le B .
\]
Since $C\ge1$ and $s\le fa$, we get
$C^{6(kf+s)}\le C^{6(kf+fa)}\le e^B/4^f$. The spanning family therefore satisfies the required upper bound.\leanfn{See \href{https://github.com/lsernau/grh-lean/blob/72fead102f663d5c7a407964a13dc0cf9d977207/grh/src/New/UpperBound/Size/ShortColumns.lean\#L977}{Explicit constants for the full source space}.}
\end{proof}

\thmsource{Lean source: \leansourcelink{src/New/UpperBound/Size/ShortColumns.lean}{942}.}
\begin{lemma}[Dimension of the ramified source space]
\label{lem:a-ramified-dimension}
\emitthmsource
Let $U_0$ be the full source space, $U$ the ramified source space, and $N$ the level prime. Then
\[
 \dim_{\mathbb C}U\ge
 \left(1-\frac{f}{N-1}\right)\dim_{\mathbb C}U_0 .
\]
\end{lemma}
\begin{proof}
Let $K_0$ be the level of $X_0$, and $\pi=\pi_N$ the level integer with $\pi\bar\pi=N$. For $w\mid\pi$ denote the chosen scalar adelic point and its translation operator by $z_w$. Its residue at $w$ generates the multiplicative group of order $N-1$, and at the other places dividing $\pi$ it is $1$. For a positive integer $i$, the scalar matrix $z_w^i$ belongs to $K_0$ exactly when its bottom row satisfies the level condition at each place dividing $\pi$. This requires its residue at $w$ to be $1$, or equivalently $N-1\mid i$. Thus $1,z_w,\dots,z_w^{N-2}$ represent distinct $K_0$-cosets in $K_w=\langle K_0,z_w\rangle$.\leanfn{See \href{https://github.com/lsernau/grh-lean/blob/72fead102f663d5c7a407964a13dc0cf9d977207/grh/src/New/UpperBound/Size/ShortColumns.lean\#L632}{Local order of a scalar translation}.}

As $K_0$ and $z_w$ lie in the chosen principal congruence level, the rational stabilizers for $K_w$ and $K_0$ are trivial. Fix a double coset $\Gamma gK_w$ for the larger level $K_w$. The smaller-level double cosets $\Gamma gkK_0$ correspond bijectively to cosets $kK_0\in K_w/K_0$. In fact, $\Gamma gk_1K_0=\Gamma gk_2K_0$ gives $\gamma\in\Gamma$ and $t\in K_0$ with $\gamma gk_1=gk_2t$, so $\gamma\in\Gamma\cap gK_wg^{-1}=\{1\}$ and hence $k_1K_0=k_2K_0$. At a torsion-free level one can assign any value in the same Weyl module to each representative. Thus the dimension is the number of double cosets times the Weyl module dimension. The $K_0$-level forms fixed by $z_w$ are exactly the $K_w$-level forms, so\leanfn{See \href{https://github.com/lsernau/grh-lean/blob/72fead102f663d5c7a407964a13dc0cf9d977207/grh/src/New/UpperBound/Size/ShortColumns.lean\#L514}{Dimensions and indices for torsion-free levels}, \href{https://github.com/lsernau/grh-lean/blob/72fead102f663d5c7a407964a13dc0cf9d977207/grh/src/New/UpperBound/Size/ShortColumns.lean\#L819}{Dimension bound for the fixed subspace}.}
\[
 \dim_{\mathbb C}\ker(1-z_w)
 =\frac{\dim_{\mathbb C}U_0}{[K_w:K_0]}
 \le \frac{\dim_{\mathbb C}U_0}{N-1}.
\]

The $(N-1)$-th power of each $z_w$ is the identity on $U_0$. The operators $z_w$ commute, and $X^{N-1}-1$ has simple roots over $\mathbb C$, so $U_0$ splits into their simultaneous eigenspaces. Ramification requires the operator $\sum_{i=0}^{N-2}z_w^i$ to annihilate the vector for each $w$. On a common eigenspace the sum acts by $N-1$ when the eigenvalue of $z_w$ is $1$ and by zero otherwise. The ramified space is thus the sum of eigenspaces where each $z_w$ has eigenvalue different from $1$, or equivalently the image of $\prod_{w\mid\pi}(1-z_w)$. Every excluded eigenspace is contained in some $\ker(1-z_w)$. Since $N$ splits completely in $E$, the $2f$ finite places above $N$ occur in conjugate pairs, exactly one of each pair dividing $\pi$. Thus\leanfn{See \href{https://github.com/lsernau/grh-lean/blob/72fead102f663d5c7a407964a13dc0cf9d977207/grh/src/New/UpperBound/Size/ShortColumns.lean\#L712}{The ramification condition and the product operator}, \href{https://github.com/lsernau/grh-lean/blob/72fead102f663d5c7a407964a13dc0cf9d977207/grh/src/New/UpperBound/Objects/Columns.lean\#L601}{Counting places above the level prime}.}
\[
 \dim_{\mathbb C}U_0-\dim_{\mathbb C}U
 \le \sum_{w\mid\pi}\dim_{\mathbb C}\ker(1-z_w)
 \le \frac{f}{N-1}\dim_{\mathbb C}U_0 .
\]
\end{proof}

\thmsource{Lean source: \leansourcelink{src/New/UpperBound/Size/ShortColumns.lean}{954}.}
\begin{lemma}[A scale bound for the codimension of the ramified space]
\label{lem:a-ramified-loss-scale}
\emitthmsource
For admissible data $D$, we have $f/(N-1)\le b^{-8}$.
\end{lemma}
\begin{proof}
The scale conditions give $0<b$, $f\le b$, and $b^9+1\le N$. Hence $N-1\ge b^9>0$, so\leanfn{See \href{https://github.com/lsernau/grh-lean/blob/72fead102f663d5c7a407964a13dc0cf9d977207/grh/src/New/UpperBound/Size/ShortColumns.lean\#L954}{The level prime and scale conditions}.}
\[
 \frac{f}{N-1}\le\frac{b}{b^9}=b^{-8}.
\]
\end{proof}

\thmsource{Lean source: \leansourcelink{src/New/UpperBound/Size/ArchimedeanSize.lean}{426}.}
\begin{lemma}[Archimedean size of the composite lift]
\label{lem:a-composite-size}
\emitthmsource
Let $\Psi$ be \emph{chain datum} for the trivial twist, with $\sigma\in\operatorname{Aut}(\mathbb C)$ and $u\in\mathbb R$. If all $\sigma$-conjugate point values of the full source form $\phi$ have squared Weyl norm at most $e^u$, then for each finite adelic point $v$,
\[
 \bigl\|\sigma\bigl((\operatorname{Comp}_{1,\Psi}\phi)(v)\bigr)
 \bigr\|^2_{\mathrm{Weyl}}
 \le \exp\bigl(u+H-mS+2G_0+B\bigr),
 \qquad
 G_0=\sum_{p\in Q}\gamma_p\log p .
\]
Here $Q$ is the set of primes less than $x_1$ that are nonsplit in $\mathbb Q(\sqrt{-d})$.
\end{lemma}
\begin{proof}
By \statementref{Lemma}{lem:a-weak-chain-replacement} the composite lift can be written as a restricted lift for some \emph{weak chain datum} $\Psi'$ satisfying $\sum_w|\sigma(\Psi'(w))|\le\exp(x^{60})$. Let $S_j$ be the sum of the $s_\infty$ in the first $n_j$ rows at the infinite places for the upper-chain Young diagram at step $j$. The single-step Weyl norm bound for the kernel is $\exp(x^{50}+S_j/2)$. Applying the operator norm bounds successively, at any tuple of finite adelic points and translation parameter the product term obtained by acting on the source value with the $m$ kernels has norm bounded by
\[
 R=\exp(u/2)\prod_{j=0}^{m-1}\exp(x^{50}+S_j/2).
\]
The value before restriction is the sum of these product terms. At each step the class number is at most $\exp(x^{30})$, so the number of tuples of classes is at most $\exp(mx^{30})$. Apply the triangle inequality over the classes and translation parameters, then the bound on the coefficient sum. The norm before restriction is at most $R\exp(mx^{30}+x^{60})$.\leanfn{See \href{https://github.com/lsernau/grh-lean/blob/72fead102f663d5c7a407964a13dc0cf9d977207/grh/src/New/UpperBound/Size/ArchimedeanSize.lean\#L359}{Product of kernels and class-sum estimates}.}

Write $h_*=r_0^8$, half the difference of the weights of consecutive lifts. The squared norm estimate for restriction to $Y_m$ multiplies the bound by
\[
 \exp\bigl(x^{70}-(|\Phi_a|a-fmh_*)\log N\bigr).
\]
Here the branching exponents sum to $|\Phi_a|a-fmh_*$. Combining the restriction estimate with the square of the previous bound gives the required exponent for the squared norm.\leanfn{See \href{https://github.com/lsernau/grh-lean/blob/72fead102f663d5c7a407964a13dc0cf9d977207/grh/src/New/UpperBound/Size/ArchimedeanSize.lean\#L402}{Restriction and the combined error terms}.}
\[
 u+\sum_{j=0}^{m-1} S_j- |\Phi_a|a\log N+E_0,
 \qquad
 E_0=x^{70}+fmh_*\log N+2(mx^{30}+x^{60}+mx^{50}).
\]

Let $A_0=m\bigl(|\Phi_a|s_\infty(a)+k\sum_{\tau\in\Phi}s_\infty(\kappa_\tau)\bigr)$ denote the main term in the size at step $m$. The total row size and the compensation from level primes satisfy
\[
 \begin{aligned}
 \sum_{j=0}^{m-1} S_j&\le A_0+B/2,\\
 A_0-|\Phi_a|a\log N&\le H-mS+2G_0+32fa\ell .
 \end{aligned}
\]
For the first bound, consider a row of main length $r$ with at most $h$ extra boxes. By $r+h\le2a$, the difference from the derivative size at the main length is at most
$h[\log(2a+1)+\log(M/\Omega)]\le25h\ell$. The main length contributes $|\Phi_a|s_\infty(a)+(k-1)\sum_\tau s_\infty(\kappa_\tau)$ at each step. Adding over the extra boxes gives total error at most
\[
 25mf\ell\bigl(b+a/x_1+n_mh_*^2\bigr)\le B/2 .
\]
By $s_\infty(\kappa_\tau)\ge0$, replacing $k-1$ by the larger $k$ bounds the sum of the main lengths by $A_0$. The second formula follows from the identity
\[
 H-mS+2G_0
 =m|\Phi_a|s_{\rm main}
   +mk\sum_{\tau}s_\infty(\kappa_\tau),
\]
with $A_0=m|\Phi_a|s_\infty(a)+mk\sum_\tau s_\infty(\kappa_\tau)$, which makes the second term on the right cancel in the difference. The size estimates for the level primes give
\[
 |\Phi_a|\bigl(m(s_\infty(a)-s_{\rm main})-a\log N\bigr)\le32fa\ell .
\]
The admissible scales also give $E_0+32fa\ell\le B/2$. Substituting the three estimates into the bound for the exponent gives $u+H-mS+2G_0+B$.\leanfn{See \href{https://github.com/lsernau/grh-lean/blob/72fead102f663d5c7a407964a13dc0cf9d977207/grh/src/New/UpperBound/Size/ArchimedeanSize.lean\#L126}{Level compensation and the gain identity}, \href{https://github.com/lsernau/grh-lean/blob/72fead102f663d5c7a407964a13dc0cf9d977207/grh/src/New/UpperBound/Size/ArchimedeanSize.lean\#L284}{Row sizes and the bound for the full chain}.}
\end{proof}

\section{Exceptional primes and elementary divisors of edge Hecke operators}
\label{sec:a-hecke}

We use the cyclotomic orders at exceptional primes to construct full markings, then bound the colength of the edge restriction pairing. An approximate inverse reduces the problem to bounding the number of extra generators, and the filtration by deep forms gives the needed dimension bound.

\thmsource{Lean source: \leansourcelink{src/New/UpperBound/Divisibility/MarkedPrimes.lean}{146}.}
\begin{lemma}[The cyclotomic order at an exceptional prime]
\label{lem:a-exceptional-order}
\emitthmsource
Let $p$ be an \emph{exceptional prime}, so $\chi_K(p)=-1$ and $p\not\equiv\pm1\pmod m$ and the order $o_0$ of $-p$ modulo $q_0$ is odd and at most $r_p=\lceil\log_p m\rceil$. Let $\rho_p$ be the residue class of $-p$ in $(\mathbb Z/q_1)^\times$, and let $o=\operatorname{ord}(\rho_p)$. Then $m\mid o$.
\end{lemma}
\begin{proof}
Every prime factor of $q_1$ splits in $K$, whereas $\chi_K(d)=0$. By $\chi_K(p)=-1$ we have $p\nmid q_1d$. In particular, $2\mid q_1$ makes $p$ odd. Write $m=2^e$. The admissible data give the tower exponent $s=22e$ and $2^{s+1}\mid q_1$, so $u=-p$ satisfies $u^o\equiv1\pmod{2^{s+1}}$. Write $o$ as $o=2^\nu v$, with $v$ odd. Factoring $u^o-1$ shows that the odd power $v$ does not change the $2$-adic valuation of $u^{2^\nu}-1$.

When $\nu=0$, the valuation is $v_2(p+1)\le e-1$ by $p\not\equiv-1\pmod m$, which cannot give the required congruence. Hence $\nu\ge1$, and further factorization gives
\[
 v_2(u^o-1)=v_2(u-1)+v_2(u+1)+(\nu-1).
\]
For each $j\ge1$, the factor $u^{2^j}+1$ contains exactly one factor of $2$. One of $p-1$ and $p+1$ has $2$-adic valuation $1$, and by $p\not\equiv\pm1\pmod m$ the other has valuation at most $e-1$. Therefore
\[
 s+1\le v_2(u^o-1)\le e+\nu-1.
\]
Combining this with $s=22e$ gives $\nu\ge e$, that is, $m\mid o$.\leanfn{See \href{https://github.com/lsernau/grh-lean/blob/72fead102f663d5c7a407964a13dc0cf9d977207/grh/src/New/UpperBound/Divisibility/MarkedPrimes.lean\#L136}{The order and $2$-adic valuation at an exceptional prime}, \href{https://github.com/lsernau/grh-lean/blob/72fead102f663d5c7a407964a13dc0cf9d977207/grh/src/New/Derivatives/Field/PrimeDecomposition.lean\#L723}{Inert primes and the conductor}.}
\end{proof}

Fix an exceptional prime $p<x_1$ and a $p$-adic embedding $\iota$. For each $\tau\in\Phi$ take the representative complex embedding agreeing with the fixed embedding on $K$, and write $w_\tau$ for the place induced by $\iota\circ\tau$ on $E$. Call a set $\Lambda$ of places a marking at $p$ if its places are all induced by $\Phi_a$ and $\Lambda\cap\bar\Lambda=\varnothing$. It is a \emph{full mark} if $w_\tau\in\Lambda$ or $\bar w_\tau\in\Lambda$ for every $\tau\in\Phi_a$.

\thmsource{Lean source: \leansourcelink{src/New/UpperBound/Divisibility/MarkedPrimes.lean}{225}.}
\begin{lemma}[Full markings at exceptional primes]
\label{lem:a-full-mark}
\emitthmsource
For every exceptional prime $p<x_1$ and $p$-adic embedding $\iota\colon\mathbb C\hookrightarrow\mathbb C_p$ there is a full marking.
\end{lemma}
\begin{proof}
Retain the notation of \statementref{Lemma}{lem:a-exceptional-order} and let $c$ denote complex conjugation on $E$. Since $p$ is inert in $K$, the Frobenius for $E/\mathbb Q$ acts by complex conjugation on $K$ and by the $p$-th power on the cyclotomic part. The element $\rho_p$ of $\operatorname{Gal}(E/K)$ acts by the $-p$-th power on that part, so the Frobenius is $c\rho_p$. For any place $w$ above $p$, Frobenius fixes $w$, hence $\rho_p$ sends $w$ to $\bar w$. The same relation holds for the pullback action on places.

By $m\mid o$ and the evenness of $m$, the integer $o$ is even. If $w=\bar w$, then $c$ belongs to the decomposition group. By $p\nmid q_1d$ this group is generated by Frobenius, so the element can be written as $c=(c\rho_p)^j$. Restriction to $K$ shows that $j$ is odd, and then $\rho_p^j=1$, contrary to evenness of $o$. Thus the places above $p$ form pairs of distinct complex conjugates. Since $o$ is even, $c\rho_p$ also has order $o$, and each residue field has $p^o$ elements.\leanfn{See \href{https://github.com/lsernau/grh-lean/blob/72fead102f663d5c7a407964a13dc0cf9d977207/grh/src/New/Derivatives/Field/PrimeDecomposition.lean\#L733}{The Frobenius action above an inert prime}, \href{https://github.com/lsernau/grh-lean/blob/72fead102f663d5c7a407964a13dc0cf9d977207/grh/src/New/UpperBound/Divisibility/MarkedPrimes.lean\#L171}{Residue fields and conjugate pairs}.}

If $\tau\in\Phi_a$, the residue class of $\tau$ modulo $q_0$ is $1$. As $\rho_p^{o_0}$ is also $1$ modulo $q_0$, the place $\rho_p^{o_0}\tau$ still belongs to $\Phi_a$. By the pullback definition of the induced place and the oddness of $o_0$,
\[
 w_{\rho_p^{o_0}\tau}=\overline{w_\tau}.
\]
Thus the finite set $\{w_\tau:\tau\in\Phi_a\}$ is stable under conjugation. Choose one place from each conjugate pair $\{w,\bar w\}$ to obtain a set $\Lambda$ meeting all three requirements. Also $\chi_K(p)=-1$ and $p<x_1$, so $p$ is a small nonsplit prime meeting the size conditions in the definition of a marking.\leanfn{See \href{https://github.com/lsernau/grh-lean/blob/72fead102f663d5c7a407964a13dc0cf9d977207/grh/src/New/UpperBound/Divisibility/MarkedPrimes.lean\#L205}{Induced places and construction of a full marking}.}
\end{proof}

\thmsource{Lean source: \leansourcelink{src/New/UpperBound/Divisibility/MarkedPrimes.lean}{486}.}
\begin{lemma}[Full markings eliminate collisions]
\label{lem:a-full-mark-collision}
\emitthmsource
Let $\Lambda$ be a full marking at $p$, and suppose $S_\Lambda=\{\tau\in\Phi:w_\tau\in\Lambda\}$. Then
\[
 C_p(S_\Lambda)=0.
\]
\end{lemma}
\begin{proof}
Each pair $(\tau,j)$ in the collision sum must have $j$ odd and $\tau,\rho_p^j\tau$ both in $\Phi_a\setminus S_\Lambda$. Fullness makes $w_\tau\notin\Lambda$ imply $\bar w_\tau\in\Lambda$. Since $j$ is odd, however, $w_{\rho_p^j\tau}=\bar w_\tau$ and hence $\rho_p^j\tau\in S_\Lambda$, a contradiction. Therefore all terms of the sum vanish.\leanfn{See \href{https://github.com/lsernau/grh-lean/blob/72fead102f663d5c7a407964a13dc0cf9d977207/grh/src/New/UpperBound/Divisibility/MarkedPrimes.lean\#L481}{Full markings and the collision sum}.}
\end{proof}

\thmsource{Lean source: \leansourcelink{src/New/UpperBound/ElementaryDivisors/ElementaryDivisorTheorem.lean}{156}.}
\begin{lemma}[An approximate inverse and a rank bound for additional generators]
\label{lem:a-near-inverse-few-generators}
\emitthmsource
Let $O$ be a valuation ring with maximal ideal $\mathfrak m$, and let $\mathbf k=O/\mathfrak m$. Take a finite index set $J,F,I$ and linear maps
\[
 Q,B:O^J\longrightarrow O^{F\times I},
 \qquad R,A:O^{F\times I}\longrightarrow O^J .
\]
Let $RQ=1$, let $C\subseteq O^I$ be a submodule and $S\subseteq\mathbf k^I$ a subspace. Suppose the images of $Q$ and $B$ have each $F$-component in $C$. For a scalar $a,z,w\in O$ assume $w\in\mathfrak m$ and $aw\mid z$ and suppose
\[
 AB\equiv a\operatorname{Id}\pmod z,\qquad
 zh\in C\ \Longrightarrow\ \overline h\in S .
\]
Then there is a finite set $\mathcal S\subseteq O^J$ such that
\[
 |J|+|\mathcal S|\le |F|\dim_\mathbf kS,\qquad
 aO^J\subseteq \operatorname{im}(AQ)+\langle\mathcal S\rangle_O .
\]
\end{lemma}
\begin{proof}
By $RQ=1$ the operator $1-QR$ projects onto $\ker R$. Take the Smith decomposition of the square matrix $((1-QR)B)R$. Since $R$ is surjective, its image equals that of $(1-QR)B$. Thus there exist an invertible matrix $U$ and diagonal entries $t_i$ such that
\[
 \operatorname{im}Q+\operatorname{im}B
   =\operatorname{im}Q+\big\langle t_iU_i:i\in F\times I\big\rangle_O ,
\]
Here $U_i$ is column $i$ of $U$. By $t_iU_i\in\ker R$ and integrality in the ring, $t_i\ne0$ implies $RU_i=0$.\leanfn{See \href{https://github.com/lsernau/grh-lean/blob/72fead102f663d5c7a407964a13dc0cf9d977207/grh/src/New/UpperBound/ElementaryDivisors/ElementaryDivisorTheorem.lean\#L64}{Smith decomposition on the complementary summand}.}

Take $J_0=\{i:t_i\ne0,\ t_i\mid z\}$. For each $i\in J_0$, membership of $t_iU_i$ in the left-hand side above implies that $zU_i$ lies componentwise in $C$. Similarly, multiplying each column of $Q$ by $z$ gives componentwise membership in $C$. Therefore the columns of $Q$ and the columns $U_i$ of $J_0$ reduce into $S^F$. They are linearly independent over the residue field $\mathbf k$: left multiplication of a relation by $\overline R$ forces the coefficients of the columns of $\overline Q$ to vanish. Invertibility of $\overline U$ then forces all remaining coefficients to vanish. Hence $|J|+|J_0|\le |F|\dim_\mathbf kS$.\leanfn{See \href{https://github.com/lsernau/grh-lean/blob/72fead102f663d5c7a407964a13dc0cf9d977207/grh/src/New/UpperBound/ElementaryDivisors/ElementaryDivisorTheorem.lean\#L117}{Independence of the two sets of columns over the residue field}.}

Put $\mathcal S=\{AU_i:i\in J_0\}$ and let $G=\operatorname{im}(AQ)+\langle\mathcal S\rangle_O$. Divisibility in the valuation ring is a total preorder, so $i\notin J_0$ implies $z\mid t_i$. Thus after applying $A$ to $\operatorname{im}Q+\operatorname{im}B$ the images of all $t_iU_i$ lie in $G$ or $zO^J$. Together with $AB\equiv a\operatorname{Id}\pmod z$ this gives
\[
 aO^J\subseteq G+zO^J\subseteq G+\mathfrak m\,aO^J .
\]
Apply Nakayama's lemma to the finitely generated module $(G+aO^J)/G$ to get $aO^J\subseteq G$. By $|\mathcal S|\le|J_0|$ the number of generators satisfies the required bound as well.\leanfn{See \href{https://github.com/lsernau/grh-lean/blob/72fead102f663d5c7a407964a13dc0cf9d977207/grh/src/New/UpperBound/ElementaryDivisors/ElementaryDivisorTheorem.lean\#L144}{Removing the precision error by Nakayama's lemma}, \href{https://github.com/lsernau/grh-lean/blob/72fead102f663d5c7a407964a13dc0cf9d977207/grh/src/New/UpperBound/ElementaryDivisors/ElementaryDivisorTheorem.lean\#L162}{A generator bound from the approximate inverse}.}
\end{proof}

\thmsource{Lean source: \leansourcelink{src/New/UpperBound/ElementaryDivisors/Depth.lean}{752}.}
\begin{lemma}[Total dimension of the deep-form filtration]
\label{lem:a-edge-deep-form-count}
\emitthmsource
Let $T$ be a \emph{Frobenius cycle} at the prime $p$, with coefficient field $\mathbf k$ of characteristic $p$. Write $\nu$ for the cycle length and $P$ for the Frobenius power. Take $0<2r\le k$ and degree $n$, and let $\mathcal B$ index monomials of degree $n$ in each of $\nu$ copies of $k$ variables. The row index $b$ of variable $(b,i)$ starts at zero, and the parity of the copy index $i$ alternates along the cycle. A variable is deep if $i$ is even and $b\ge r$, or if $i$ is odd and $b<r$. The depth of a monomial is its degree in the deep variables. Let $\mathcal D_\delta$ be the space of coefficient vectors whose every $g\in\mathrm{GL}_k(W)$ by $\rho(g)$ contains only monomials of depth at least $\delta$. Here $\rho(g)$ is the tensor product of the $n$-th symmetric powers, using $g$ under the embedding of every even copy and $(g^{-1})^\top$ for the odd copies. Then for every integer $x\ge0$,
\[
 \sum_{\delta=1}^x\dim_{\mathbf k}\mathcal D_\delta
 \le c_{\rm div}|\mathcal B|,
 \qquad
 c_{\rm div}=\nu n\left(\frac1{P^{2r}}+\frac{r^2}{kP}\right).
\]
\end{lemma}
\begin{proof}
Give each variable $(b,i)$ a real weight: $0$ for a variable that is not deep, $P^{-1}$ for a deep variable when $b<2r$, and $P^{-2r}$ when $b\ge2r$. Write $\operatorname{wt}(\beta)$ for the resulting weighted sum of the exponents of a monomial $\beta$. The weighted-degree lemma for deep forms shows that every nonzero $F\in\mathcal D_\delta$ has a nonzero coefficient at some monomial $\beta$ with $\operatorname{wt}(\beta)\ge\delta$. Projecting $F$ to the coefficients of these monomials is thus injective, giving
\[
 \dim\mathcal D_\delta
 \le \#\{\beta\in\mathcal B:\operatorname{wt}(\beta)\ge\delta\}.
\]
Sum over $\delta=1,\ldots,x$. Each fixed $\beta$ is counted at most $\lfloor\operatorname{wt}(\beta)\rfloor$ times, so the left-hand side is at most $\sum_\beta\operatorname{wt}(\beta)$.\leanfn{See \href{https://github.com/lsernau/grh-lean/blob/72fead102f663d5c7a407964a13dc0cf9d977207/grh/src/New/UpperBound/ElementaryDivisors/Depth.lean\#L612}{Weighted degree of deep forms}, \href{https://github.com/lsernau/grh-lean/blob/72fead102f663d5c7a407964a13dc0cf9d977207/grh/src/New/UpperBound/ElementaryDivisors/Depth.lean\#L755}{Summing dimensions}.}

In a single copy, permuting any two rows preserves $\mathcal B$, and each monomial has degree $n$ in that copy. The mean exponent of each variable is therefore $n/k$. Since the parity alternates, each row is deep in exactly half the copies, so
\[
 \begin{aligned}
 \frac1{|\mathcal B|}\sum_\beta\operatorname{wt}(\beta)
   &=\frac nk\cdot\frac{\nu}{2}
       \left(\frac{2r}{P}+\frac{k-2r}{P^{2r}}\right)\\
   &\le \nu n\left(\frac1{P^{2r}}+\frac{r^2}{kP}\right),
 \end{aligned}
\]
Here $r\ge1$ was used to deduce $r\le r^2$.\leanfn{See \href{https://github.com/lsernau/grh-lean/blob/72fead102f663d5c7a407964a13dc0cf9d977207/grh/src/New/UpperBound/ElementaryDivisors/Depth.lean\#L62}{The number of deep copies for each row}, \href{https://github.com/lsernau/grh-lean/blob/72fead102f663d5c7a407964a13dc0cf9d977207/grh/src/New/UpperBound/ElementaryDivisors/Depth.lean\#L690}{The mean exponent of one variable}.}
\end{proof}

\thmsource{Lean source: \leansourcelink{src/New/UpperBound/ElementaryDivisors/ElementaryDivisorTheorem.lean}{876}.}
\begin{lemma}[A basis of starting functions and the endpoint colength]
\label{lem:a-edge-origin-basis}
\emitthmsource
Keep the notation for the Frobenius cycle, with coefficient ring $O$ now a valuation ring of characteristic zero and $p$ a nonunit of $O$. Let $w\in O$ be the element furnished by the cycle with $w^2=p$. Assume $0<2r\le k$, $0<\theta\le1$, $n^4\le q=P^\nu$, and $(\nu k\log q/\theta)^5P^{3r}\le n$. Take a set $V$ of starting points and a set $E_0$ of edges. Write $a(e)$ for the start of edge $e$ and $g(e)\in\mathrm{GL}_k(W)$ for its basis change. Suppose $e\mapsto(a(e),[g(e)^{-1}])$ surjects onto $V\times\mathcal P$, where $\mathcal P$ is the coset set of the quasiparabolic with block sizes $r,k-r$. Put $\Delta=\operatorname{diag}(w^{\operatorname{dep}(\beta)})$ and define edge restriction by
\[
 JF(e)=\Delta\rho(g(e))F(a(e)).
\]
The starting-function module is the saturation of the image of $J$ with respect to $w$:
\[
 \mathcal F_{\mathrm{orig}}
   =\{\psi\in O^{E_0\times\mathcal B}:
                          w^a\psi\in\operatorname{im}J
                          \text{ for some }a\ge0\}.
\]
Use divisibility of the determinant to measure \emph{colength}. Then there is a matrix $Q,R,d$ satisfying $\operatorname{im}Q=\mathcal F_{\mathrm{orig}}$, $RQ=1$, and $J=Qd$, and
\[
 \det d\mid w^{|V|\lfloor c_{\rm div}|\mathcal B|\rfloor}.
\]
\end{lemma}
\begin{proof}
Surjectivity from edges to starting points shows that $J$ has at least as many rows as columns. Over the valuation ring, take the Smith decomposition of $J$ and write
\[
 JU^{-1}=Q\operatorname{diag}(t_i),\qquad RQ=1,
 \qquad d=\operatorname{diag}(t_i)U ,
\]
Here $U$ is invertible and $i$ ranges over $V\times\mathcal B$. Each $t_i$ therefore divides every coordinate in column $i$ of $JU^{-1}$. For a start $v$ and any $\gamma\in\mathrm{GL}_k(W)$ there is an edge $e$ with $a(e)=v$ and $\gamma g(e)^{-1}$ in the specified quasiparabolic. For every element $u$ of this subgroup, depth scaling gives $\Delta\rho(u)=A_u\Delta$ with $A_u$ integral. Thus each coefficient of the component $F_i(v)$ of this column at $v$ satisfies
\[
 t_i\mid w^{\operatorname{dep}(\beta)}
               \bigl(\rho(\gamma)F_i(v)\bigr)_\beta .
\]
Let $x$ bound all depths. Taking $\gamma=1$ shows that $t_i$ divides $w^x$ times every coordinate of column $i$ of $U^{-1}$. Pairing with row $i$ of $U$ gives $t_i\mid w^x$.\leanfn{See \href{https://github.com/lsernau/grh-lean/blob/72fead102f663d5c7a407964a13dc0cf9d977207/grh/src/New/UpperBound/ElementaryDivisors/Edges.lean\#L323}{The quasiparabolic preserves the halfway depth lattice}, \href{https://github.com/lsernau/grh-lean/blob/72fead102f663d5c7a407964a13dc0cf9d977207/grh/src/New/UpperBound/ElementaryDivisors/Edges.lean\#L612}{Smith decomposition of edge restriction}.}

By the $t_i\mid w^x$, multiplying any $QF$ by $w^x$ puts it in $\operatorname{im}J$, so $\operatorname{im}Q\subseteq\mathcal F_{\mathrm{orig}}$. Conversely, if $w^a\psi\in\operatorname{im}J$ then $QR$ fixes $w^a\psi$. Cancel the nonzero scalar $w^a$ to obtain $QR\psi=\psi$.\leanfn{See \href{https://github.com/lsernau/grh-lean/blob/72fead102f663d5c7a407964a13dc0cf9d977207/grh/src/New/UpperBound/ElementaryDivisors/Edges.lean\#L676}{Projection onto the starting-function module}.}

If $t_i\nmid w^\delta$, those divisibilities imply that for each starting point $v$ and each $\gamma$, the coefficients of monomials of depth at most $\delta$ in the $\gamma$-translate of $F_i(v)$ are all nonunits. A unit would cancel and give $t_i\mid w^\delta$. Thus the reduction of column $i$ belongs pointwise to the deep-form space $\mathcal D_{\delta+1}$ over the residue field. Since $U^{-1}$ is invertible, the columns remain independent there, so
\[
 \#\{i:t_i\nmid w^\delta\}\le |V|\dim\mathcal D_{\delta+1}.
\]
Put $s_i=\min\{s\ge0:t_i\mid w^s\}\le x$. Then $\sum_i s_i=\sum_{\delta=0}^{x-1}\#\{i:t_i\nmid w^\delta\}$, and $\det d$ agrees with $\prod_i t_i$ up to a unit. Hence
\[
 \det d\mid w^{\,|V|\sum_{\delta=1}^x\dim\mathcal D_\delta}.
\]
The residue field has characteristic $p$ and reduction preserves the Frobenius power $P$ of $T$. \statementref{Lemma}{lem:a-edge-deep-form-count} and rounding the sum of integer dimensions give the required exponent.\leanfn{See \href{https://github.com/lsernau/grh-lean/blob/72fead102f663d5c7a407964a13dc0cf9d977207/grh/src/New/UpperBound/ElementaryDivisors/Edges.lean\#L654}{Counting deep forms from the Smith diagonal}, \href{https://github.com/lsernau/grh-lean/blob/72fead102f663d5c7a407964a13dc0cf9d977207/grh/src/New/UpperBound/ElementaryDivisors/ElementaryDivisorTheorem.lean\#L886}{The determinant bound for the endpoint basis decomposition}.}
\end{proof}

In the Frobenius cycle, the coefficient representation $\rho(g)$ acts on even copies by the embedded $g$ and on odd copies by the embedded $(g^{-1})^{\mathsf T}$. For adjacent rows $r-1,r$, write $u_i=(r,i)$ and $\bar u_i=(r-1,i)$ and swap these two variables in the odd copies. If the cycle has length $\nu$ and $q=P^\nu$, then the \emph{defect} of an exponent $\beta$ along this root in copy $i$ is
\[
 \operatorname{def}_i(\beta)
   =\beta_{u_i}+\frac1{q-1}\sum_{j\in\mathbb Z/\nu}
        P^{[j-i]_\nu}\bigl(\beta_{u_j}-\beta_{\bar u_j}\bigr),
\]
Here $[j-i]_\nu$ is the representative in $0,\ldots,\nu-1$.\leanfn{See \href{https://github.com/lsernau/grh-lean/blob/72fead102f663d5c7a407964a13dc0cf9d977207/grh/src/New/UpperBound/ElementaryDivisors/TwistingData.lean\#L166}{Root defect weights and coefficient conditions}.}

\thmsource{Lean source: \leansourcelink{src/New/UpperBound/ElementaryDivisors/AveragingMap.lean}{566}.}
\begin{lemma}[Eliminating bad row weights with a Frobenius averaging matrix]
\label{lem:a-edge-average}
\emitthmsource
Let $W$ be a Frobenius cycle of $\nu=2h$ copies in the valuation ring $O$ of characteristic zero. Suppose $w^2=p$, with $p$ a nonunit in $O$, and $q=P^\nu$ the size of the residue field of $W$.
Let $\mathcal B$ index monomials of degree $n$ in each of $\nu$ copies of $k$ variables.
Use the block sizes $r,k-r$ with $1\le r<k$ and compute defects along the adjacent rows $r-1,r$.
Let $\nu k(h+2)n+n+k<P^h$. For each $c\ge0$ write $c_\ast=3kc$ and $c_0=k^2\log_p q$.
Then for the weighted pairing $\langle\ ,\ \rangle_b$ satisfying edge adjointness, there is an integral matrix $\mathrm{Av}$ such that
\[
 \begin{aligned}
 w^c&\mid (\rho(g)\mathrm{Av} F)_\beta
     &&\text{if $g$ is block diagonal and some $\operatorname{def}_i(\beta)>3c_\ast$},\\
 \langle\mathrm{Av} F,G\rangle_b
   &\equiv p^{\nu c_0}\langle F,G\rangle_b \pmod{w^c}
     &&\text{if $G$ is in the halfway lattice of every outgoing edge}.
 \end{aligned}
\]
\end{lemma}

\begin{proof}
Put $D=\nu nk(h+2)P^h$, so $D+k<P^{2h}=q$. Substitute generic matrices in the variables of each copy and apply the Frobenius substitutions to obtain polynomials in the matrix coefficients. Use finite-field averaging on generic matrices to choose a linear form $I_j$ for each $j$. It is invariant under block diagonal conjugation and kills monomials of row weight greater than $2c_\ast P^h$. On polynomials of degree at most $D$ whose bad row-weight coefficients vanish modulo $w^c$ for every conjugation, it agrees with $e=p^{c_0}\operatorname{ev}_{1}$ modulo $w^c$. The row weight is the degree in row $r$ minus that in row $r-1$.\leanfn{See \href{https://github.com/lsernau/grh-lean/blob/72fead102f663d5c7a407964a13dc0cf9d977207/grh/src/New/UpperBound/ElementaryDivisors/AveragingMap.lean\#L579}{Degree budget and selection of the averages}.}

For $J\subseteq\mathbb Z/\nu$ let $\mathrm{Av}_J$ be the coefficient operator obtained by applying $I_j$ to copy $j\in J$ and $e$ to the copies $j\notin J$, and take $\mathrm{Av}=\mathrm{Av}_{\mathbb Z/\nu}$. Consider a bad exponent $\beta$ of defect exceeding $3c_\ast$ in copy $i$. Keep copy $j=i-h$ and contract all others first. In the resulting generic matrix polynomial, the row weight divided by $P^h$ differs from $\operatorname{def}_i(\beta)$ by at most $c_\ast$ for any coefficient not divisible by $w^c$. These row weights are among those killed by $I_j$. Thus applying $I_j$ gives a value divisible by $w^c$. Invariance under block diagonal conjugation gives the calculation for every $\rho(g)\mathrm{Av}F$.\leanfn{See \href{https://github.com/lsernau/grh-lean/blob/72fead102f663d5c7a407964a13dc0cf9d977207/grh/src/New/UpperBound/ElementaryDivisors/AveragingMap.lean\#L426}{Comparing row weights and defects}, \href{https://github.com/lsernau/grh-lean/blob/72fead102f663d5c7a407964a13dc0cf9d977207/grh/src/New/UpperBound/ElementaryDivisors/AveragingMap.lean\#L477}{Eliminating coefficients of large defect}.}

If $j\notin J$, the difference of the two pairings is $(I_j-e)\Psi$, where $\Psi$ is a polynomial of degree at most $D$ obtained by keeping copy $j$ and pairing with $G$. In each block diagonal conjugate of $\Psi$, a coefficient of row weight exceeding $2c_\ast P^h$ is a sum whose terms either have contraction coefficient divisible by $w^c$ or correspond to exponent $\beta$ of defect exceeding $c_\ast$ in copy $j+h$. In the latter case the large-defect pairing bound for the halfway lattice makes its pairing with $G$ divisible by $w^c$. Thus $(I_j-e)\Psi\equiv0\pmod{w^c}$. Add the copies successively to $J$ and sum the congruences. Using $\mathrm{Av}_{\varnothing}=p^{\nu c_0}I$ gives the second formula.\leanfn{See \href{https://github.com/lsernau/grh-lean/blob/72fead102f663d5c7a407964a13dc0cf9d977207/grh/src/New/UpperBound/ElementaryDivisors/AveragingMap.lean\#L512}{Congruence of the pairing when adding one copy}.}
\end{proof}

\thmsource{Lean source: \leansourcelink{src/New/UpperBound/ElementaryDivisors/DefectCount.lean}{921}.}
\begin{lemma}[The quasiparabolic index and the dimension of bounded-defect forms]
\label{lem:a-edge-defect-count}
\emitthmsource
Over a coefficient field $\mathbf k$ of characteristic $p$ take Frobenius data with cycle length $\nu$ and residue field size $q=P^\nu$. Let $C_t$ be the space of degree $n$ in each copy with bounded defect: after any translate by the Levi group of block sizes $r,k-r$, the coefficients of defect exceeding $3t$ along any root $(r-1,r)$ in any copy vanish. Suppose $n>0$, $t\ge2(k+1)$, $\theta\le1$, and $8\nu kt P^{2r}\le\theta n$. If $\mathcal P$ is the corresponding quasiparabolic coset set, then
\[
 |\mathcal P|\dim_{\mathbf k} C_t
 \le (1+\theta)\binom{n+k-1}{k-1}^{\nu}.
\]
\end{lemma}

\begin{proof}
\phantomsection\label{par:a-stepanov-count}
\emph{Stepanov's count in the root lemma.} Let $W$ be a local ring with $q$ residue classes, let $\iota:W\to R$ be a homomorphism to a commutative ring, and let $a$ range over a set of representatives for the residue classes. If a nonzero polynomial $\Phi\in R[T]$ is divisible by $(T-\iota(a))^e$ for every $a$, where $e$ is a nonnegative integer, then $qe\le\deg\Phi$. Indeed, distinct representatives differ by units in $W$, so the monic factors are pairwise coprime. Their product divides $\Phi$ and has degree $qe$. This is the local form of the \hyperref[par:stepanov-method]{zero count in the preface}.\leanfn{See \href{https://github.com/lsernau/grh-lean/blob/72fead102f663d5c7a407964a13dc0cf9d977207/grh/src/New/UpperBound/ElementaryDivisors/RootLemma.lean\#L58}{Counting zeros with multiplicity at residue classes}.}

Among the translates under $\mathrm{SL}_2(W)$ acting on the two adjacent rows, the root lemma picks a term $X^\beta$ minimizing the sum of the coefficient's $w$-adic order and the monomial weight. Write $\ell$ for the $w$-adic order of its coefficient. The proof makes a Frobenius substitution in the original polynomial and extracts the coefficient associated with $\beta$ to obtain a univariate auxiliary polynomial $\Phi$ modulo $w^{\ell+1}$. Let $d$ and $e$ be the pole and zero weights of $\beta$, respectively. Expansion at infinity gives $\deg\Phi\le d$ and identifies the coefficient of $T^d$ with the reduction of the selected coefficient, so $\Phi\ne0$. The Taylor expansions at the residue representatives give vanishing to order at least $e$. The two weights satisfy
\[
 \operatorname{def}_i(\beta)=\frac{e-d/q}{q-1}.
\]
Thus $qe\le\deg\Phi\le d$ implies $\operatorname{def}_i(\beta)\le0$. In the characteristic-$p$ case used here, applying this to the simple roots within the blocks yields a nonzero term of nonpositive defect along all these roots in a translate by a representative of a Levi flag. The lattice-point count below uses this fact.\leanfn{See \href{https://github.com/lsernau/grh-lean/blob/72fead102f663d5c7a407964a13dc0cf9d977207/grh/src/New/UpperBound/ElementaryDivisors/RootLemma.lean\#L873}{The auxiliary polynomial and degree comparison}, \href{https://github.com/lsernau/grh-lean/blob/72fead102f663d5c7a407964a13dc0cf9d977207/grh/src/New/UpperBound/ElementaryDivisors/RootLemma.lean\#L942}{The root lemma}, and \href{https://github.com/lsernau/grh-lean/blob/72fead102f663d5c7a407964a13dc0cf9d977207/grh/src/New/UpperBound/ElementaryDivisors/DefectCount.lean\#L450}{The choice of a Levi flag representative}.}

Write $E=4tP^{2r}$, let $\mathcal L$ be the Levi group, and let $\mathcal I$ be the lower Iwahori, consisting of matrices reducing to lower triangular ones. Counting lattice points in the defect cone gives the upper bounds $D_i\in[1,n+E]$ for the copies, such that
\[
 \dim_{\mathbf k} C_t
 \le [\mathcal L:\mathcal L\cap\mathcal I]\,
       \operatorname{vol}(\mathcal D_0\cap S_D).
\]
Here $S_D$ is the region of nonnegative exponent vectors whose degrees in each copy lie in $[D_i-1,D_i]$, and $\mathcal D_0$ requires nonpositive defect at every adjacent root. The estimate combines degree-$n$ exponents and grid indices of denominator $q^j$ into translates and adds a fixed shift taking bounded defect into the zero-defect cone. Record the coefficients for each Levi flag and grid index, divide the lattice-point dimension bound by $q^{j k\nu}$, and pass to the limit. The intersection of the cone with the degree region is bounded with boundary of measure zero, so lattice-point density converges to volume.\leanfn{See \href{https://github.com/lsernau/grh-lean/blob/72fead102f663d5c7a407964a13dc0cf9d977207/grh/src/New/UpperBound/ElementaryDivisors/DefectCount.lean\#L752}{Flags, grid indices, and the dimension bound}.}

Cyclic differences turn the zero-defect conditions into Pieri inequalities. The Pieri cone occupies fraction $\prod_{i=1}^k(q-1)/(q^i-1)$ of the volume of $S_D$. The quasiparabolic level index times the number of Levi flags is at most $\prod_{i=1}^k(q^i-1)/(q-1)$, since the block flag and parabolic coset determine the full flag. Thus multiplication by $|\mathcal P|$ leaves just $\operatorname{vol}(S_D)$.\leanfn{See \href{https://github.com/lsernau/grh-lean/blob/72fead102f663d5c7a407964a13dc0cf9d977207/grh/src/New/UpperBound/ElementaryDivisors/DefectCount.lean\#L872}{The flag index}, \href{https://github.com/lsernau/grh-lean/blob/72fead102f663d5c7a407964a13dc0cf9d977207/grh/src/New/UpperBound/ElementaryDivisors/DefectCount.lean\#L909}{The volume fraction of the cone in the degree region}.}

The interval width in each copy is $1$. Integrate over the last coordinate in the nonnegative orthant to bound the volume of each copy by $D_i^{k-1}/(k-1)!$. Multiply and use $D_i\le ne^{E/n}$ to get
\[
 \operatorname{vol}(S_D)\le e^{\nu(k-1) E/n}
           \bigl(n^{k-1}/(k-1)!\bigr)^{\nu}.
\]
By $2\nu k E\le\theta n$ the exponent $z=\nu(k-1)E/n$ satisfies $0\le z\le\theta/2\le 1/2$, hence $e^z\le1+2z\le1+\theta$. The claim follows from $n^{k-1}/(k-1)!\le\binom{n+k-1}{k-1}$.\leanfn{See \href{https://github.com/lsernau/grh-lean/blob/72fead102f663d5c7a407964a13dc0cf9d977207/grh/src/New/UpperBound/ElementaryDivisors/DefectCount.lean\#L939}{The budget for increasing the degree}.}
\end{proof}

\thmsource{Lean source: \leansourcelink{src/New/UpperBound/ElementaryDivisors/ElementaryDivisorTheorem.lean}{672}.}
\begin{lemma}[Edge congruence conditions and reduction spaces]
\label{lem:a-edge-condition}
\emitthmsource
Let $T$ be a Frobenius cycle from the local ring $W$ to the valuation ring $O$ of characteristic zero, of cycle length $\nu$, with $w^2=p$ and prime $p$ a nonunit in $O$. The cycle gives $\nu$ ring maps $\iota_i:W\to O$ satisfying $\iota_{i+1}(x)\equiv\iota_i(x)^P\pmod p$, where $P$ is a power of $p$ with $P\ge5$, and the residue field of $W$ has size $q=P^\nu$. Let $\mathcal B$ index monomials of degree $n$ in each of $\nu$ copies of $k$ variables. Assume
\[
 0<2r\le k,\quad 0<\theta\le1,\quad n^4\le q,\quad
 (\nu k\log q/\theta)^5 P^{3r}\le n .
\]
Put $c_0=k^2\log_p q$ and $c=5\nu c_0$, and let $\mathcal P$ denote the coset set of the quasiparabolic with block sizes $r,k-r$. Suppose weight $b,b'$ gives the adjoint pairing
\(\langle F,G\rangle_b=G^\top\operatorname{diag}(b)F\), satisfying $b_\beta b'_\beta=p^{\nu c_0}$. In the standard halfway lattice coordinates, put
\[
 \mathcal D=\{F\in O^{\mathcal B}:
       \Delta\rho(g)\Delta^{-1}F\in O^{\mathcal B}
            \text{ for every }g\in\mathrm{GL}_k(W)\},
 \qquad
 \Delta=\operatorname{diag}(w^{\operatorname{dep}(\beta)}).
\]
where $\Delta^{-1}$ is defined over the fraction field. This says $F$ belongs to the halfway lattices of all outgoing edges. Then there is a submodule $C\subseteq O^{\mathcal B}$, an averaging matrix $\mathrm{Av}$, and a subspace $S\subseteq\mathbf k^{\mathcal B}$ over the residue field $\mathbf k$ satisfying
\[
 \begin{gathered}
 \mathcal D+w^cO^{\mathcal B}\subseteq C,\qquad
 \mathrm{Av}\bigl(O^{\mathcal B}\bigr)\subseteq C,\\
 \langle \mathrm{Av}F,G\rangle_b
     \equiv p^{\nu c_0}\langle F,G\rangle_b \pmod{w^c}
          \quad (G\in\mathcal D),\\
 w^{c-1}F\in C\ \Longrightarrow\ \overline F\in S,\qquad
 |\mathcal P|\dim_\mathbf k S\le (1+\theta)|\mathcal B|.
 \end{gathered}
\]
\end{lemma}
\begin{proof}
Take $c_\ast=3kc$. Define $C$ by requiring that after each translate by the block diagonal group $\mathrm{GL}_r(W)\times\mathrm{GL}_{k-r}(W)$, every coefficient whose defect in some copy exceeds $3c_\ast$ belongs to $(w^c)$. Impose the same conditions on the reduced cycle but require those coefficients to vanish. This gives $S$, where the reduced cycle is obtained by composing each embedding of $T$ with $O\to\mathbf k$, and $\mathbf k$ has characteristic $p$.\leanfn{See \href{https://github.com/lsernau/grh-lean/blob/72fead102f663d5c7a407964a13dc0cf9d977207/grh/src/New/UpperBound/ElementaryDivisors/ElementaryDivisorTheorem.lean\#L667}{Characteristic of the residue field}.}

The large-defect pairing estimate gives $\mathcal D\subseteq C$. Since $\rho(g)$ is integral, we also have $w^cO^{\mathcal B}\subseteq C$. If $w^{c-1}F\in C$, for each relevant translate and monomial we obtain $w^c\mid w^{c-1}(\rho(g)F)_\beta$. Canceling $w^{c-1}$ gives $w\mid(\rho(g)F)_\beta$. Because $p$ is a nonunit, $w$ lies in the maximal ideal. Thus all these coefficients reduce to zero, which gives $\overline F\in S$.\leanfn{See \href{https://github.com/lsernau/grh-lean/blob/72fead102f663d5c7a407964a13dc0cf9d977207/grh/src/New/UpperBound/ElementaryDivisors/Edges.lean\#L773}{Coefficient conditions for the halfway lattice}, \href{https://github.com/lsernau/grh-lean/blob/72fead102f663d5c7a407964a13dc0cf9d977207/grh/src/New/UpperBound/ElementaryDivisors/ElementaryDivisorTheorem.lean\#L647}{Reduction after removing one level of precision}.}

As the parity alternates along the cycle, put $h=\nu/2$, an integer. The size assumptions give
\[
 \nu k(h+2)n+n+k<P^h.
\]
Indeed $n^4\le P^{2h}$ gives $n^2\le P^h$. By $n\ge(\nu k)^5$ and $\nu,k\ge2$, we have $h+2\le\nu k$, $k<n$, and $(\nu k)^2+2\le n$, so the left side is strictly below $((\nu k)^2+2)n\le n^2$. This meets the hypotheses for the averaging matrix. By \statementref{Lemma}{lem:a-edge-average}, the image of $\mathrm{Av}$ lies in $C$ and satisfies the claimed congruence modulo $w^c$ when paired with $\mathcal D$.\leanfn{See \href{https://github.com/lsernau/grh-lean/blob/72fead102f663d5c7a407964a13dc0cf9d977207/grh/src/New/UpperBound/ElementaryDivisors/ElementaryDivisorTheorem.lean\#L495}{The degree inequality for the averaging matrix}, \href{https://github.com/lsernau/grh-lean/blob/72fead102f663d5c7a407964a13dc0cf9d977207/grh/src/New/UpperBound/ElementaryDivisors/AveragingMap.lean\#L579}{The averaging matrix and its pairing congruence}.}

Finally the degree condition gives $8\nu k c_\ast P^{2r}\le\theta n$. We also have $c>0$ and $k\ge2$, so $c_\ast\ge 2(k+1)$. \statementref{Lemma}{lem:a-edge-defect-count} gives $|\mathcal P|\dim_\mathbf kS\le(1+\theta)|\mathcal B|$, and the preceding steps furnish the required data.\leanfn{See \href{https://github.com/lsernau/grh-lean/blob/72fead102f663d5c7a407964a13dc0cf9d977207/grh/src/New/UpperBound/ElementaryDivisors/ElementaryDivisorTheorem.lean\#L577}{Degree bound for the defect threshold}, \href{https://github.com/lsernau/grh-lean/blob/72fead102f663d5c7a407964a13dc0cf9d977207/grh/src/New/UpperBound/ElementaryDivisors/DefectCount.lean\#L939}{Dimension of the reduced conditions and the quasiparabolic index}, \href{https://github.com/lsernau/grh-lean/blob/72fead102f663d5c7a407964a13dc0cf9d977207/grh/src/New/UpperBound/ElementaryDivisors/ElementaryDivisorTheorem.lean\#L688}{Constructing the edge congruence conditions}.}
\end{proof}

\thmsource{Lean source: \leansourcelink{src/New/UpperBound/ElementaryDivisors/ElementaryDivisorTheorem.lean}{944}.}
\begin{lemma}[An approximate inverse of the edge pairing]
\label{lem:a-edge-pairing-near-inverse}
\emitthmsource
Let $C,S,\mathrm{Av}$ satisfy \statementref{Lemma}{lem:a-edge-condition} and let $J$ be a finite coordinate set. Suppose $Q':O^J\to O^{(H\times\mathcal E)\times\mathcal B}$ has left inverse $R'$ and each local component of $Q'$ lies in $\mathcal D$. For each edge $e\in\mathcal E$ take an integral invertible auxiliary transformation $u_e$. Let $P_b$ apply $u_e\otimes\operatorname{diag}(b)$ edge by edge, and define
\[
 \mathrm{Add}=Q'^\top P_b,\qquad
 \mathrm{Inv}=\Bigl(\bigoplus_e
          u_e^{-1}\otimes\bigl(\mathrm{Av}\operatorname{diag}(b')\bigr)\Bigr)R'^\top ,
\]
where the local tensor coordinates use the same ordering. The image of $\mathrm{Inv}$ lies componentwise in $C$, and
\[
 \mathrm{Add}\,\mathrm{Inv}\equiv p^{2\nu c_0}\operatorname{Id}
         \pmod{w^c}.
\]
\end{lemma}
\begin{proof}
The factors $u_e$ and $u_e^{-1}$ in the product cancel. Thus entry $(i,j)$ of $\mathrm{Add}\,\mathrm{Inv}$ is the sum over $x\in H\times\mathcal E$ of
\[
 \left\langle\mathrm{Av}\bigl(\operatorname{diag}(b')R'_j(x)\bigr),
                      Q'_i(x)\right\rangle_b ,
\]
Here $Q'_i(x)$ and $R'_j(x)$ are the coefficient vectors of column $i$ of $Q'$ and row $j$ of $R'$, respectively, at $x$. Since $Q'_i(x)\in\mathcal D$, modulo $w^c$ each $\mathrm{Av}$ may be replaced by $p^{\nu c_0}$ times the identity. By $b_\beta b'_\beta=p^{\nu c_0}$ the same scalar appears again. The sum is therefore congruent modulo $w^c$ to
\[
 p^{2\nu c_0}\sum_x Q'_i(x)^\top R'_j(x)
    =p^{2\nu c_0}(R'Q')_{ji}
    =p^{2\nu c_0}\delta_{ij}.
\]
The approximate inverse first takes integral linear combinations with $u_e^{-1}$ on auxiliary coordinates edge by edge, then applies the averaging matrix on the coefficient coordinates. The result lies componentwise in $C$.\leanfn{See \href{https://github.com/lsernau/grh-lean/blob/72fead102f663d5c7a407964a13dc0cf9d977207/grh/src/New/UpperBound/ElementaryDivisors/ElementaryDivisorTheorem.lean\#L312}{Entrywise product of the pairing and the approximate inverse}, \href{https://github.com/lsernau/grh-lean/blob/72fead102f663d5c7a407964a13dc0cf9d977207/grh/src/New/UpperBound/ElementaryDivisors/ElementaryDivisorTheorem.lean\#L958}{The approximate inverse congruence from the edge conditions}.}
\end{proof}

\thmsource{Lean source: \leansourcelink{src/New/UpperBound/ElementaryDivisors/ElementaryDivisorTheorem.lean}{990}.}
\begin{lemma}[Additional generators for the edge pairing]
\label{lem:a-edge-few-generators}
\emitthmsource
Retain the setting of \statementref{Lemma}{lem:a-edge-condition}. Let $V,H,\mathcal E$ be a finite set with $|H\times\mathcal E|=|V|\,|\mathcal P|$ and put $\mathcal L=O^{V\times\mathcal B}$. Let $Q,Q'$ be maps from $\mathcal L$ to $O^{(H\times\mathcal E)\times\mathcal B}$, each with a left inverse and with image lying locally in $\mathcal D$. Continue to let $P_b$ apply $u_e\otimes\operatorname{diag}(b)$ edge by edge and put $A_0=Q'^\top P_bQ$. There is a finite set $\mathcal S\subseteq\mathcal L$ such that
\[
 |\mathcal S|\le\theta\,|V|\,|\mathcal B|,\qquad
 p^{2\nu c_0}\mathcal L\ \subseteq\ A_0\mathcal L+\langle\mathcal S\rangle_O .
\]
\end{lemma}
\begin{proof}
Use the left inverse of $Q'$ to construct $\mathrm{Inv}$ in \statementref{Lemma}{lem:a-edge-pairing-near-inverse}, and put $a=p^{2\nu c_0}$ and $z=w^{c-1}$. From $\nu c_0\ge2$ and $c=5\nu c_0$ we get
\[
 aw=w^{4\nu c_0+1}\mid w^{5\nu c_0-1}=z .
\]
Here $w$ lies in the maximal ideal and $\mathrm{Add}\,\mathrm{Inv}\equiv a\operatorname{Id}\pmod{w^c}$. The congruence therefore also holds modulo $z$. The image of $Q$ lies componentwise in $\mathcal D\subseteq C$, and that of $\mathrm{Inv}$ lies componentwise in $C$. Also $zh\in C$ implies $\overline h\in S$. This gives the hypotheses of the linear algebra lemma for the approximate inverse. It gives $\mathcal S$ satisfying the stated inclusion and
\[
 |V|\,|\mathcal B|+|\mathcal S|
   \le |H\times\mathcal E|\dim_\mathbf k S
   \le (1+\theta)|V|\,|\mathcal B|.
\]
The last step uses $|H\times\mathcal E|=|V|\,|\mathcal P|$ and the edge dimension bound. Rearranging gives the bound on $|\mathcal S|$.\leanfn{See \href{https://github.com/lsernau/grh-lean/blob/72fead102f663d5c7a407964a13dc0cf9d977207/grh/src/New/UpperBound/ElementaryDivisors/ElementaryDivisorTheorem.lean\#L413}{The pairing exponent is at least two}, \href{https://github.com/lsernau/grh-lean/blob/72fead102f663d5c7a407964a13dc0cf9d977207/grh/src/New/UpperBound/ElementaryDivisors/ElementaryDivisorTheorem.lean\#L1011}{Precision of the approximate inverse and additional generators}.}
\end{proof}

\thmsource{Lean source: \leansourcelink{src/New/UpperBound/ElementaryDivisors/ElementaryDivisorTheorem.lean}{1066}.}
\begin{lemma}[Pairings and colengths of edge restriction]
\label{lem:a-edge-pairing-colength}
\emitthmsource
Let $T$ be a Frobenius cycle from a local ring $W$ to a valuation ring $O$ of characteristic zero, of cycle length $\nu$, with $p$ prime and a nonunit in $O$. Write $w=\sqrt p\in O$. Let $P$ be the Frobenius power of the cycle and $q=P^\nu$ the size of the residue field of $W$. Take rank $k$, degree $n$ in each copy, first block size $r$, and real $\theta$ satisfying
\[
 0<2r\le k,\qquad 0<\theta\le1,\qquad n^4\le q,
 \qquad \left(\frac{\nu k\log q}{\theta}\right)^5 P^{3r}\le n .
\]
Let $\mathcal B$ index monomials of degree $n$ in each of $\nu$ copies of $k$ variables, and put
\[
  c_0=k^2\log_p q,\qquad
  c_{\rm div}=\nu n\left(\frac1{P^{2r}}+\frac{r^2}{kP}\right).
\]
Let $\mathcal P$ be the coset set of the quasiparabolic of block sizes $r,k-r$, the subgroup of $\mathrm{GL}_k(W)$ whose upper right block reduces to zero. Take a finite set $V,H,\mathcal E$. For $e\in H\times\mathcal E$ give the starting points $a(e),a'(e)\in V$ of the original and dual edges and the respective basis changes $g(e),g'(e)\in\mathrm{GL}_k(W)$. Suppose the maps
\[
 e\longmapsto\bigl(a(e),[g(e)^{-1}]\bigr),\qquad
 e\longmapsto\bigl(a'(e),[g'(e)^{-1}]\bigr)
\]
both surject onto $V\times\mathcal P$ and $|H\times\mathcal E|=|V|\,|\mathcal P|$. Let $\rho$ denote the degree-$n$ representation of $T$ and $D_T$ the depth matrix with diagonal entries $w^{\operatorname{dep}(\beta)}$. Define the edge restriction $E,E'$ by
\[
 (EF)(e)=D_T\rho(g(e))F(a(e)),\qquad
 (E'F)(e)=D_T\rho(g'(e))F(a'(e)).
\]
Suppose weight $b,b'$ satisfies $b_\beta b'_\beta=p^{\nu c_0}$, and $\rho(g)$ and $\rho(g^\top)$ are adjoint for the pairing of weight $b$. For each original edge $e\in\mathcal E$ choose an integral invertible auxiliary transformation $u_e$. Let $P_b$ be the block operator applying $u_e$ and $\operatorname{diag}(b)$ on each edge. Put $\mathcal L=O^{V\times\mathcal B}$ and $A=E'^\top P_bE$. Then for every integer $s\ge0$ there exists $\lambda\in\mathbb N$ such that
\[
 \lambda\le(3c_{\rm div}+\theta s)\,|V|\,|\mathcal B|,
 \qquad
 \operatorname{ColengthDvd}\bigl(\mathcal L,A\mathcal L+p^s\mathcal L;p^\lambda\bigr).
\]
Here colength divisibility means that each finite family in $\mathcal L$ can be sent into the indicated submodule by an integral matrix whose determinant divides $p^\lambda$.
\end{lemma}
\begin{proof}
Put $d_{\mathcal L}=|V|\,|\mathcal B|$ and $K=\lfloor c_{\rm div}|\mathcal B|\rfloor$. Apply the starting-function basis lemma to each side to get $E=Qd$ and $E'=Q'd'$, where $Q,Q'$ each has a left inverse, their images are the starting-function modules on the two sides, and
\[
 \det d\mid w^{|V|K},\qquad
 \det d'\mid w^{|V|K},
\]
Hence $\det d'\det d\mid p^{|V|K}$. If $\psi$ lies in the starting-function module for the original edges, there exists $a,F$ with $w^a\psi=EF$. Fix edge $e$ and put $G=\rho(g(e))F(a(e))$. The original edges run over all parabolic cosets for each starting point, so $w^a\mid(\Delta\rho(\gamma)G)_\beta$ for every $\gamma,\beta$. Thus $\psi(e)=w^{-a}\Delta G$ lies in the intersection $\mathcal D$ of the halfway lattices. The dual edges are treated in the same way. The images of $Q,Q'$ therefore satisfy the edge congruence conditions at each edge. The pairing matrix factors as
\[
 A=d'^\top A_0d,\qquad A_0=Q'^\top P_bQ .
\]
Take the congruence conditions of precision $5\nu c_0$ for the edge data and the corresponding averaging operator. The additional-generator lemma for edge pairings gives a finite family $\mathcal S\subseteq\mathcal L$ satisfying\leanfn{See \href{https://github.com/lsernau/grh-lean/blob/72fead102f663d5c7a407964a13dc0cf9d977207/grh/src/New/UpperBound/ElementaryDivisors/Edges.lean\#L578}{Values of starting functions in the intersection of the halfway lattices}, \href{https://github.com/lsernau/grh-lean/blob/72fead102f663d5c7a407964a13dc0cf9d977207/grh/src/New/UpperBound/ElementaryDivisors/ElementaryDivisorTheorem.lean\#L1095}{Basis changes at the endpoints and generators for the middle factor}.}
\[
 |\mathcal S|\le\theta d_{\mathcal L},\qquad
 p^{2\nu c_0}\mathcal L\subseteq A_0\mathcal L+\langle\mathcal S\rangle_O .
\]

To combine the bounds, first $p^{2\nu c_0}\mathcal L\subseteq A_0\mathcal L+\langle\mathcal S\rangle_O$ gives colength dividing $(p^{2\nu c_0})^{d_{\mathcal L}}$. The image of matrix $d$ has colength in $\mathcal L$ dividing $\det d$. Map it by $A_0$ and add $\langle\mathcal S\rangle_O$ to reach $A_0d\mathcal L+\langle\mathcal S\rangle_O$, multiplying the colength by $\det d$. Next multiply each additional generator by $p^s$ to replace it by $p^s\mathcal L$ at the cost of another factor $(p^s)^{|\mathcal S|}$. Finally map by $d'^\top$ and enlarge $p^sd'^\top\mathcal L$ to $p^s\mathcal L$. The preceding colength for reaching $d'^\top\mathcal L$ divides $\det d'$. The product of the divisibility bounds is therefore\leanfn{See \href{https://github.com/lsernau/grh-lean/blob/72fead102f663d5c7a407964a13dc0cf9d977207/grh/src/New/UpperBound/ElementaryDivisors/ElementaryDivisorTheorem.lean\#L232}{Multiplication of colengths under the pairing factorization}, \href{https://github.com/lsernau/grh-lean/blob/72fead102f663d5c7a407964a13dc0cf9d977207/grh/src/New/Reduction/Lattices/Colength.lean\#L199}{The colength of the additional generators}.}
\[
 \det d'\det d\;p^{2\nu c_0 d_{\mathcal L}+s|\mathcal S|}.
\]

Thus we can take $\lambda=|V|K+2\nu c_0 d_{\mathcal L}+s|\mathcal S|$. The degree condition implies $\nu c_0\le c_{\rm div}$. Substitute $K\le c_{\rm div}|\mathcal B|$ and $|\mathcal S|\le\theta d_{\mathcal L}$ to obtain\leanfn{See \href{https://github.com/lsernau/grh-lean/blob/72fead102f663d5c7a407964a13dc0cf9d977207/grh/src/New/UpperBound/ElementaryDivisors/ElementaryDivisorTheorem.lean\#L561}{Upper bound for the polynomial pairing exponent}, \href{https://github.com/lsernau/grh-lean/blob/72fead102f663d5c7a407964a13dc0cf9d977207/grh/src/New/UpperBound/ElementaryDivisors/ElementaryDivisorTheorem.lean\#L1080}{Colength formula for the edge restriction}.}
\[
 \lambda\le c_{\rm div}d_{\mathcal L}
        +2c_{\rm div}d_{\mathcal L}+\theta s\,d_{\mathcal L}.
\]
\end{proof}

\thmsource{Lean source: \leansourcelink{src/New/UpperBound/ElementaryDivisors/ElementaryDivisorTheorem.lean}{1158}.}
\begin{theorem}[An elementary divisor bound for the edge Hecke operator]
\label{thm:a-edge-elementary-divisors}
\emitthmsource
Use the cycle, size conditions, and notation of \statementref{Lemma}{lem:a-edge-pairing-colength}. Let $X$ be a finite vertex set and $H$ the basis index set of the free auxiliary module at each vertex, with $J:I\simeq\mathcal B\times H$. Index the edges entering $x\in X$ by $c\in\mathcal P$. Edge $e=(x,c)$ starts at $a(e)$, has source and target basis changes $\gamma_e$ and $h_c$, and acts on the auxiliary module by $u_e\in\mathrm{GL}_H(O)$. Suppose $[h_c^{-1}]=c$ and that $a(e)$ and $[(\gamma_e^\top)^{-1}]$ uniquely determine each edge. Define the \emph{edge Hecke operator} $A_T$ by sending the starting value along edge $e=(x,c)$ via
\[
 \bigl(\rho(\gamma_e)D_T^2\rho(h_c)\bigr)^\top\otimes u_e
\]
to the target, then summing over the edges entering $x$ and reordering the coordinates by $J$.

Suppose $f:(X\to O^I)\to\mathcal M$ and $Y:\mathcal M\to\mathcal M$ are $O$-linear maps satisfying $Yf=fA_T$, and write $L=\operatorname{im}f$. Then $YL\subseteq L$, and for every integer $s\ge0$ there exists $\lambda\in\mathbb N$ such that
\[
 \operatorname{ColengthDvd}(L,YL+p^sL;p^\lambda),
 \qquad
 \lambda\le (3c_{\rm div}+\theta s)\,|X|\,|I|.
\]
\end{theorem}
\begin{proof}
The size conditions give $1<n$, which together with $n^4\le q$ gives $n<q$. Choose a monomial pairing weight $b,b'$ of degree $n$ satisfying $b_\beta b'_\beta=p^{\nu c_0}$ and $\operatorname{diag}(b)\rho(g^\top)=\rho(g)^\top\operatorname{diag}(b)$. Put $V=H\times X$. Regard edge $e=(x,c)$ and auxiliary basis index $h\in H$ together as an edge coordinate. Take the original starting point and basis change to be $(h,a(e))$ and $\gamma_e^\top$, and the dual ones to be $(h,x)$ and $h_c$. The uniqueness assumption makes $(x,c)\mapsto (a(e),[(\gamma_e^\top)^{-1}])$ an injection of the finite set $X\times\mathcal P$ into itself and hence a surjection. The corresponding dual map is $(x,c)\mapsto(x,[h_c^{-1}])=(x,c)$. After adding $H$, the two sides still satisfy the surjectivity and cardinality conditions in the lemma.\leanfn{See \href{https://github.com/lsernau/grh-lean/blob/72fead102f663d5c7a407964a13dc0cf9d977207/grh/src/New/UpperBound/ElementaryDivisors/Edges.lean\#L458}{Pairing weights and adjointness}, \href{https://github.com/lsernau/grh-lean/blob/72fead102f663d5c7a407964a13dc0cf9d977207/grh/src/New/UpperBound/ElementaryDivisors/ElementaryDivisorTheorem.lean\#L1148}{Exhausting the cosets at the starts of original edges}.}

Write $B_b=\operatorname{diag}(b)$. Since $D_T$ and $B_b$ are diagonal, the two halfway depth matrices on an edge multiply to give
\[
  (D_T\rho(h_c))^\top B_b(D_T\rho(\gamma_e^\top))
    =\bigl(\rho(\gamma_e)D_T^2\rho(h_c)\bigr)^\top B_b .
\]
The auxiliary factor is still $u_e$. Summing over the edges gives the following identity for $A=E'^\top P_bE$ in \statementref{Lemma}{lem:a-edge-pairing-colength}.
\[
  \operatorname{regroup}(AF)
     = A_T\bigl(\operatorname{regroup}(B_bF)\bigr)
 \qquad\text{for every }F\in O^{V\times\mathcal B},
\]
Here $\operatorname{regroup}:O^{V\times\mathcal B}\simeq X\to O^I$ is induced by $J$, and $B_b$ acts by the monomial indices at each vertex.\leanfn{See \href{https://github.com/lsernau/grh-lean/blob/72fead102f663d5c7a407964a13dc0cf9d977207/grh/src/New/UpperBound/ElementaryDivisors/ElementaryDivisorTheorem.lean\#L776}{The depth matrix identity on one edge}, \href{https://github.com/lsernau/grh-lean/blob/72fead102f663d5c7a407964a13dc0cf9d977207/grh/src/New/UpperBound/ElementaryDivisors/ElementaryDivisorTheorem.lean\#L807}{Reordering after summing over all edges}.}

Take $\pi=f\circ\operatorname{regroup}$. By $Yf=fA_T$, its image $L$ is stable under $Y$. The identity also gives $\pi(AF)=Y\pi(B_bF)$ and hence $\pi(A\mathcal L)\subseteq YL$. Since $\pi$ surjects onto $L$, lift any finite family in $L$ to one in $\mathcal L$. Correct the lifted family by the matrix from \statementref{Lemma}{lem:a-edge-pairing-colength} and apply $\pi$. The corrected family then lies in $YL+p^sL$ and the determinant condition is preserved. Finally,
\[
 |V|\,|\mathcal B|=|H|\,|X|\,|\mathcal B|=|X|\,|I|,
\]
This gives the stated upper bound for the same $\lambda$.\leanfn{See \href{https://github.com/lsernau/grh-lean/blob/72fead102f663d5c7a407964a13dc0cf9d977207/grh/src/New/UpperBound/ElementaryDivisors/ElementaryDivisorTheorem.lean\#L1166}{The elementary divisor theorem for the edge Hecke operator}.}
\end{proof}

\thmsource{Lean source: \leansourcelink{src/New/UpperBound/Divisibility/MarkOperator.lean}{336}.}
\begin{lemma}[Colength of the Hecke operator for one marking]
\label{lem:a-mark-single-colength}
\emitthmsource
Fix a $p$-adic embedding. Let $\Lambda$ be a marking at $p$, with $w\in\Lambda$, and let $T_w=p^{A_w}H_{w,r_p}$ be the scaled Hecke operator, where $r_p=\lceil\log_p m\rceil$ and $A_w=a\,\#\{\tau\in\Phi_a:w_\tau=w\}$. Write $M_0$ for the local integral module of the full source space and let $u_0=\dim_{\mathbb C}U_0$. Then $T_w$ preserves $M_0$. For each $j_0\in\mathbb N$ there exists $i_0\in\mathbb N$ such that every finite family in $M_0$ can be sent into $T_wM_0+p^{j_0}M_0$ by an integral matrix with determinant dividing $p^{i_0}$, where
\[
 i_0\le
 \left(\frac{12A_w}{mp}+\frac{j_0}{m^2f}\right)u_0 .
\]
\end{lemma}
\begin{proof}
Write $o=\operatorname{ord}(\rho_p)$, $o_0=\operatorname{ord}(\rho_p\bmod q_0)$, and $h=o/o_0$, with $\rho_p$ as in \statementref{Lemma}{lem:a-exceptional-order}. Choose a place $\tau_0\in\Phi_a$ inducing $w$. The decomposition group of $E/\mathbb Q$ at $w$ is generated by Frobenius $c\rho_p$, with $c$ complex conjugation. Thus $\tau$ induces $w$ or $\bar w$ exactly when there exists $j$ with $\tau=\rho_p^j\tau_0$. These places lie in $\Phi_a$ exactly when $\rho_p^j\equiv1\pmod{q_0}$, or $o_0\mid j$. Thus the places of $\Phi_a$ inducing $w$ or $\bar w$ are precisely
\[
 \tau_i=\rho_p^{o_0 i}\tau_0\qquad (i\in\mathbb Z/h\mathbb Z).
\]
The places are distinct, so $h\le f$. Since $o_0$ is odd, even $i$ corresponds to $w$ and odd $i$ to $\bar w$. Since $m\mid o$, the number $h$ is even, and hence $A_w=ah/2$.\leanfn{See \href{https://github.com/lsernau/grh-lean/blob/72fead102f663d5c7a407964a13dc0cf9d977207/grh/src/New/UpperBound/Divisibility/MarkOperator.lean\#L157}{Cyclotomic orbit inducing the conjugate pair}.}

Regard each $\tau_i$ as a complex embedding. When $i$ is odd, first compose $\tau_i$ with $c$, then compose with the fixed $p$-adic embedding. The resulting maps all induce the same place $w$. Each therefore maps the valuation ring of $w$ into the fixed $p$-adic valuation ring. Each successive map is obtained from the previous one by precomposing with the $o_0$-th Frobenius at $w$. Since $p$ is unramified in $E$, the element $p$ itself is a uniformizer at $w$. Modulo $p$ the adjacent maps differ by the $P_*=p^{o_0}$-th power, and the residue field has size $Q_*=p^o=P_*^h$. These maps with their alternating parities form the Frobenius cycle required by the elementary divisor theorem.\leanfn{See \href{https://github.com/lsernau/grh-lean/blob/72fead102f663d5c7a407964a13dc0cf9d977207/grh/src/New/UpperBound/Divisibility/MarkOperator.lean\#L196}{The Frobenius cycle in the valuation ring}.}

In the coefficient valuation ring, $p$ is a nonunit and $\sqrt p$ is integral. Put $\theta=(m^2f)^{-1}$. Applying the elementary divisor theorem requires checking $0<r_p$, $2r_p\le k$, $0<\theta\le1$, and
\[
 a^4\le Q_*,
 \qquad
 \left(\frac{h k\log Q_*}{\theta}\right)^5P_*^{3r_p}\le a.
\]
The first three follow from the admissible scales and $1\le o_0\le r_p$. Also $m\mid o$ and $p\ge13$ give $Q_*\ge p^m\ge e^m$, and $a^4\le e^{5\ell}\le e^m$. Write $l_0=3^{\lfloor\log_3 s\rfloor+4}$. By $r_p\le l_0$, $P_*^{3r_p}\le x_1^{3l_0^2}\le\sqrt b$. Using $p<x_1$ gives $\log Q_*\le o_0h\log x_1$. Together with the scale bounds
\[
 \left(\frac{h k(o_0h\log x_1)}{\theta}\right)^5\sqrt b\le a,
\]
this proves both size conditions.\leanfn{See \href{https://github.com/lsernau/grh-lean/blob/72fead102f663d5c7a407964a13dc0cf9d977207/grh/src/New/UpperBound/Divisibility/MarkOperator.lean\#L43}{Size conditions for the Frobenius cycle}.}

Choose a representative integral above $p$ for every double coset of the source level. Triviality of rational stabilizers identifies $M_0$ with the product of integral Weyl coordinates at these representatives. On the orbit $\{\tau_i\}$ the Weyl part of the source weight is a tensor product of $h$ symmetric powers of degree $a$. The weights at other places form the auxiliary factor. In the right-coset expansion of the Hecke operator, each rational edge matrix is integral invertible on the auxiliary factor. On the orbit factor, the matrix factors into integral basis changes at the ends and a depth operator in the middle, with two complementary depth scalings for the even and odd embeddings. Multiplying by $p^{A_w}$ turns the Hecke operator into an integral edge operator. \statementref{Theorem}{thm:a-edge-elementary-divisors} applies with degree $a$, first block size $r_p$, and Frobenius power $P_*$. It gives preservation of $M_0$ and a colength exponent at most
\[
 (3c_T+\theta j_0)u_0,\qquad
 c_T=h a\left(\frac{1}{P_*^{2r_p}}+\frac{r_p^2}{kP_*}\right),
\]
Here $u_0$ is the number of double cosets times the dimension of the Weyl module.\leanfn{See \href{https://github.com/lsernau/grh-lean/blob/72fead102f663d5c7a407964a13dc0cf9d977207/grh/src/New/UpperBound/Divisibility/HeckeEdges.lean\#L725}{Source coordinates and the edge Hecke operator}, \href{https://github.com/lsernau/grh-lean/blob/72fead102f663d5c7a407964a13dc0cf9d977207/grh/src/New/UpperBound/ElementaryDivisors/ElementaryDivisorTheorem.lean\#L1166}{Elementary divisor bound for the edge operator}.}

Finally $P_*\ge p$, $p^{r_p}\ge m$, and $r_p\ge1$ give $P_*^{2r_p}\ge mp$. The scales also give $mr_p^2\le k$ and hence $r_p^2/(kP_*)\le1/(mp)$. Thus
\[
 3c_T\le \frac{6ha}{mp}=\frac{12A_w}{mp}.
\]
Substituting in the elementary divisor bound gives the result.\leanfn{See \href{https://github.com/lsernau/grh-lean/blob/72fead102f663d5c7a407964a13dc0cf9d977207/grh/src/New/UpperBound/Divisibility/MarkOperator.lean\#L108}{Simplifying the elementary divisor constant}.}
\end{proof}

\thmsource{Lean source: \leansourcelink{src/New/UpperBound/Divisibility/LocalDivisibility.lean}{326}.}
\begin{lemma}[Colength of a product of Hecke operators at the marked places]
\label{lem:a-exceptional-hecke-colength}
\emitthmsource
Fix a $p$-adic embedding and let $M_0$ be the local integral module of the full source space $U_0$, with $u_0=\dim_{\mathbb C}U_0$. Let $\Lambda$ be a marking at $p$, with $A_w=a\,\#\{\tau\in\Phi_a:w_\tau=w\}$ and $r_p=\lceil\log_p m\rceil$. For $w\in\Lambda$ write $T_w$ for the $r_p$-th Hecke operator multiplied by $p^{A_w}$. Let $Y_\Lambda$ be their ordered product and put $j_p=\lceil2mfa/p\rceil$ and $d_p=fa\ell^2/p$. Then $Y_\Lambda^m$ preserves $M_0$, and there is $c\in\mathbb N$ with $c\le d_pu_0$ such that every finite family in $M_0$ can be sent by an integral matrix into
\[
 Y_\Lambda^m M_0+p^{j_p}M_0,
 \qquad\text{where the determinant of the matrix divides }p^c.
\]
\end{lemma}
\begin{proof}
By \statementref{Lemma}{lem:a-mark-single-colength}, $T_w$ preserves $M_0$. There is a nonnegative integer $c_w$ such that each finite family can be sent into $T_w M_0+p^{j_p}M_0$ by an integral matrix with determinant dividing $p^{c_w}$, and\leanfn{See \href{https://github.com/lsernau/grh-lean/blob/72fead102f663d5c7a407964a13dc0cf9d977207/grh/src/New/UpperBound/Divisibility/MarkOperator.lean\#L334}{Hecke colength for one marking}.}
\[
 c_w\le
  \left(\frac{12A_w}{mp}+\frac{j_p}{m^2f}\right)u_0 .
\]

First consider how to combine bounds under products. Suppose operators $V_1,V_2$ preserving $M_0$ have exponent bounds $c_1,c_2$ and the same allowed error $p^{j_p}M_0$. For a finite family $\boldsymbol v$ take an integral matrix $A$ such that
$A\boldsymbol v=V_1\boldsymbol x+p^{j_p}\boldsymbol e$ with each member of $\boldsymbol x,\boldsymbol e$ in $M_0$. Then take an integral matrix $B$ for $\boldsymbol x$ so that the members of $B\boldsymbol x$ lie in $V_2M_0+p^{j_p}M_0$. Since $V_1$ preserves $M_0$ and commutes with the action of the coefficient matrices,
\[
 BA\boldsymbol v=V_1(B\boldsymbol x)+p^{j_p}B\boldsymbol e
    \ \in\ (V_1V_2M_0+p^{j_p}M_0)^{\#\boldsymbol v}.
\]
Also $\det(BA)=(\det B)(\det A)\mid p^{c_1+c_2}$. Applying this argument successively in the product order gives an exponent at most $m\sum_{w\in\Lambda}c_w$.\leanfn{See \href{https://github.com/lsernau/grh-lean/blob/72fead102f663d5c7a407964a13dc0cf9d977207/grh/src/New/Reduction/Lattices/Colength.lean\#L83}{Additivity of colength exponents under products}.}

Each $w\in\Lambda$ is induced by a place of $\Phi_a$, so $\#\Lambda\le f$. The subsets of $\Phi_a$ corresponding to the different $w$ are disjoint, hence $\sum_{w\in\Lambda} A_w\le af$. Using $j_p\le 2mfa/p+1$ gives
\[
 \begin{aligned}
 m\sum_{w\in\Lambda}
       \left(\frac{12A_w}{mp}+\frac{j_p}{m^2f}\right)
 &\le \frac{12af}{p}+\frac{1}{m}\left(\frac{2mfa}{p}+1\right)\\
 &=\frac{14fa+p/m}{p}
 \le \frac{fa\ell^2}{p}
 =d_p .
 \end{aligned}
\]
The penultimate inequality uses $p<x_1\le a$, $f,m\ge1$, and $\ell=\log b\ge128$. The last equality uses the exceptional-prime condition $p\not\equiv\pm1\pmod m$. Multiply by $u_0$ and use the exponent bound for one operator to obtain the required $c$.\leanfn{See \href{https://github.com/lsernau/grh-lean/blob/72fead102f663d5c7a407964a13dc0cf9d977207/grh/src/New/UpperBound/Divisibility/LocalDivisibility.lean\#L285}{Computing the cost of the product}.}
\end{proof}

\section{Glued tensors, transfer operators, and local gains}
\label{sec:a-local-divisibility}

Expanding the glued tensor into tensor products of translated vacua lets us estimate the Weyl gauge along the chain, one lift at a time. At a marked place, Hecke transfer moves powers of the Hecke operators from the source forms to the kernel data. \statementref{Lemma}{lem:a-charged-schur-jet} of \chapref{chap:derivatives} then bounds the kernel derivatives.

\thmsource{Lean source: \leansourcelink{src/New/UpperBound/Divisibility/ChainDatumExpansion.lean}{324}.}
\begin{lemma}[Translation expansion of the glued tensor]
\label{lem:a-glued-translate-expansion}
\emitthmsource
Let $\Psi$ be chain data for twist $\vartheta$, with $\pi_N\bar\pi_N=N$, and write $q_\vartheta$ for the twist level norm. For a tuple $t=(t_j)_j$ of translations, write $\mathcal U(t)$ for the tensor of the vacua translated at the separate steps and $\mathcal G(t)$ for the tensor obtained by gluing adjacent steps in pairs. Put
\[
 \mathcal Y=\prod_{0\le i<m/2}\prod_{w\mid\pi_N}
                \mathbb F_N^{\,n_{2i+1}}.
\]
Each $y\in\mathcal Y$ defines a translation tuple $g(y)$ supported above $N$. At step $2i$ the last column has value $-\widetilde y_{i,w}/N$ at $w\mid\pi_N$, where $\widetilde y_{i,w}$ is a chosen integral lift of $y_{i,w}$. The last row at step $2i+1$ is its conjugate transpose, and all other entries are zero. Then
\[
 \sum_t\Psi(t)\mathcal G(t)
  =\sum_{t\in\operatorname{supp}\Psi}\Psi(t)
     \sum_{y\in\mathcal Y}\mathcal U\bigl(t+g(y)\bigr),
 \qquad
 \#\mathcal Y\le N^{fmn_m/2}.
\]
\end{lemma}
\begin{proof}
Each $t\in\operatorname{supp}\Psi$ is supported above $Mq_\vartheta$. Since $N$ and $Mq_\vartheta$ are coprime, $t_j$ is zero at every place above $N$. The column translation in the gluing is supported above $\pi_N$, and the conjugate row translation is supported above $\bar\pi_N$. Write $z_j$ for this translation at a given step. The Heisenberg product formula gives the phase from $\overline{t_j}\cdot z_j$. Conjugation preserves the places above $N$, so the pairing vanishes at each place. Two successive translations thus equal the single translation by $t_j+z_j$.\leanfn{See \href{https://github.com/lsernau/grh-lean/blob/72fead102f663d5c7a407964a13dc0cf9d977207/grh/src/New/UpperBound/Divisibility/ChainDatumExpansion.lean\#L117}{Product of translations with disjoint supports}.}

For double step $i$, expand the gluing operator as a finite sum over $(y_{i,w})_{w\mid\pi_N}$. By the preceding argument the two factors in each term are the vacua translated by $t_{2i}+g(y)_{2i}$ and $t_{2i+1}+g(y)_{2i+1}$, respectively. Multilinearity gives $\mathcal Y$ as the index set over all double steps, so $\mathcal G(t)=\sum_y\mathcal U(t+g(y))$. Weight and sum with $\Psi(t)$ to get the identity. There are also $f$ places above $\pi_N$, each with residue field $\mathbb F_N$, hence
\[
 \#\mathcal Y=N^{f\sum_{i<m/2}n_{2i+1}}
       \le N^{fmn_m/2},
\]
Here we use monotonicity of the ranks, $n_{2i+1}\le n_m$.\leanfn{See \href{https://github.com/lsernau/grh-lean/blob/72fead102f663d5c7a407964a13dc0cf9d977207/grh/src/New/UpperBound/Divisibility/ChainDatumExpansion.lean\#L134}{Expansion of one double step}, \href{https://github.com/lsernau/grh-lean/blob/72fead102f663d5c7a407964a13dc0cf9d977207/grh/src/New/UpperBound/Divisibility/ChainDatumExpansion.lean\#L198}{Counting the gluing classes}, \href{https://github.com/lsernau/grh-lean/blob/72fead102f663d5c7a407964a13dc0cf9d977207/grh/src/New/UpperBound/Divisibility/ChainDatumExpansion.lean\#L338}{Expansion of the full tensor}.}
\end{proof}

\thmsource{Lean source: \leansourcelink{src/New/UpperBound/Divisibility/ChainDatumExpansion.lean}{352}.}
\begin{lemma}[Weak chain data preserving the composite lift]
\label{lem:a-weak-chain-replacement}
\emitthmsource
For any chain data $\Psi$ for twist $\vartheta$, there are weak chain data $\Psi'$ with the following property for each source form $u\in U_0$. The composite lift from gluing $\Psi$ equals the restriction to $Y_m$ of the linear combination with coefficients $\Psi'$ of the chains of single-step lifts.
\end{lemma}
\begin{proof}
Retain the translation tuples and gluing classes of \statementref{Lemma}{lem:a-glued-translate-expansion} and define the finitely supported coefficients
\[
  \Psi'(z)=
    \sum_{\substack{t\in\operatorname{supp}\Psi,\ y\in\mathcal Y\\
                    t+g(y)=z}}
           \Psi(t).
\]
If $z\in\operatorname{supp}\Psi'$, then $g(y)$ vanishes at every finite place outside $q_\vartheta N$. Hence $\bar\pi_M z_j=\bar\pi_M t_j$ is integral at each such place, the required support condition for weak chain data. For any $p\nmid q_\vartheta$ and any $p$-adic place, each $\Psi(t)$ is integral, so $\Psi'(z)$ is integral. When the twist has no finite ramification, for every automorphism $\sigma$ of the complex field the triangle inequality and the original bound for the chain data give
\[
 \sum_z|\sigma(\Psi'(z))|
 \le \#\mathcal Y\sum_t|\sigma(\Psi(t))|
 \le M^{3fmn_m^2}N^{fmn_m}
 \le e^{x^{60}} .
\]
The last expression is the logarithmic size condition for $M,N$ and $x$ in the admissible data.\leanfn{See \href{https://github.com/lsernau/grh-lean/blob/72fead102f663d5c7a407964a13dc0cf9d977207/grh/src/New/UpperBound/Divisibility/ChainDatumExpansion.lean\#L249}{Admissible support after adding the translations}, \href{https://github.com/lsernau/grh-lean/blob/72fead102f663d5c7a407964a13dc0cf9d977207/grh/src/New/UpperBound/Divisibility/ChainDatumExpansion.lean\#L357}{The coefficients and their size bound}.}

Collecting terms with the same $z=t+g(y)$ gives
$\sum_z\Psi'(z)\mathcal U(z)=\sum_t\Psi(t)\mathcal G(t)$.
By definition of the chain data, the right-hand tensor is fixed by the joint action of the source and intermediate levels, with the identity at the last level. Splitting each pair of consecutive factors, the associativity isomorphism for tensor products carries each level action on a double step to the tensor product of the single-step Weil actions. Thus the single-step tensor on the left is also jointly invariant. With the coefficient properties, this makes $\Psi'$ weak chain data.\leanfn{See \href{https://github.com/lsernau/grh-lean/blob/72fead102f663d5c7a407964a13dc0cf9d977207/grh/src/New/UpperBound/Divisibility/ChainDatumExpansion.lean\#L295}{Joint invariance for double-step and single-step tensors}.}

Fix $u$ and send the two equal tensors into the same linear lifting map, which applies the theta kernels in succession on a pure tensor. The left side then gives the sum of chains of single-step lifts weighted by $\Psi'$. The right side gives the original glued lift. Apply the restriction from $X_m$ to $Y_m$ to both sides.\leanfn{See \href{https://github.com/lsernau/grh-lean/blob/72fead102f663d5c7a407964a13dc0cf9d977207/grh/src/New/UpperBound/Divisibility/ChainDatumExpansion.lean\#L376}{The lift identity from the tensor identity}.}
\end{proof}

\thmsource{Lean source: \leansourcelink{src/New/UpperBound/Divisibility/ChainDatumExpansion.lean}{414}.}
\begin{lemma}[Vacuum projection and the weak chain conditions on the original coefficients]
\label{lem:a-weak-chain-vacuum-projection}
\emitthmsource
For any chain data $\Psi$ for twist $\vartheta$, the unglued tensor $\sum_t\Psi(t)\mathcal U(t)$ formed with the original coefficients is jointly invariant under the levels. The original coefficients $\Psi$ meet the weak chain conditions.
\end{lemma}
\begin{proof}
At each step take the integral vacuum projection above $N$, averaging integral Heisenberg translations at those places with normalized measure. For $t\in\operatorname{supp}\Psi$, each $t_j$ is zero there, so that average fixes the translated vacuum. With an additional translate $z_j$, the projection multiplies by the indicator that $z_j$ is integral at those places. This follows by averaging an additive character over the integral translation lattice and using self-duality of the standard lattice.

For a gluing class $y_i$ of double step $i$, a nonzero residue class forces a nonintegral entry of the form $-\widetilde y/N$ in the last column at step $2i$ at the corresponding place. This term is killed by the vacuum projection. The zero class gives the original pair of translated vacua. Thus the projection on each double step returns the original unglued data. Taking tensor products over the double steps and summing with $\Psi(t)$ gives $\sum_t\Psi(t)\mathcal U(t)$ as the projection of the full glued tensor.\leanfn{See \href{https://github.com/lsernau/grh-lean/blob/72fead102f663d5c7a407964a13dc0cf9d977207/grh/src/New/UpperBound/Divisibility/ChainDatumExpansion.lean\#L56}{Invariance of the translated vacuum}, \href{https://github.com/lsernau/grh-lean/blob/72fead102f663d5c7a407964a13dc0cf9d977207/grh/src/New/UpperBound/Divisibility/ChainDatumExpansion.lean\#L86}{Eliminating nonzero gluing classes}.}

The Weil action of each level preserves the integral translation lattice. Change variables in the averaging integral to see that the vacuum projection commutes with each level action, preserving joint invariance. Integrality of the original coefficients at all $p\nmid q_\vartheta$ is part of the chain data. Admissibility of the translations away from $q_\vartheta$ implies admissibility away from $q_\vartheta N$. In the absence of finite ramification of the twist, the original bound $M^{3fmn_m^2}$ is at most $e^{x^{60}}$ by the same logarithmic scale conditions. Thus $\Psi$ meets the weak chain conditions.\leanfn{See \href{https://github.com/lsernau/grh-lean/blob/72fead102f663d5c7a407964a13dc0cf9d977207/grh/src/New/UpperBound/Divisibility/ChainDatumExpansion.lean\#L390}{The projection commutes with the level action}, \href{https://github.com/lsernau/grh-lean/blob/72fead102f663d5c7a407964a13dc0cf9d977207/grh/src/New/UpperBound/Divisibility/ChainDatumExpansion.lean\#L421}{Weak chain conditions for the original coefficients}.}
\end{proof}

\thmsource{Lean source: \leansourcelink{src/New/UpperBound/Divisibility/ChainValuation.lean}{238}.}
\begin{lemma}[Multiplication of the gauges along the chain]
\label{lem:a-weak-chain-gauge}
\emitthmsource
Let $\Psi$ be weak chain data for twist $\vartheta$ and let $p$ be a prime coprime to $q_\vartheta$. Let $\varphi_j(t)$ be the finite kernel datum at step $j$, namely the vacuum translated by $t_j$. Choose gauge base points $d_j$ at the levels and linear operators $A_j$ on the kernel data so that $A_j$ commutes with the joint action of the preceding and following levels on each datum $\varphi_j(t)$ in the support. If for each $t\in\operatorname{supp}\Psi$ the kernel at step $j$ for the new datum $A_j\varphi_j(t)$ has gauge multiplier at most $p^{r_j}$ from $d_j$ to $d_{j+1}$, the $\Psi$-weighted lift for those new data sends source forms of gauge at most $1$ at $d_0$ to forms of gauge at most $p^{\sum_{j<m}r_j}$ at $d_m$. If $d_j=1$ and $p\ne N$, the same bound holds after restriction from $X_m$ to $Y_m$.
\end{lemma}
\begin{proof}
The gauge at $d_j$ means the gauge of $u(g)$ for every finite point $g$ in the Weyl lattice determined by $gd_j$. The commutation condition lets the tensor map $\bigotimes_{j<m}A_j$ carry the jointly invariant tensor to another jointly invariant tensor. Thus the compact-level averages at each step combine in the class integral. At each stage $X_j$ with $j<m$, choose representatives $c_j$ of the rational-point and level double cosets. Every level is inside the principal congruence level and $n_j\le n_m$. The stabilizer criterion therefore gives trivial automorphism groups, so the weights in the class integral are all $1$. Thus at any finite point $g_m$ of $X_m$ the weighted lift chain has value
\[
 \sum_{(c_0,\ldots,c_{m-1})}\ \
 \sum_{t\in\operatorname{supp}\Psi}
 \Psi(t)\,K^{A,t}_{m-1}(g_m,c_{m-1})\cdots
 K^{A,t}_0(c_1,c_0)\,u(c_0),
\]
Here $K^{A,t}_j(g_{j+1},g_j)$ is the linear action of the single-step kernel for the new data $A_j\varphi_j(t)$ at the two specified points.\leanfn{See \href{https://github.com/lsernau/grh-lean/blob/72fead102f663d5c7a407964a13dc0cf9d977207/grh/src/New/UpperBound/Divisibility/ChainValuation.lean\#L156}{The class expansion at each stage and the trivial stabilizers}, \href{https://github.com/lsernau/grh-lean/blob/72fead102f663d5c7a407964a13dc0cf9d977207/grh/src/New/UpperBound/Divisibility/ChainValuation.lean\#L248}{Joint invariance of the new data}.}

In one term start from the gauge bound $1$ for $u(c_0)$. Each kernel increases it by a factor at most $p^{r_j}$, giving gauge at most $\prod_{j<m}p^{r_j}$ in the lattice at $g_md_m$ for the product. The coefficient $\Psi(t)$ is integral at the given $p$-adic place, so it does not increase the bound. The finite sum has the same bound by the ultrametric inequality. This holds for every $g_m$.\leanfn{See \href{https://github.com/lsernau/grh-lean/blob/72fead102f663d5c7a407964a13dc0cf9d977207/grh/src/New/UpperBound/Divisibility/ChainValuation.lean\#L214}{Local gauges of products of kernels}.}

When $p\ne N$ the restricted value at $h$ is the value of $X_m$ at $\operatorname{diag}(h,1)$ times the integral branching matrix and a scalar formed from powers of embeddings of $\bar\pi_N$. The integral branching matrix maps the Weyl lattice into the target lattice. By $\pi_N\bar\pi_N=N$ the scalar is a unit at the $p$-adic place. Thus restriction does not increase the gauge with base point $1$.\leanfn{See \href{https://github.com/lsernau/grh-lean/blob/72fead102f663d5c7a407964a13dc0cf9d977207/grh/src/New/UpperBound/Divisibility/ChainValuation.lean\#L349}{Restriction to $Y_m$}, \href{https://github.com/lsernau/grh-lean/blob/72fead102f663d5c7a407964a13dc0cf9d977207/grh/src/New/UpperBound/Divisibility/LocalDivisibility.lean\#L146}{Stepwise gauge bound for the restricted lift}.}
\end{proof}

\thmsource{Lean source: \leansourcelink{src/New/UpperBound/Divisibility/LocalDivisibility.lean}{155}.}
\begin{lemma}[Standard data preserve integrality away from the level]
\label{lem:a-weak-chain-integral-away}
\emitthmsource
Let $\Psi$ be weak chain data for twist $\vartheta$ with $p\nmid q_\vartheta$ and $p\notin\{N,d\}$. Then the restricted lift for $\Psi$ sends each integral source form to one with $p$-adic gauge at most $1$ at every point.
\end{lemma}
\begin{proof}
For any $t\in\operatorname{supp}\Psi$ the support condition ensures that $\bar\pi_Mt_j$ is integral at every finite place away from $q_\vartheta N$, in particular above $p$. Choose $\delta=\bar\pi_M$ clearing the translation denominators. The vacuum translated by $t_j$ hence belongs to the $\delta$-standard data, the module over the valuation ring generated by translates of the vacuum whose translation vectors become integral above $p$ after multiplication by $\delta$.\leanfn{See \href{https://github.com/lsernau/grh-lean/blob/72fead102f663d5c7a407964a13dc0cf9d977207/grh/src/New/UpperBound/Divisibility/ChainValuation.lean\#L174}{Standard data from the support of weak chain data}.}

In the single-step gauge estimate for standard data, use extra Taylor scaling $\bar\delta$ to bound the multiplier by
\[
  \left(\frac{|\bar\delta|_p}{|d\bar\delta|_p}\right)^{b_{j+1}},
\]
Here $b_{j+1}$ counts the boxes in columns of height less than $n_{j+1}$ in the weight diagram at stage $j+1$. The absolute value is normalized by $|p|_p=p^{-1}$. By $\bar\delta\ne0$ and $p\ne d$ the ratio above is $1$, so we can take $r_j=0$ at every step. An integral source form has gauge at most $1$ at the identity base point. Apply \statementref{Lemma}{lem:a-weak-chain-gauge}.\leanfn{See \href{https://github.com/lsernau/grh-lean/blob/72fead102f663d5c7a407964a13dc0cf9d977207/grh/src/New/UpperBound/Divisibility/LocalDivisibility.lean\#L161}{Integrality of the chain away from the level}, \href{https://github.com/lsernau/grh-lean/blob/72fead102f663d5c7a407964a13dc0cf9d977207/grh/src/New/UpperBound/Kernels/ThetaLiftGauge.lean\#L298}{Single-step gauge for standard data}.}
\end{proof}

\thmsource{Lean source: \leansourcelink{src/New/UpperBound/Divisibility/LocalDivisibility.lean}{183}.}
\begin{lemma}[Unmarked divisibility at nonsplit primes]
\label{lem:a-weak-chain-unmarked-valuation}
\emitthmsource
Let $\Psi$ be weak chain data for twist $\vartheta$ with restricted lift $\mathcal L_\Psi$. Take $p<x_1$, $\chi_K(p)\ne1$, and $p\nmid q_\vartheta$, and put
\[
 d_p=\frac{fa\ell^2}{p}
  \begin{cases}
   m,&p\equiv\pm1\pmod m,\\
   1,&p\not\equiv\pm1\pmod m,
  \end{cases}
 \qquad j_p=\left\lceil\frac{2mfa}{p}\right\rceil .
\]
For each integral source form $u$ we have at every point
\[
 \operatorname{Gauge}_p\bigl(\mathcal L_\Psi(u)\bigr)
 \le p^{m C_p(\varnothing)+d_p-\gamma_p},
 \qquad
 \operatorname{Gauge}_p\bigl(\mathcal L_\Psi(p^{j_p}u)\bigr)
 \le p^{d_p-\gamma_p}.
\]
\end{lemma}
\begin{proof}
By $p<x_1$, we have $p\nmid MN$. For every $t$ in the support, $\bar\pi_M$ is a unit at all places above $p$. Thus $t_j$ is itself integral there and each kernel is $1$-standard. Let $D_{p,j}(S')$ be the predicted divisibility bound at step $j$ after deleting the columns of height $1$ on $S'$, including the compensation for deleting columns and the logarithmic error. The standard-data estimate at a nonsplit prime gives a single-step gauge multiplier at most $p^{-D_{p,j}(\varnothing)}$.\leanfn{See \href{https://github.com/lsernau/grh-lean/blob/72fead102f663d5c7a407964a13dc0cf9d977207/grh/src/New/UpperBound/Divisibility/ChainValuation.lean\#L174}{Standard data without denominator clearing}, \href{https://github.com/lsernau/grh-lean/blob/72fead102f663d5c7a407964a13dc0cf9d977207/grh/src/New/UpperBound/Divisibility/ValuationAtNonSplitPrimes.lean\#L66}{The single-step gain at a nonsplit prime}.}

Comparing the single-step predicted divisibility with the local gain gives
\[
 D_{p,j}(\varnothing)
 \ge \frac{\gamma_p}{m}-C_p(\varnothing)-\frac{fa\ell^2}{mp}
 \ge \frac{\gamma_p-d_p}{m}-C_p(\varnothing),
\]
The second inequality follows from $d_p\ge fa\ell^2/p$. The exponent from \statementref{Lemma}{lem:a-weak-chain-gauge} is $\sum_{j<m}-D_{p,j}(\varnothing)$ and so is at most $mC_p(\varnothing)+d_p-\gamma_p$.\leanfn{See \href{https://github.com/lsernau/grh-lean/blob/72fead102f663d5c7a407964a13dc0cf9d977207/grh/src/New/UpperBound/Divisibility/JetExponent.lean\#L504}{Single-step predicted divisibility and collision terms}, \href{https://github.com/lsernau/grh-lean/blob/72fead102f663d5c7a407964a13dc0cf9d977207/grh/src/New/UpperBound/Divisibility/LocalDivisibility.lean\#L190}{Divisibility along the unmarked chain}.}

For each fixed infinite place the collision terms use only exponents $e\ge1$, so sum to at most $\sum_{e\ge1}a/p^e=a/(p-1)\le 2a/p$. With at most $f$ places this gives $mC_p(\varnothing)\le j_p$. By linearity we can move $p^{j_p}$ to the output, multiplying its gauge by $p^{-j_p}$. Substitute in the first inequality to get the second.\leanfn{See \href{https://github.com/lsernau/grh-lean/blob/72fead102f663d5c7a407964a13dc0cf9d977207/grh/src/New/UpperBound/Divisibility/LocalDivisibility.lean\#L222}{An integer power that cancels all collisions}.}
\end{proof}

\thmsource{Lean source: \leansourcelink{src/Lib/Weil/Unramified/HeckeTransfer.lean}{495}. See also \href{https://github.com/lsernau/grh-lean/blob/72fead102f663d5c7a407964a13dc0cf9d977207/grh/src/Lib/Weil/Unramified/HeckeTransfer.lean\#L413}{Hecke operator identities for the Weil kernel}.}
\begin{lemma}[Transfer of Hecke polynomials for an unramified local pairing]
\label{lem:a-unramified-pairing-hecke}
\emitthmsource
We use the Hecke operator form of the unramified theta correspondence. See Howe~\cite{Howe1979} and Kudla~\cite{Kudla1986} for the representation-theoretic background. Let $w$ be a split finite place with residue field of size $q$. Take Hermitian spaces $X,X'$ of ranks $n$ and $n+\delta$ and a Weil action $\omega$ determined by two base-change characters unramified at $w$. Write $a,b$ for the characters and $\upsilon$ for the splitting characters. Let their Frobenius scalars be $\alpha=q^{n/2}a^\flat_n(\varpi_w)$ and $\beta=q^{(n+\delta)/2}b^\flat_{n+\delta}(\varpi_w)$, with $a^\flat_N(y)=a(y/\bar y)\upsilon(y)^{-N}$ and the analogous definition of $b^\flat_N$. Suppose the finite data $\psi$ are fixed by the integral vacuum projection on $w,\bar w$. Let $L$ be a bilinear lift on a linear data space containing the Weil orbit of $\psi$, satisfying the following for the source representation $\pi$ and target representation $\rho$.
\[
 \rho(h)L(\psi,v)=L\bigl(\omega(g,h)\psi,\pi(g)v\bigr),
\]
Then whenever $v$ is fixed by the source maximal compact at $w$, the lift $L(\psi,v)$ is fixed by the corresponding target maximal compact. Define the source Hecke polynomial by
\[
 P_{w,X}(U)=\sum_{i=0}^{n}(-U)^i q^{\binom i2}T_{w,i}^{X},
\]
Define the target polynomial likewise. For $v$ as above,
\[
 P_{w,X'}(U)L(\psi,v)=
 \prod_{j=0}^{\delta-1}(1-q^j\alpha U)\,
 L\bigl(\psi,P_{w,X}(q^\delta\alpha\beta^{-1}U)v\bigr).
\]
\end{lemma}
\begin{proof}
First compute the Hecke polynomial acting on $\psi$ in the local Weil model. The double-coset sum of index $i$ is indexed by $i$-dimensional subspaces of the residue vector space. For an integral matrix $x:X\to X'$, inverse translation by a representative $\gamma_{U_0}$ of a subspace $U_0$ gives an integral matrix exactly when $U_0\subseteq\ker\bar x$. Conversely, since $\gamma_{U_0}$ is integral, integrality after inverse translation implies integrality of the original matrix. The target-group calculation uses the transpose $\bar x^\top$. The character value at the uniformizer and the vacuum scaling contribute weights $\alpha^{\dim U_0}$ and $\beta^{\dim U_0}$, respectively. Thus both polynomial actions are multipliers on the local matrix variable.\leanfn{See \href{https://github.com/lsernau/grh-lean/blob/72fead102f663d5c7a407964a13dc0cf9d977207/grh/src/Lib/Weil/Unramified/HeckeTransfer.lean\#L305}{Subspace coset sums and vacuum translations}.}

Let $\bar x$ have rank $a_x$ and write $(Y;q)_d=\prod_{j=0}^{d-1}(1-q^jY)$. On integral matrices the multipliers for the target polynomial and the source polynomial at the conjugate place $\bar w$ are respectively
\[
 (\alpha U;q)_{n+\delta-a_x},\qquad
 (\beta U;q)_{n-a_x}.
\]
In fact, in the sum over subspace dimensions each $i$-dimensional subspace contributes $(-Y)^i q^{\binom i2}$. The finite-field binomial identity gives $(Y;q)_d$ after summing over all subspaces, where $d$ is the dimension of the kernel. Both multipliers vanish on nonintegral matrices. The factorization
\[
 (Y;q)_{n+\delta-a_x}
   =(Y;q)_\delta(q^\delta Y;q)_{n-a_x}
\]
gives the transfer identity on the local Weil data. The determinant characters on the integral vacuum are trivial on local units, so the maximal compacts on both sides fix $\psi$.\leanfn{See \href{https://github.com/lsernau/grh-lean/blob/72fead102f663d5c7a407964a13dc0cf9d977207/grh/src/Lib/Weil/Unramified/HeckeTransfer.lean\#L413}{Frobenius scalars and the product factorization}.}

View $L(\psi,\cdot)$ as a linear map from the source space to the target. Equivariance turns the action of $(g,h)$ on the datum into composition with $\rho(h)$ on the left and $\pi(g)^{-1}$ on the right. The target polynomial therefore acts on the image and the source polynomial on the input. Inversion in the source action also exchanges the conjugate place $\bar w$ with $w$. In a split unitary group, the $\bar w$-component is the inverse transpose of the $w$-component. Inverting the double-coset representatives thus matches the Hecke sums. For $v$ spherical at $w$, compact factors do not affect the values. Apply the local data identity to $v$ to get transfer and sphericity.\leanfn{See \href{https://github.com/lsernau/grh-lean/blob/72fead102f663d5c7a407964a13dc0cf9d977207/grh/src/Lib/Weil/Unramified/SplitHecke.lean\#L436}{The adjoint of the source Hecke polynomial for the pairing}, \href{https://github.com/lsernau/grh-lean/blob/72fead102f663d5c7a407964a13dc0cf9d977207/grh/src/Lib/Weil/Unramified/HeckeTransfer.lean\#L495}{Transfer through the bilinear lift}.}
\end{proof}

\thmsource{Lean source: \leansourcelink{src/New/UpperBound/Divisibility/ValuationAtNonSplitPrimes.lean}{466}.}
\begin{lemma}[An integral Hecke polynomial for one theta lift]
\label{lem:a-marked-hecke-polynomial}
\emitthmsource
Let $\Psi$ be weak chain data for twist $\vartheta$ with $p\nmid q_\vartheta$, and let $\Lambda$ be a marking at $p$ with $w\in\Lambda$. Denote by $q_w=z^2$ the residue field size of $w$, with $z$ a positive integer. Fix a $p$-adic embedding and put $\mathcal O_p=\{u\in\mathbb C:|\iota_p(u)|_p\le1\}$. Let $\Theta_e^x$ be the lift at step $e$ with translation datum $x$ in the support, and let $\delta=t+1$ be the rank difference. There is a polynomial $R_e(U)\in\mathcal O_p[U]$ independent of $x$ satisfying $R_e(0)=1$. Every input $\phi$ spherical at $w$ has spherical lift and satisfies
\[
 P_{w,X_{e+1}}(U)\Theta_e^x\phi
   =R_e(U)\,\Theta_e^x\bigl(P_{w,X_e}(z^\delta U)\phi\bigr).
\]
\end{lemma}
\begin{proof}
The marked place $w$ and its conjugate $\bar w$ are distinct. At both, $Mq_\vartheta N$ is a unit. The weak-chain support condition makes $x_e$ integral there, so the data are fixed by both vacuum projections. The extra tame data occur only at the chosen level and twist places, so the rank-twisted characters on both sides are unramified at $w$. Apply \statementref{Lemma}{lem:a-unramified-pairing-hecke} to $\Theta_e^x$, using Weil equivariance of the theta kernel at this step.\leanfn{See \href{https://github.com/lsernau/grh-lean/blob/72fead102f663d5c7a407964a13dc0cf9d977207/grh/src/New/UpperBound/Divisibility/ValuationAtNonSplitPrimes.lean\#L387}{Unramifiedness of the rank-twisted characters}, \href{https://github.com/lsernau/grh-lean/blob/72fead102f663d5c7a407964a13dc0cf9d977207/grh/src/New/UpperBound/Divisibility/ValuationAtNonSplitPrimes.lean\#L402}{The finite vacuum at marked places}.}

Write $\chi$ for the value of the base change of the step character at a uniformizer of $w$. It is a $p$-adic unit. By definition it is the character value at $a/\bar a$, where the finite idele $a$ is the uniformizer only at $w$. The character weight formula gives the $p$-adic absolute-value exponent as the difference between the total weights of the infinite places inducing $w$ and $\bar w$, which is zero. Since both consecutive ranks are even, the rank-twisted Frobenius scalars are $z^{n_e}\chi$ and $z^{n_{e+1}}\chi$. In the transfer formula the source variable is then scaled by
\[
 q_w^\delta\frac{z^{n_e}\chi}{z^{n_{e+1}}\chi}=z^\delta.
\]
The extra polynomial factor is
\[
 R_e(U)=\prod_{j=0}^{\delta-1}\bigl(1-q_w^jz^{n_e}\chi\,U\bigr),
\]
Its coefficients are integral and its constant term is $1$. It depends only on the character and place for this step, hence is independent of $x$.\leanfn{See \href{https://github.com/lsernau/grh-lean/blob/72fead102f663d5c7a407964a13dc0cf9d977207/grh/src/New/UpperBound/Divisibility/ValuationAtNonSplitPrimes.lean\#L420}{Norm of the character value at the uniformizer}, \href{https://github.com/lsernau/grh-lean/blob/72fead102f663d5c7a407964a13dc0cf9d977207/grh/src/New/UpperBound/Divisibility/ValuationAtNonSplitPrimes.lean\#L475}{The integral Hecke polynomial for a single step}.}
\end{proof}

\thmsource{Lean source: \leansourcelink{src/New/UpperBound/Divisibility/ValuationAtNonSplitPrimes.lean}{623}.}
\begin{lemma}[Coefficients and cost in powers of $p$ for one transfer step]
\label{lem:a-marked-hecke-one-step-cost}
\emitthmsource
Retain the preceding setting, put $r_p=\lceil\log_p m\rceil$ and $b_p=fr_pt$, and write $T_i^{(e)}$ for the index-$i$ Hecke operator at $w$. Spherical inputs have spherical images under $\Theta_e^x$. For each step $e$ there are integral coefficients $d_e(i,i')\in\mathcal O_p$ independent of $x$ such that, for $0\le i\le r_p$ and each $\phi$ spherical at $w$,
\[
 \sum_{i'=0}^{i}d_e(i,i')\,T_{i'}^{(e+1)}(\Theta_e^x\phi)
      =p^{b_p}\Theta_e^x(T_i^{(e)}\phi).
\]
\end{lemma}
\begin{proof}
By $R_e(0)=1$, the series $R_e(U)$ is invertible in $\mathcal O_p[[U]]$. Write $a_h=[U^h]R_e(U)^{-1}$. Multiply the single-step polynomial identity by $R_e(U)^{-1}$, extract the coefficient of $U^i$ and multiply by $(-1)^i$. Retaining only terms with $i'\le i$ gives
\[
 \sum_{i'=0}^{i}
    (-1)^{i+i'}a_{i-i'}q_w^{\binom{i'}2}\,T_{i'}^{(e+1)}(\Theta_e^x\phi)
  =z^{(t+1)i}q_w^{\binom i2}\,\Theta_e^x(T_i^{(e)}\phi)
 \quad (i\le n_e).
\]
All coefficients on the left belong to $\mathcal O_p$ and are independent of $x$. To compare coefficients in the identity for all $U$, apply any linear functional on the target. This gives a scalar polynomial identity over the infinite field $\mathbb C$. Inverting the formal power series completes the argument.\leanfn{See \href{https://github.com/lsernau/grh-lean/blob/72fead102f663d5c7a407964a13dc0cf9d977207/grh/src/New/UpperBound/Divisibility/ValuationAtNonSplitPrimes.lean\#L550}{Coefficients of the inverse power series}.}

The local degree of a marked place is even, so write it as $q_w=p^{2h}$. It is the order of $-p$ in the group of infinite places, hence $2h\le f$. Take $z=p^h$. The scaling on the right above simplifies to
\[
 z^{(t+1)i}q_w^{\binom i2}
  =p^{h i(t+i)}.
\]
The scales give $r_p\le t$ and $r_p\le k\le n_e$, hence when $i\le r_p$,
\[
 h i(t+i)\le2h\,r_pt\le fr_pt=b_p.
\]
Multiply the coefficients by the integer $p^{b_p-h i(t+i)}$ to obtain the required integral coefficients $d_e(i,i')$.\leanfn{See \href{https://github.com/lsernau/grh-lean/blob/72fead102f663d5c7a407964a13dc0cf9d977207/grh/src/New/UpperBound/Divisibility/ValuationAtNonSplitPrimes.lean\#L457}{Even local degree at a marked place}, \href{https://github.com/lsernau/grh-lean/blob/72fead102f663d5c7a407964a13dc0cf9d977207/grh/src/New/UpperBound/Divisibility/ValuationAtNonSplitPrimes.lean\#L608}{Integer power bound for single-step scaling}, \href{https://github.com/lsernau/grh-lean/blob/72fead102f663d5c7a407964a13dc0cf9d977207/grh/src/New/UpperBound/Divisibility/ValuationAtNonSplitPrimes.lean\#L634}{A uniform single-step transfer cost}.}
\end{proof}

\thmsource{Lean source: \leansourcelink{src/New/UpperBound/Divisibility/ValuationAtNonSplitPrimes.lean}{668}.}
\begin{lemma}[Iterating Hecke transfer along the lift chain]
\label{lem:a-marked-hecke-chain-transfer}
\emitthmsource
Keep the marked place and weak chain data, and write $C_e^x$ for the first $e$ lifts with data $x$. For each $e$ there are coefficients $c_e(i)\in\mathcal O_p$ independent of $x$ such that for every source form $\phi$ spherical at $w$, the image $C_e^x\phi$ is still spherical and
\[
 \sum_{i=0}^{r_p}c_e(i)\,T_i^{(e)}(C_e^x\phi)
     =p^{e b_p}\,C_e^x(T_{r_p}^{(0)}\phi).
\]
\end{lemma}
\begin{proof}
When $e=0$ take $c_0(i)=\mathbf 1_{i=r_p}$. Suppose the identity holds for $e$. The single-step sphericity result makes $C_{e+1}^x\phi=\Theta_e^x(C_e^x\phi)$ spherical. Apply $\Theta_e^x$ to the identity at step $e$ and multiply by $p^{b_p}$. Apply \statementref{Lemma}{lem:a-marked-hecke-one-step-cost} to each $i$ and exchange the two finite sums to give coefficient $i'$ as
\[
 c_{e+1}(i')=\sum_{i=i'}^{r_p}c_e(i)d_e(i,i').
\]
The terms are integral and independent of $x$, and so are the new coefficients. The right-hand exponent becomes $(e+1)b_p$, completing the induction.\leanfn{See \href{https://github.com/lsernau/grh-lean/blob/72fead102f663d5c7a407964a13dc0cf9d977207/grh/src/New/UpperBound/Divisibility/ValuationAtNonSplitPrimes.lean\#L679}{Preserving sphericity and collecting coefficients step by step}.}
\end{proof}

\thmsource{Lean source: \leansourcelink{src/New/UpperBound/Divisibility/ValuationAtNonSplitPrimes.lean}{758}.}
\begin{lemma}[A uniform Hecke transfer exponent for one marked place]
\label{lem:a-single-mark-hecke-transfer}
\emitthmsource
Retain the setting and write $E_p^{\rm tr}=r_pf(n_m+r_p)$. Each original datum $\varphi_j(x)$ is fixed by the target maximal compact at $w$. For each $j$ there are coefficients $c_i$ independent of $x$, all in $\mathcal O_p$ when $j<m$, such that every source form $\phi$ spherical at $w$ satisfies
\[
 \sum_{i=0}^{r_p} c_i\,T_i^{(j+1)}(C_{j+1}^x\phi)
    =p^{E_p^{\rm tr}} C_{j+1}^x(T_{r_p}^{(0)}\phi).
\]
\end{lemma}
\begin{proof}
At $w$ the original datum is the vacuum. The target maximal compact acts on it through an unramified determinant character, which takes the value $1$ on local units. Thus the datum is spherical. The preceding lemma for the chain up to stage $j+1$ supplies coefficients $c_{j+1}(i)$ and transfer exponent $(j+1)b_p$. Set
\[
 c_i=p^{E_p^{\rm tr}-(j+1)b_p}c_{j+1}(i)
\]
to obtain the identity. For $j<m$, since $n_m=k+m(t+1)$, we have $(j+1)t\le mt\le n_m$ and hence
\[
 (j+1)b_p\le fr_p n_m\le fr_p(n_m+r_p)=E_p^{\rm tr}.
\]
The additional power of $p$ is nonnegative, so $c_i$ is integral.\leanfn{See \href{https://github.com/lsernau/grh-lean/blob/72fead102f663d5c7a407964a13dc0cf9d977207/grh/src/New/UpperBound/Divisibility/ValuationAtNonSplitPrimes.lean\#L741}{The vacuum at a marked place is spherical}, \href{https://github.com/lsernau/grh-lean/blob/72fead102f663d5c7a407964a13dc0cf9d977207/grh/src/New/UpperBound/Divisibility/ValuationAtNonSplitPrimes.lean\#L770}{The transfer theorem for one marked place}.}
\end{proof}

\thmsource{Lean source: \leansourcelink{src/New/UpperBound/Divisibility/ValuationAtNonSplitPrimes.lean}{251}.}
\begin{lemma}[Counting boxes at a marked neighboring point]
\label{lem:a-marked-neighbor-box-count}
\emitthmsource
Let $\Lambda$ be a marking. Fix a $p$-adic embedding and write $S_\Lambda=\{\tau:w_\tau\in\Lambda\}$ and $S_\Lambda^-=\{\tau:\bar w_\tau\in\Lambda\}$. For each $w\in\Lambda$ take a subset $J_w$ of at most $r$ basis indices of $X_{j+1}$. If $\tau\in S_\Lambda$ put $J_\tau=J_{w_\tau}$, if $\tau\in S_\Lambda^-$ put $J_\tau=J_{\bar w_\tau}$, and otherwise put $J_\tau=\varnothing$. Define
\[
 \begin{aligned}
 e(\tau,b)&=\mathbf 1_{\tau\in S_\Lambda,\ b\in J_\tau}
              -\mathbf 1_{\tau\in S_\Lambda^-,\ b\in J_\tau},\\
 c(\tau,h,b)&=\mathbf 1_{\tau\in S_\Lambda,\ h=1,\ b\notin J_\tau}
              +\mathbf 1_{\tau\in S_\Lambda^-,\ b\in J_\tau}.
 \end{aligned}
\]
Then for every target tableau $I'$, the sum of $e+c$ over its boxes is at most
\[
 \sum_{w\in\Lambda}s_w+r\sum_{\tau\in S_\Lambda}c_j(\tau),
 \qquad s_w=a\,\#\{\tau\in\Phi_a:w_\tau=w\},
\]
Here $c_j(\tau)$ counts the columns of height below $n_{j+1}$ and different from $1$.
\end{lemma}
\begin{proof}
If $\tau\in S_\Lambda^-$, the contribution $e=-1$ occurs exactly in the marked rows, whose correction is $1$. The sum is zero, and both terms vanish outside the two subsets. For $\tau\in S_\Lambda$ a column of height $1$ contributes $1$ whether or not its row entry is in $J_\tau$. These columns occur only in $\Phi_a$, with $a$ columns at each place. Other low columns contribute $1$ only for row entries in $J_\tau$. Since row entries within each tableau column are distinct, one column contributes at most $|J_\tau|\le r$. Summing over places and grouping columns of height $1$ by induced place gives $\sum_w s_w$. The remaining low columns give the second term.\leanfn{See \href{https://github.com/lsernau/grh-lean/blob/72fead102f663d5c7a407964a13dc0cf9d977207/grh/src/New/UpperBound/Divisibility/ValuationAtNonSplitPrimes.lean\#L115}{Counting boxes in one column}, \href{https://github.com/lsernau/grh-lean/blob/72fead102f663d5c7a407964a13dc0cf9d977207/grh/src/New/UpperBound/Divisibility/ValuationAtNonSplitPrimes.lean\#L256}{Exponents and total corrections at the neighboring point}.}
\end{proof}

\thmsource{Lean source: \leansourcelink{src/New/UpperBound/Divisibility/ValuationAtNonSplitPrimes.lean}{292}.}
\begin{lemma}[The single-step gauge at marked double-coset points]
\label{lem:a-marked-neighbor-jet-gauge}
\emitthmsource
Let $\Lambda$ be a marking and suppose $r_p=\lceil\log_p m\rceil$. For each $w\in\Lambda$ choose a component in the $\operatorname{diag}(\varpi_w1_{i_w},1)$ double coset with $i_w\le r_p$. Use the unitary condition to take the corresponding component at the conjugate place and take the identity at other places. This gives a finite target point $g$. If the finite datum $\psi$ at step $j$ is $1$-standard, the theta kernel for $\omega(1,g)\psi$ has gauge multiplier between the identity base points at most
\[
 p^{\,\sum_{w\in\Lambda}s_w+
      r_p\sum_{\tau\in S_\Lambda}c_j(\tau)-D_{p,j}(S_\Lambda)}.
\]
Here $D_{p,j}(S_\Lambda)$ is the predicted lower bound after removing height-$1$ columns on $S_\Lambda$ and including the deletion correction and logarithmic error.
\end{lemma}
\begin{proof}
The marked place is unramified, so we may take $p$ as local uniformizer. At each $w$ write $g_w$ as an integral invertible factor, a diagonal factor, and another integral invertible factor. On a set of indices with $|J_w|=i_w$ the diagonal entries are $p$, and elsewhere they are $1$. The unitary condition gives entries $p^{-1}$ at the conjugate place. The weight embedding reads the conjugate of $w_\tau$, so place $\tau$ sees the entry $p^{-e(\tau,b)}$, with $e$ from \statementref{Lemma}{lem:a-marked-neighbor-box-count}. There are two contributions in the Weyl representation. First, for each marked finite place the infinite places inducing it and those inducing its conjugate have the same distribution of lift character weights, so the total determinant exponent is zero. The remaining gauge factor is $p^{\sum_{I'} e(\tau,b)}$, with the sum over the boxes of the target tableau.\leanfn{See \href{https://github.com/lsernau/grh-lean/blob/72fead102f663d5c7a407964a13dc0cf9d977207/grh/src/New/UpperBound/Divisibility/ValuationAtNonSplitPrimes.lean\#L179}{Integral basis changes for the double coset}, \href{https://github.com/lsernau/grh-lean/blob/72fead102f663d5c7a407964a13dc0cf9d977207/grh/src/New/UpperBound/Divisibility/ValuationAtNonSplitPrimes.lean\#L213}{Images of the diagonal entries under the weight embeddings}, \href{https://github.com/lsernau/grh-lean/blob/72fead102f663d5c7a407964a13dc0cf9d977207/grh/src/New/UpperBound/Divisibility/ValuationAtNonSplitPrimes.lean\#L240}{Cancellation of determinant weights}.}

Use $c(\tau,h,b)$ of \statementref{Lemma}{lem:a-marked-neighbor-box-count} for the corrections on the boxes. As the marked prime is inert in $K$, every allowed partner exponent is odd and the places induced by $\tau'$ and $w_\tau$ are conjugate. Thus $\tau\in S_\Lambda$ implies $\tau'\in S_\Lambda^-$ and $J_\tau=J_{\tau'}$. If $b\notin J_\tau$, columns of height $1$ on $\tau$ supply the correction, while if $b\in J_\tau$ every height on $\tau'$ does. If $h=1$ and $\tau'\in S_\Lambda$, swap the two places. This proves the partner conditions for the weighted Schur derivatives.\leanfn{See \href{https://github.com/lsernau/grh-lean/blob/72fead102f663d5c7a407964a13dc0cf9d977207/grh/src/New/UpperBound/Divisibility/ValuationAtNonSplitPrimes.lean\#L141}{Valuation corrections for partner places}.}

The integral basis changes preserve standard data. By \statementref{Lemma}{lem:a-charged-schur-jet} the $(I',I)$ entry of the derivative matrix between integral bases has absolute value at most $p^{\sum_{I'}c(\tau,h,b)-D_{p,j}(S_\Lambda)}$. The diagonal Weyl factor changes the exponent to $\sum_{I'}(e+c)-D_{p,j}(S_\Lambda)$. Counting boxes and using $|J_w|\le r_p$ gives the claimed exponent. Finally, weak approximation at finite places and kernel equivariance reduce the gauge at any point to this integral-basis matrix norm estimate, so the bound is uniform.\leanfn{See \href{https://github.com/lsernau/grh-lean/blob/72fead102f663d5c7a407964a13dc0cf9d977207/grh/src/New/UpperBound/Kernels/ThetaLiftGauge.lean\#L186}{Bounding the kernel gauge by a diagonally weighted derivative matrix}, \href{https://github.com/lsernau/grh-lean/blob/72fead102f663d5c7a407964a13dc0cf9d977207/grh/src/New/UpperBound/Divisibility/ValuationAtNonSplitPrimes.lean\#L304}{The single-step gauge at a marked neighboring point}.}
\end{proof}

\thmsource{Lean source: \leansourcelink{src/New/UpperBound/Divisibility/ValuationAtNonSplitPrimes.lean}{868}.}
\begin{lemma}[Transferring marked Hecke operators to the single-step kernel]
\label{lem:a-marked-hecke-datum-transfer}
\emitthmsource
Let $\Psi$ be weak chain data for twist $\vartheta$ with $p\nmid q_\vartheta$ and let $\Lambda$ be a marking at $p$, so $p<x_1$ is exceptional and nonsplit, $\Lambda$ is disjoint from its conjugate, and each of its places is induced by an infinite place of $\Phi_a$. Put $r_p=\lceil\log_p m\rceil$ and $E_p^{\rm tr}=r_p f(n_m+r_p)$ and let $s_w=a\#\{\tau\in\Phi_a:w_\tau=w\}$ denote the Hecke scaling exponent of $w\in\Lambda$. Put $Y_w=p^{s_w}T_{w,r_p}^{(0)}$, where $T_{w,i}^{(j)}$ is the Hecke operator on $X_j$ at $w$ defined by $\operatorname{diag}(\varpi_w1_i,1)$. Let $S_\Lambda$ be the set of infinite places inducing places in $\Lambda$ and set
\[
 H_{\Lambda,j}
   =r_p\sum_{\tau\in S_\Lambda}c_j(\tau)-D_{p,j}(S_\Lambda),
\]
Here $c_j(\tau)$ counts columns of height below $n_{j+1}$ and different from $1$. For an ordered list $L$ of distinct places in $\Lambda$, write $Y_L=\prod_{w\in L}Y_w$ with product in the specified order. Fix step $j$, write $\varphi_j(t)$ for the original finite datum at step $j$, and $C_j^t$ for the first $j$ lifts with the original data. There is a linear operator $A_{j,L}$ on kernel data satisfying the following for each $t\in\operatorname{supp}\Psi$.
\begin{enumerate}
\item On $\varphi_j(t)$, the operator $A_{j,L}$ commutes with the Weil actions of the full finite unitary group of $X_j$ and the standard integral level of $X_{j+1}$.
\item For every source form $u\in U_0$, lifting $C_j^t u$ with $A_{j,L}\varphi_j(t)$ gives precisely $C_{j+1}^t(Y_Lu)$.
\item Suppose $j<m$. Use the identity at $w\in L$ and take at each $w\in\Lambda\setminus L$ a point of the double coset corresponding to $0\le i_w\le r_p$. Assemble these and their unitary conjugate components into a finite target point $g$. After applying the Weil action of $g$ to $A_{j,L}\varphi_j(t)$, the kernel has gauge multiplier at most $p^{\beta_L}$, where
\[
  \beta_L=\sum_{w\in L}(E_p^{\rm tr}-s_w)
                +\sum_{w\in\Lambda}s_w+H_{\Lambda,j}.
\]
\end{enumerate}
\end{lemma}
\begin{proof}
Induct on the ordered list $L$. When $L$ is empty take $A_{j,\varnothing}=1$, which gives the first two properties by definition. Since the small nonsplit prime $p$ does not divide $MN$ and $p\nmid q_\vartheta$, the original data are $1$-standard. For the double-coset translations at all marked places, \statementref{Lemma}{lem:a-marked-neighbor-jet-gauge} gives gauge multiplier
\[
 p^{\sum_{w\in\Lambda}s_w+H_{\Lambda,j}} .
\]
The estimate factors each double-coset representative into left and right integral factors and a diagonal factor, then uses the weight exponents of the boxes and the valuations of the Taylor coefficients. The total cost of the two contributions is at most
$\sum_{w\in\Lambda}s_w+r_p\sum_{\tau\in S_\Lambda}c_j(\tau)$, with divisibility $D_{p,j}(S_\Lambda)$ from the Taylor coefficients.\leanfn{See \href{https://github.com/lsernau/grh-lean/blob/72fead102f663d5c7a407964a13dc0cf9d977207/grh/src/New/UpperBound/Divisibility/ValuationAtNonSplitPrimes.lean\#L304}{The single-step estimate at double-coset points}.}

Now suppose $L=(w,L')$ with $w\notin L'$ and put $\varphi'=A_{j,L'}\varphi_j(t)$. The original datum is fixed by the target maximal compact at $w$. The first induction property makes $\varphi'$ fixed by it as well. \statementref{Lemma}{lem:a-single-mark-hecke-transfer} gives coefficients $c_i$ independent of $t$, all $p$-adically integral when $j<m$, such that for any source form $z$ spherical at $w$,
\[
 \sum_{i=0}^{r_p}c_iT_{w,i}^{(j+1)}(C_{j+1}^t z)
   =p^{E_p^{\rm tr}}C_{j+1}^t(T_{w,r_p}^{(0)}z).
\]
The target Weil representation gives the Hecke operator for the same double coset on the kernel data. Write it as $\mathsf T_{w,i}$ and define
\[
 A_{j,L}
   =\left(\sum_{i=0}^{r_p}
       c_i p^{s_w-E_p^{\rm tr}}\mathsf T_{w,i}\right) A_{j,L'} .
\]
The Hecke sum commutes with the source action. On $\varphi'$ spherical at $w$ it also commutes with the target integral level, which permutes the compact cosets used in the sum. Together with the first induction property, this gives the required commutation for $A_{j,L}$.\leanfn{See \href{https://github.com/lsernau/grh-lean/blob/72fead102f663d5c7a407964a13dc0cf9d977207/grh/src/New/UpperBound/Divisibility/ValuationAtNonSplitPrimes.lean\#L770}{Single-marking transfer and integrality of coefficients}, \href{https://github.com/lsernau/grh-lean/blob/72fead102f663d5c7a407964a13dc0cf9d977207/grh/src/New/UpperBound/Divisibility/ValuationAtNonSplitPrimes.lean\#L812}{The target integral level and the Hecke sum}.}

Target translation equivariance of the theta kernel moves $\mathsf T_{w,i}$ from the data to the lift, giving $T_{w,i}^{(j+1)}$. The second induction property says that lifting $C_j^tu$ with $\varphi'$ gives $C_{j+1}^t(Y_{L'}u)$. The source form $Y_{L'}u$ is still of source level, so it is spherical at $w$. Apply transfer and multiply by $p^{s_w-E_p^{\rm tr}}$ to obtain $C_{j+1}^t(Y_wY_{L'}u)$.\leanfn{See \href{https://github.com/lsernau/grh-lean/blob/72fead102f663d5c7a407964a13dc0cf9d977207/grh/src/New/UpperBound/Divisibility/ValuationAtNonSplitPrimes.lean\#L830}{Equivariance moves the Hecke sum to the lift value}.}

Finally, take $g$ from the third property. In expanding $\mathsf T_{w,i}$ each coset representative can be absorbed into the $w$ component of $g$. Since the original $w$ component is $1$, the new point remains the identity at $L'$ and lies in the index-$i$ double coset at $w$. Apply the third induction property. Multiplication by $c_i p^{s_w-E_p^{\rm tr}}$ increases the gauge by a factor at most $p^{E_p^{\rm tr}-s_w}$. The ultrametric inequality gives
$p^{\beta_{L'}+E_p^{\rm tr}-s_w}=p^{\beta_L}$. This completes the induction.\leanfn{See \href{https://github.com/lsernau/grh-lean/blob/72fead102f663d5c7a407964a13dc0cf9d977207/grh/src/New/UpperBound/Divisibility/ValuationAtNonSplitPrimes.lean\#L844}{Combining the double-coset representatives}, \href{https://github.com/lsernau/grh-lean/blob/72fead102f663d5c7a407964a13dc0cf9d977207/grh/src/New/UpperBound/Divisibility/ValuationAtNonSplitPrimes.lean\#L888}{Induction on the list of marked places}.}
\end{proof}

\thmsource{Lean source: \leansourcelink{src/New/UpperBound/Divisibility/ValuationAtNonSplitPrimes.lean}{975}.}
\begin{lemma}[Restricted lifts of marked Hecke powers]
\label{lem:a-marked-hecke-chain-gauge}
\emitthmsource
Let $\Psi$ be weak chain data for twist $\vartheta$ with $p\nmid q_\vartheta$, and let $\Lambda$ be a marking at $p$. In the notation above take $Y_\Lambda=\prod_{w\in\Lambda}Y_w$ and write $\mathcal L_\Psi$ for the restricted lift given by $\Psi$. For every integral source form $u$, at every point we have
\[
 \operatorname{Gauge}_p\bigl(\mathcal L_\Psi(Y_\Lambda^m u)\bigr)
 \le p^{\sum_{j<m}(\#\Lambda\,E_p^{\rm tr}+H_{\Lambda,j})}.
\]
\end{lemma}
\begin{proof}
Take an ordered list of all $\Lambda$ and let $A_j$ be the operator from \statementref{Lemma}{lem:a-marked-hecke-datum-transfer}. The exponent in its third property for $g=1$ simplifies to
\[
 \sum_{w\in\Lambda}(E_p^{\rm tr}-s_w)+\sum_{w\in\Lambda}s_w+
 H_{\Lambda,j}=\#\Lambda\,E_p^{\rm tr}+H_{\Lambda,j}.
\]
The preceding and following levels lie in their standard integral levels. The first property therefore gives commutation of $A_j$ with those levels on the original data.

Fix $t\in\operatorname{supp}\Psi$ and write $\widehat C_j^t$ for the operator of the first $j$ lifts with new data $A_i\varphi_i(t)$. Induction on $j$ gives
\[
 \widehat C_j^t u=C_j^t(Y_\Lambda^j u).
\]
For $j=0$ both sides are $u$. Suppose the identity holds for $j$. Substitute the right side into the next lift and apply the second property of \statementref{Lemma}{lem:a-marked-hecke-datum-transfer} to the source form $Y_\Lambda^ju$. This gives $C_{j+1}^t(Y_\Lambda^{j+1}u)$. Take $j=m$, sum with weights $t$, and restrict to $Y_m$. Thus $\mathcal L_\Psi(Y_\Lambda^mu)$ is the restriction obtained by lifting $u$ with the new data.\leanfn{See \href{https://github.com/lsernau/grh-lean/blob/72fead102f663d5c7a407964a13dc0cf9d977207/grh/src/New/UpperBound/Divisibility/ValuationAtNonSplitPrimes.lean\#L1011}{Moving the Hecke powers to the data step by step}.}

Since $u$ is integral, its gauge at the identity base point is at most $1$. Apply the chain gauge bound of \statementref{Lemma}{lem:a-weak-chain-gauge} to obtain the sum of the single-step exponents. Since the marked prime satisfies $p\ne N$, restriction preserves the bound.\leanfn{See \href{https://github.com/lsernau/grh-lean/blob/72fead102f663d5c7a407964a13dc0cf9d977207/grh/src/New/UpperBound/Divisibility/ValuationAtNonSplitPrimes.lean\#L984}{The gauge bound for marked Hecke powers}.}
\end{proof}

\thmsource{Lean source: \leansourcelink{src/New/UpperBound/Divisibility/LocalDivisibility.lean}{372}.}
\begin{lemma}[Composite-lift divisibility on the full-marking submodule]
\label{lem:a-full-mark-valuation}
\emitthmsource
Let $\Lambda$ be a full marking at $p$, with $Y_\Lambda$, $j_p$, $M_0$, and $d_p$ as in \statementref{Lemma}{lem:a-exceptional-hecke-colength}, and retain the local gain $\gamma_p$. For every twist with level norm coprime to $p$ and each of its chain data, the composite lift on
\[
 Y_\Lambda^m M_0+p^{j_p}M_0
\]
has pointwise $p$-adic Weyl gauge at most $p^{2d_p-\gamma_p}$.
\end{lemma}
\begin{proof}
By \statementref{Lemma}{lem:a-weak-chain-replacement} the composite lift for the original chain data is the restricted lift of some weak chain data. Write $\mathcal L$ for this lift and set $S_\Lambda=\{\tau\in\Phi:w_\tau\in\Lambda\}$ and $b_\Lambda=\#(\Phi_a\cap S_\Lambda)$. At each step $j<m$, remove the columns of height $1$ on $S_\Lambda$ and include the deletion correction and logarithmic error. The predicted divisibility lower bound $D_{p,j}(S_\Lambda)$ satisfies
\[
 D_{p,j}(S_\Lambda)
 \ge \frac{\gamma_p}{m}+b_\Lambda\frac{a}{x_1}-\frac{d_p}{m}.
\]
Fullness of the marking makes the collision sum zero, and $p$ is inert in $K$, so the uniformizer has $p$-adic valuation $1$. Substituting gives the estimate.\leanfn{See \href{https://github.com/lsernau/grh-lean/blob/72fead102f663d5c7a407964a13dc0cf9d977207/grh/src/New/UpperBound/Divisibility/JetExponent.lean\#L491}{Single-step predicted divisibility}, \href{https://github.com/lsernau/grh-lean/blob/72fead102f663d5c7a407964a13dc0cf9d977207/grh/src/New/UpperBound/Divisibility/LocalDivisibility.lean\#L254}{Lower bound with a full marking}.}

\statementref{Lemma}{lem:a-marked-hecke-chain-gauge} bounds the gauge exponent of $\mathcal L(Y_\Lambda^m u)$ by the sum of $m$ single-step exponents, with step $j$ contributing
\[
 \#\Lambda\,E_p^{\rm tr}
 +r_p\sum_{\tau\in S_\Lambda}c_j(\tau)-D_{p,j}(S_\Lambda),
\]
Here $E_p^{\rm tr}=r_p f(n_m+r_p)$ is the transfer exponent and $c_j(\tau)$ the number of the relevant low columns. The low-column and transfer cost estimates give
\[
 \#\Lambda\,E_p^{\rm tr}+r_p\sum_{\tau\in S_\Lambda}c_j(\tau)
     \le \frac{d_p}{m}+b_\Lambda\frac a{x_1}.
\]
Subtracting the two estimates cancels $b_\Lambda a/x_1$, giving a single-step exponent at most $(2d_p-\gamma_p)/m$. Summing gives, for $u\in M_0$,\leanfn{See \href{https://github.com/lsernau/grh-lean/blob/72fead102f663d5c7a407964a13dc0cf9d977207/grh/src/New/UpperBound/Divisibility/ValuationAtNonSplitPrimes.lean\#L969}{Hecke transfer and the gauge at each step}, \href{https://github.com/lsernau/grh-lean/blob/72fead102f663d5c7a407964a13dc0cf9d977207/grh/src/New/UpperBound/Divisibility/MarkedPrimes.lean\#L513}{Cost of the low columns and transfer}.}
\[
 \operatorname{Gauge}_p\bigl(\mathcal L(Y_\Lambda^m u)\bigr)
 \le p^{2d_p-\gamma_p}.
\]

Now take $v\in M_0$. By \statementref{Lemma}{lem:a-weak-chain-unmarked-valuation}, with no marking the collision cost for each initial infinite place is at most $\sum_{j\ge1}a/p^j=a/(p-1)\le 2a/p$, hence at most $2mfa/p\le j_p$ over all $m$ steps. The unmarked gauge bound gives exponent at most
\[
 m C_p(\varnothing)+d_p-\gamma_p.
\]
Multiplication by $p^{j_p}$ reduces the exponent by $j_p$. Thus
\[
 \operatorname{Gauge}_p\bigl(\mathcal L(p^{j_p}v)\bigr)\le p^{d_p-\gamma_p}.
\]
By $d_p\ge0$ this bound is at most $p^{2d_p-\gamma_p}$ as well. The ultrametric triangle inequality for the local Weyl gauge gives the same bound for $\mathcal L(Y_\Lambda^m u+p^{j_p}v)$. Every vector of the submodule admits this decomposition.\leanfn{See \href{https://github.com/lsernau/grh-lean/blob/72fead102f663d5c7a407964a13dc0cf9d977207/grh/src/New/UpperBound/Divisibility/LocalDivisibility.lean\#L183}{Unmarked collision cost and the power correction}, \href{https://github.com/lsernau/grh-lean/blob/72fead102f663d5c7a407964a13dc0cf9d977207/grh/src/New/UpperBound/Divisibility/LocalDivisibility.lean\#L371}{Combining the two summands after changing to weak chain data}.}
\end{proof}

\thmsource{Lean source: \leansourcelink{src/New/UpperBound/Divisibility/JetExponent.lean}{163}.}
\begin{lemma}[Collision losses at nonexceptional primes]
\label{lem:a-nonexceptional-collision}
\emitthmsource
Let $p<x_1$ be a small nonsplit prime that is nonexceptional. Then
\[
 m\,C_p(\varnothing)\le d_p,\qquad
 d_p=\frac{fa\ell^2}{p}
 \begin{cases}
  m,&p\equiv\pm1\pmod m,\\
  1,&p\not\equiv\pm1\pmod m.
 \end{cases}
\]
\end{lemma}
\begin{proof}
Collision terms are indexed by an infinite place $\tau$ and a positive integer $j$ and have size $a/p^j$. There is no collision when $p=d$. Such a collision would require $\tau$ and $-d^j\tau$ to both lie in the class $1$ modulo $q_0$. But $d\equiv q_0/2-1\pmod{q_0}$ and $4\mid q_0$ give $d^2\equiv1\pmod{q_0}$, so $d^j$ modulo $q_0$ can only be $1$ or $q_0/2-1$. Neither equals $-1$, ruling out the required $-d^j\equiv1\pmod{q_0}$.

Suppose $p\ne d$ and a collision occurs. The partner exponent condition gives $\chi_K(p)=-1$ with $j$ odd, hence $-p^j=(-p)^j$. Membership of both places in $\Phi_a$ forces $o_0\mid j$, where $o_0$ is the order of $-p$ modulo $q_0$. If all collision exponents are at least $J$, sum over $\tau$ and enlarge the range of the geometric series to obtain\leanfn{See \href{https://github.com/lsernau/grh-lean/blob/72fead102f663d5c7a407964a13dc0cf9d977207/grh/src/New/Derivatives/Sharpness/NearPairs.lean\#L53}{Excluding collisions at the ramified prime}, \href{https://github.com/lsernau/grh-lean/blob/72fead102f663d5c7a407964a13dc0cf9d977207/grh/src/New/UpperBound/Divisibility/JetExponent.lean\#L124}{Collision exponents and the geometric series}.}
\[
 C_p(\varnothing)
 \le f\sum_{j\ge J}\frac{a}{p^j}
 =\frac{fa}{p^J}\frac{p}{p-1}
 \le\frac{2fa}{p^J}.
\]

If $p\equiv\pm1\pmod m$, take $J=1$. Then $d_p=mfa\ell^2/p$, and $\ell^2\ge2$ gives the bound. If $p\not\equiv\pm1\pmod m$, any collision forces $o_0$ odd. Since $p$ is nonexceptional, $o_0>r_p$ with $r_p=\lceil\log_p m\rceil$. Thus each collision exponent is at least $r_p+1$. Take $J=r_p+1$ and use $p^{r_p}\ge m$ to get
\[
 m\,C_p(\varnothing)
 \le \frac{2mfa}{p^{r_p+1}}
 \le \frac{2fa}{p}
 \le \frac{fa\ell^2}{p}=d_p.
\]
If there are no collisions, the left side is zero and the bound holds.\leanfn{See \href{https://github.com/lsernau/grh-lean/blob/72fead102f663d5c7a407964a13dc0cf9d977207/grh/src/New/UpperBound/Divisibility/JetExponent.lean\#L162}{Total loss in the nonexceptional case}.}
\end{proof}

\section{The level prime and uniform local losses}
\label{sec:a-level-prime}

At the level prime $N$, gluing changes the lattice in which integrality is measured. We bound the cost of correcting source forms into the required lattice, then combine it with the losses at the other primes.

Fix a $N$-adic place $\iota$ and write $\mathcal O_\iota$ for its valuation ring, with $|N|_\iota=N^{-1}$. Let $\Sigma_N$ be the conjugate set of places induced by $\iota$. It contains one place from each conjugate pair above $N$, each with residue field $\mathbb F_N$. Write $\pi_N\bar\pi_N=N$. For upper-chain space $X_j$ at stage $j$ let $\delta_j(S)$ be the finite adelic unitary point given by $\operatorname{diag}(1,\ldots,1,N)$ at $\Sigma_N\cap S$ and the identity at $\Sigma_N\setminus S$, and write
\[
 \delta_j^{\mathrm{all}}=\delta_j(\Sigma_N),\qquad
 \delta_j=\delta_j(\{w:w\mid\pi_N\}),\qquad
 \delta_j^\vee=\delta_j(\delta_j^{\mathrm{all}})^{-1}.
\]
At $\Sigma_N$, the points $\delta_j$ and $\delta_j^\vee$ send the standard lattice to $L_j$ and its dual, respectively. A stage-$j$ form $\phi$ is integral at $d$ if, for every finite adelic point $g$, the Weyl action of $(gd)^{-1}$ moving $\phi(g)$ back to the identity gives all coordinates in $\mathcal O_\iota$. Write $\mathcal A_\iota,\mathcal A^L,\mathcal A^\vee$ for the full source form modules integral at $1,\delta_0,\delta_0^\vee$, respectively.\leanfn{See \href{https://github.com/lsernau/grh-lean/blob/72fead102f663d5c7a407964a13dc0cf9d977207/grh/src/New/UpperBound/Divisibility/SetupAtLevel.lean\#L32}{Lattices and adelic points at the level prime}.}

\thmsource{Lean source: \leansourcelink{src/New/UpperBound/Divisibility/ColengthAtLevel.lean}{98}.}
\begin{lemma}[Mean degree in the last coordinate of the source space]
\label{lem:a-level-mean-degree}
\emitthmsource
Let $\mathcal J$ be the product of source-level double cosets for $X_0$ with the Weyl basis index set. For $i\in\mathcal J$ and a source coordinate variable $v$, let $d_i(v)$ be the total degree of letter $v$ in the tableaux of the Weyl basis vector at all infinite places. Put $s(v)=\sum_{i\in\mathcal J}d_i(v)$. Then
\[
 k\,s(v)\le u_0fa,\qquad u_0=\dim_{\mathbb C}U_0.
\]
\end{lemma}
\begin{proof}
Fix two coordinate variables $v,v'$. Swapping them at all infinite places conjugates the diagonal matrix with diagonal coefficient $2$ only at $v$ into the one with $2$ only at $v'$. Their Weyl representations thus have equal determinants. A diagonal matrix acts on the tableau basis by the degrees of the letters, with eigenvalues $2^{d_i(v)}$ on the module for each double coset. The source level has trivial stabilizers, so the values at each double coset lie in the same Weyl module. Multiplying determinants over the cosets gives $2^{s(v)}=2^{s(v')}$, and $s(v)$ is independent of $v$.

The degree of the tableau at one infinite place is the number of boxes in the columns of height below $k$. The source weight has at most $a$ such boxes per infinite place, so $\sum_v d_i(v)\le fa$. Summing over $i$ and using $|\mathcal J|=u_0$ gives\leanfn{See \href{https://github.com/lsernau/grh-lean/blob/72fead102f663d5c7a407964a13dc0cf9d977207/grh/src/New/UpperBound/Divisibility/ColengthAtLevel.lean\#L75}{Tableau degrees, transpositions, and mean degrees}.}
\[
 k\,s(v)=\sum_{i\in\mathcal J}\sum_v d_i(v)\le u_0fa .
\]
\end{proof}

\thmsource{Lean source: \leansourcelink{src/New/UpperBound/Divisibility/ColengthAtLevel.lean}{156}.}
\begin{lemma}[Colength between the dual lattice and the $L_0$-lattice]
\label{lem:a-level-diagonal-colength}
\emitthmsource
Let $s_N=s(v_{\mathrm{last}})$ be the total degree of the last coordinate in the source space. Any finite family in $\mathcal A^\vee$ has a correction by a matrix over $\mathcal O_\iota$ of determinant dividing $N^{s_N}$ that sends the family into $\mathcal A^L$.
\end{lemma}
\begin{proof}
Given a representative $c$ of a source-level double coset, use weak approximation at finite places for $c\delta_0^\vee$, also requiring its conjugate under $(\delta_0^{\mathrm{all}})^{-1}$ to be integral above $N$. Replace $c$ within the double coset by $r_c$, with
\[
 r_c\delta_0^\vee\text{ integral and invertible above }N\text{ },\qquad
 (\delta_0^{\mathrm{all}})^{-1}r_c\delta_0\text{ integral and invertible above }N\text{ },
\]
The second equality uses $\delta_0^\vee \delta_0^{\mathrm{all}}=\delta_0$. The Weyl gauge is invariant under rational left translation and under integral invertible right multiplication above $N$. Thus for every $\phi\in\mathcal A^\vee$ the coordinates at the representatives $r_c$ belong to $\mathcal O_\iota$.\leanfn{See \href{https://github.com/lsernau/grh-lean/blob/72fead102f663d5c7a407964a13dc0cf9d977207/grh/src/Lib/Adelic/AutomorphicForms/PAdicGauge.lean\#L84}{Weak approximation at finite places}, \href{https://github.com/lsernau/grh-lean/blob/72fead102f663d5c7a407964a13dc0cf9d977207/grh/src/New/UpperBound/Divisibility/ColengthAtLevel.lean\#L145}{Representatives adapted to both lattices}.}

Now consider the gauge at $\delta_0$. If $k_0$ lies in the source level, then $\delta_0^{-1}k_0\delta_0$ is integral invertible above $N$. At $\Sigma_N$ the matrix for the source level in the $L_0$ basis is integral invertible and its last column reduces to the last basis vector modulo $N$. At the conjugate place of $\Sigma_N$ this follows from the unitary condition. Write any point as $\gamma r_ck_0$ with $\gamma$ rational. The factorization
\[
 r_ck_0\delta_0
   =\delta_0^{\mathrm{all}}\bigl((\delta_0^{\mathrm{all}})^{-1}r_c\delta_0\bigr)
                  \bigl(\delta_0^{-1}k_0\delta_0\bigr)
\]
and the two gauge invariance properties reduce the check of $\phi(r_c)$ to $\delta_0^{\mathrm{all}}$. Each weight embedding reads one place of $\Sigma_N$, where $\delta_0^{\mathrm{all}}$ multiplies the last coordinate by $N$. The determinant power of the source weight is zero, so the Weyl action of $(\delta_0^{\mathrm{all}})^{-1}$ multiplies coordinate $i=(c,\alpha)$ by $N^{-d_i(v_{\mathrm{last}})}$. If the value coordinates at every $i$ are divisible by $N^{d_i(v_{\mathrm{last}})}$, then $\phi\in\mathcal A^L$.\leanfn{See \href{https://github.com/lsernau/grh-lean/blob/72fead102f663d5c7a407964a13dc0cf9d977207/grh/src/New/UpperBound/Divisibility/SetupAtLevel.lean\#L191}{The source level in the $L_0$ basis}, \href{https://github.com/lsernau/grh-lean/blob/72fead102f663d5c7a407964a13dc0cf9d977207/grh/src/New/UpperBound/Divisibility/ColengthAtLevel.lean\#L182}{The corrected values satisfy Weyl integrality}.}

Correct the value coordinates of the given family in order of the indices $i$. If the current coordinate does not vanish for the whole family, pick a vector with smallest valuation at that coordinate. Its entry divides the other entries, so integral elementary row operations clear them. Multiply the chosen vector by $N^{d_i(v_{\mathrm{last}})}$. If all entries are zero, do nothing. Each operation preserves divisibility of the earlier coordinates, and the resulting matrix has determinant dividing
\[
 \prod_{i\in\mathcal J} N^{d_i(v_{\mathrm{last}})}=N^{s_N}.
\]
The gauge criterion puts the corrected family in $\mathcal A^L$.\leanfn{See \href{https://github.com/lsernau/grh-lean/blob/72fead102f663d5c7a407964a13dc0cf9d977207/grh/src/New/UpperBound/Divisibility/ColengthAtLevel.lean\#L159}{Colength for the diagonal coordinate lattice}, \href{https://github.com/lsernau/grh-lean/blob/72fead102f663d5c7a407964a13dc0cf9d977207/grh/src/New/Reduction/Lattices/Colength.lean\#L134}{Coordinate correction over a valuation ring}.}
\end{proof}

\thmsource{Lean source: \leansourcelink{src/New/UpperBound/Divisibility/ColengthAtLevel.lean}{243}.}
\begin{lemma}[Colength for iterates of the raising operator]
\label{lem:a-level-up-iterate}
\emitthmsource
Let $\mathsf U_N$ be the raising operator at the level, read in the $L_0$ basis. For $i\ge0$ put $\mathcal A_i=(\mathsf U_N^i)^{-1}(\mathcal A^L)$. Any finite family in $\mathcal A_\iota$ can be sent into $\mathcal A_i$ by an integral matrix of determinant dividing $N^{s_N(i+1)}$.
\end{lemma}
\begin{proof}
First $\mathcal A_\iota\subseteq\mathcal A^\vee$. Indeed $(\delta_0^\vee)^{-1}=\delta_0^{-1}\delta_0^{\mathrm{all}}$ is an integral diagonal matrix under all the weight embeddings used for the gauge. The Weyl representation of the source weight is polynomial, so the matrix preserves coordinate integrality.

Next prove $\mathsf U_N(\mathcal A^L)\subseteq\mathcal A^\vee$. The raising operator preserves source-level invariance. Above $N$ this can be checked locally: above $N$ the source level in the $L_0$ basis fixes the last column modulo $N$, and the coset sum preserves its invariants. At the remaining finite places the source-level action commutes with every summand. At the finite coordinate $y$ the translations in the raising operator have the form
\[
 g_y=\delta_0^\vee k_y\delta_0^{-1},
\]
where $k_y$ is integral invertible above $N$. Their coefficients are values of finite additive characters, of $N$-adic absolute value $1$. At each finite adelic point $g$,
$gg_y\delta_0=g\delta_0^\vee k_y$. Thus when $\phi$ is integral at $\delta_0$ each summand of the raising operator has gauge at most $1$ at $g\delta_0^\vee$. The ultrametric inequality gives the required $\delta_0^\vee$ integrality.\leanfn{See \href{https://github.com/lsernau/grh-lean/blob/72fead102f663d5c7a407964a13dc0cf9d977207/grh/src/New/UpperBound/Divisibility/ColengthAtLevel.lean\#L118}{The standard and dual integral modules}, \href{https://github.com/lsernau/grh-lean/blob/72fead102f663d5c7a407964a13dc0cf9d977207/grh/src/New/UpperBound/Divisibility/SetupAtLevel.lean\#L287}{The raising operator preserves the source level}, \href{https://github.com/lsernau/grh-lean/blob/72fead102f663d5c7a407964a13dc0cf9d977207/grh/src/New/UpperBound/Divisibility/ColengthAtLevel.lean\#L217}{The raising operator from the $L_0$ integral module to the dual integral module}.}

For $i=0$, $\mathcal A_\iota\subseteq\mathcal A^\vee$ and \statementref{Lemma}{lem:a-level-diagonal-colength} give exponent $s_N$. Suppose a family was corrected into $\mathcal A_i$. Apply $\mathsf U_N^{i+1}$ to get an image family in $\mathcal A^\vee$ by $\mathsf U_N(\mathcal A^L)\subseteq\mathcal A^\vee$, and use \statementref{Lemma}{lem:a-level-diagonal-colength} to send it into $\mathcal A^L$. Since the operator commutes with the coefficient matrix, the second matrix sends the initial family in $\mathcal A_i$ into $\mathcal A_{i+1}$. The determinants multiply, increasing the exponent from $s_N(i+1)$ to $s_N(i+2)$, which completes the induction.\leanfn{See \href{https://github.com/lsernau/grh-lean/blob/72fead102f663d5c7a407964a13dc0cf9d977207/grh/src/New/UpperBound/Divisibility/ColengthAtLevel.lean\#L243}{Induction for inverse images of the raising operator}.}
\end{proof}

\thmsource{Lean source: \leansourcelink{src/New/UpperBound/Divisibility/ColengthAtLevel.lean}{265}.}
\begin{lemma}[Colength of the full integral source module at the level prime]
\label{lem:a-level-up-colength}
\emitthmsource
Let $\mathcal A'=(\mathsf U_N^{m/2})^{-1}(\mathcal A^L)$. There exists $c_N\in\mathbb N$ such that any finite family in $\mathcal A_\iota$ can be sent into $\mathcal A'$ by an integral matrix with determinant dividing $N^{c_N}$, and
\[
 c_N\le m f b\,u_0.
\]
\end{lemma}
\begin{proof}
Take $i=m/2$ in \statementref{Lemma}{lem:a-level-up-iterate} and put $c_N=s_N(m/2+1)$. \statementref{Lemma}{lem:a-level-mean-degree} gives $ks_N\le u_0fa$ and the scales give $(m/2+1)a\le mbk$. Thus
\[
 kc_N \le u_0fa(m/2+1)\le k\,mfb\,u_0.
\]
Since the source rank is $k>0$, cancel $k$ to get the bound.\leanfn{See \href{https://github.com/lsernau/grh-lean/blob/72fead102f663d5c7a407964a13dc0cf9d977207/grh/src/New/UpperBound/Divisibility/SetupAtLevel.lean\#L321}{The integral module given by the raising-operator preimage}, \href{https://github.com/lsernau/grh-lean/blob/72fead102f663d5c7a407964a13dc0cf9d977207/grh/src/New/UpperBound/Divisibility/ColengthAtLevel.lean\#L265}{The final source-module exponent}.}
\end{proof}

\thmsource{Lean source: \leansourcelink{src/New/UpperBound/Divisibility/ValuationAtLevel.lean}{331}.}
\begin{lemma}[Determinant character factors in the gluing basis change]
\label{lem:a-level-glued-factor}
\emitthmsource
Fix a tame character $\vartheta$. Write $\chi_j$ for the base change to $E$ of the character for step $j$ of the upper chain. For a matrix $A$ at $\Sigma_N$, put
\[
 \chi_{j,r}(A)=\chi_j(\det A)|\det A|^{-r/2},
\]
where $|\det A|$ is the finite-place modulus. Define the gluing factor $\Delta_{\vartheta,i}$ of double step $i$ as the product of the value of $\chi_{2i+2,n_{2i+1}}$ on matrix $\delta_{2i+2}(\Sigma_N\setminus\{w:w\mid\pi_N\})$ at $X_{2i+2}$ and the value of $\chi_{2i+1,n_{2i+1}}$ on $\delta_{2i}(\{w:w\mid\pi_N\})$ at $X_{2i}$. Then for $0\le i<m/2$,
\[
 |\Delta_{\vartheta,i}|_\iota\le N^{2fb}.
\]
\end{lemma}
\begin{proof}
For any $\delta_j(S)$, its determinant at $\Sigma_N\cap S$ is $N$ and elsewhere it is $1$. Thus $|\det|^{-r/2}$ is a nonnegative power of $\sqrt N$ and has $N$-adic absolute value at most $1$. For the base-change character itself, unitarity means that at the $N$-adic place corresponding to $\tau$ the determinant read is the inverse determinant at the conjugate place. Hence the $\tau$-th weight factor can exceed $1$ only when $\bar w_\tau\in S$, when it is a power of $N$ given by the weight. The two edge indices are at most $m$ and the weight at any infinite place is at most $2b$, so the factor is at most $N^{2b}$.

The two place sets of $\Sigma_N$ occurring in the gluing factor are complementary. At each $\tau$ exactly one term can supply the factor exceeding $1$. Multiplying over the $f$ infinite places gives $N^{2fb}$.\leanfn{See \href{https://github.com/lsernau/grh-lean/blob/72fead102f663d5c7a407964a13dc0cf9d977207/grh/src/New/UpperBound/Divisibility/ValuationAtLevel.lean\#L53}{Placewise absolute values of determinant characters}, \href{https://github.com/lsernau/grh-lean/blob/72fead102f663d5c7a407964a13dc0cf9d977207/grh/src/New/UpperBound/Divisibility/ValuationAtLevel.lean\#L323}{Complementary place sets and the double-step gluing factor}.}
\end{proof}

\thmsource{Lean source: \leansourcelink{src/New/UpperBound/Divisibility/ValuationAtLevel.lean}{403}.}
\begin{lemma}[Reducing the glued chain to direct lifts]
\label{lem:a-level-direct-chain}
\emitthmsource
Let $\vartheta$ be tame. For the finite Weil data $v_j$ on the upper edges, write $\mathcal G_i$ for the lift obtained by projecting $v_{2i},v_{2i+1}$ onto the vacuum above $N$ and gluing double steps. Write $\mathcal D_j$ for the lift with direct data in the $L_j,L_{j+1}$ basis, the indicator on integral matrices whose last column is the final basis vector modulo $N$. There is a constant $c_\vartheta\in\mathbb C$ depending only on $D,\vartheta$ and the chosen $N$-adic place, such that for every ramified source form $\phi$ and all finite data $v_j$,
\[
 \mathcal G_{m/2-1}\cdots\mathcal G_0\phi
 =c_\vartheta\,\mathcal D_{m-1}\cdots\mathcal D_0
                (\mathsf U_N^{m/2}\phi),\qquad
 |c_\vartheta|_\iota\le N^{mfb}.
\]
\end{lemma}
\begin{proof}
In the $L_j$ basis, write $\rho_j$ for translation of stage-$j$ forms by a matrix at $\Sigma_N$ and $\mathsf U_j$ for the raising operator, so $\mathsf U_0=\mathsf U_N$. Write $\mathcal L_j(F)$ for the lift with a matrix function $F$ in the two lattice bases. The substitution $T\mapsto B^{-1}TA$ in that function moves $A^{-1}$ to the input and $B$ to the output and introduces the determinant character ratio
\[
 \frac{\chi_{j+1,n_{j+1}}(A)}{\chi_{j+1,n_j}(B)} .
\]
This is the formula for the Weil action on the matrix argument together with lift equivariance. For integral invertible matrices the ratio is $1$ because the lift character is unramified above $N$. If $D_j^\sharp$ denotes scaling the last coordinate by $N^{-1}$ in either lattice basis, the last-row indicator $\mathbf 1_{\mathrm{row}}$ and the direct data satisfy
\[
 \rho_{j+1}(D_{j+1}^\sharp)\,\mathcal L_j(\mathbf 1_{\mathrm{row}})
   =\eta_j\,\mathcal D_j\,\rho_j(D_j^\sharp),
\]
Here $\eta_j$ is the quotient of $\chi_{j+1,n_{j+1}}(\delta_j^{\mathrm{all}})$ at $X_j$ by $\chi_{j+1,n_j}(\delta_{j+1}^{\mathrm{all}})$ at $X_{j+1}$. The matrices have the same local determinants so their character values cancel. By $n_{j+1}\ge n_j$, the remaining modulus power has $N$-adic absolute value at most $1$. Thus $|\eta_j|_\iota\le1$.\leanfn{See \href{https://github.com/lsernau/grh-lean/blob/72fead102f663d5c7a407964a13dc0cf9d977207/grh/src/New/UpperBound/Divisibility/ValuationAtLevel.lean\#L231}{Local basis changes and the lifting formula}, \href{https://github.com/lsernau/grh-lean/blob/72fead102f663d5c7a407964a13dc0cf9d977207/grh/src/New/UpperBound/Objects/Columns.lean\#L331}{Absolute values of determinant character ratios}.}

For each $w\in\Sigma_N$, a ramified vector can be written as $\psi-\rho_j(x)\psi$, where $x$ is a local integral unit scalar congruent to $1$ at the other places modulo $N$. For source forms this is the ramification condition at $w$ or its conjugate. At intermediate stages, the same scalar on the input and output leaves the matrix data unchanged, so the direct lift preserves this property. The raising operator consists of local translations commuting with the scalars and also preserves ramification.

At double step $i$ the local gluing expansion runs over vectors $y$ modulo $N$. On the first edge it uses the indicator with last column $y$. On the second it uses the finite additive character pairing the last row with $y$. When the endpoint bases are changed to $L_{2i},L_{2i+2}$, the opposite basis changes at the intermediate level cancel. This leaves the factor $\Delta_{\vartheta,i}$ and the output action $\rho_{2i+2}(D_{2i+2}^\sharp)$. If $y$ is zero at some $w$, integral unit scalars there fix the first-edge datum, and the term vanishes on a ramified vector. If $y$ is nonzero at every place, take an integral invertible basis change taking the last basis vector to $y$ over each residue field. The change cancels between the two edges, so those terms are all equal. Write $h_i$ for the number of vectors nonzero at every place. The double step is then $\Delta_{\vartheta,i}h_i$ times the direct lift followed by the reflected lift and the output translation.\leanfn{See \href{https://github.com/lsernau/grh-lean/blob/72fead102f663d5c7a407964a13dc0cf9d977207/grh/src/New/UpperBound/Divisibility/SetupAtLevel.lean\#L245}{Ramified forms in the local representation}, \href{https://github.com/lsernau/grh-lean/blob/72fead102f663d5c7a407964a13dc0cf9d977207/grh/src/New/UpperBound/Divisibility/ValuationAtLevel.lean\#L349}{Finite expansion of glued data in the lattice bases}, \href{https://github.com/lsernau/grh-lean/blob/72fead102f663d5c7a407964a13dc0cf9d977207/grh/src/New/UpperBound/Objects/LiftLawsModuloN.lean\#L770}{Contribution of vectors nonzero at every place}.}

The reflected datum is the finite additive character of the bottom right entry of an integral matrix. Its part with bottom right entry zero at some place also vanishes on ramified vectors. Grouping the remaining terms by their last column shows that the reflected lift with output translation is $\mathsf U_{j+1}\mathcal D_j$. Grouping by the last row and using its basis change identity gives $\eta_j\mathcal D_j\mathsf U_j$ for the same lift. Thus moving the raising operator past the $j$-th direct lift introduces only the factor $\eta_j$. The $i$-th glued double step on ramified vectors is therefore
\[
 \Delta_{\vartheta,i}h_i\,\eta_{2i+1}\eta_{2i}\,
    \mathcal D_{2i+1}\mathcal D_{2i}\mathsf U_{2i}.
\]
Move each $\mathsf U_{2i}$ past the $2i$ preceding direct lifts in turn. Induction gives the identity with common scalar
\[
 c_\vartheta
 =\prod_{i=0}^{m/2-1}\left(
   \Delta_{\vartheta,i}h_i\prod_{j=0}^{2i+1}\eta_j\right).
\]
Since $h_i$ is an integer and $|\eta_j|_\iota\le1$, only $\Delta_{\vartheta,i}$ can increase the gauge. \statementref{Lemma}{lem:a-level-glued-factor} bounds the total increase by $(N^{2fb})^{m/2}=N^{mfb}$.\leanfn{See \href{https://github.com/lsernau/grh-lean/blob/72fead102f663d5c7a407964a13dc0cf9d977207/grh/src/New/UpperBound/Objects/LiftLawsModuloN.lean\#L687}{Commutation of the reflected data and the raising operator}, \href{https://github.com/lsernau/grh-lean/blob/72fead102f663d5c7a407964a13dc0cf9d977207/grh/src/New/UpperBound/Objects/LiftLawsModuloN.lean\#L885}{Induction moving the raising operators toward the source}, \href{https://github.com/lsernau/grh-lean/blob/72fead102f663d5c7a407964a13dc0cf9d977207/grh/src/New/UpperBound/Divisibility/ValuationAtLevel.lean\#L403}{The common factor for the direct chain}.}
\end{proof}

\thmsource{Lean source: \leansourcelink{src/New/UpperBound/Divisibility/ValuationAtLevel.lean}{488}.}
\begin{lemma}[Single-step gauge for direct data]
\label{lem:a-level-direct-gauge}
\emitthmsource
For $j<m$ apply the direct datum to any integral vacuum translate on upper edge $j$ and read it in the $L_j,L_{j+1}$ bases. For any tame character this kernel sends a source value of gauge at most $C$ at $g\delta_j$ to a target value of gauge at most $N^{fb}C$ at $g'\delta_{j+1}$, where $g,g'$ are arbitrary finite adelic points at the two stages.
\end{lemma}
\begin{proof}
First use the standard bases. Choose an integer $\omega_K=(1+\sqrt{-d})/2$ of $K$ reducing to $\bar u\in\mathbb F_N$ modulo $N$ at the $K$-place induced by $\iota$. Take an integer representative $u$ and put $\alpha=\omega_K-u$. The $E$-places induced by $\iota$ restrict to the same place of $K$, so $\alpha$ vanishes at them all. Meanwhile $\bar\alpha$ reduces to $1-2\bar u$ there, whose square is $-d\ne0$. Thus $\bar\alpha$ vanishes on $\Sigma_N$ and $|\bar\alpha d|_\iota=1$.\leanfn{See \href{https://github.com/lsernau/grh-lean/blob/72fead102f663d5c7a407964a13dc0cf9d977207/grh/src/New/UpperBound/Divisibility/ValuationAtLevel.lean\#L96}{An integer clearing the Fourier translation denominators}.}

The direct datum is a function on matrices modulo $N$ with values in $\{0,1\}$. Expand it by finite Fourier inversion on $\Sigma_N$ (\emph{finite Fourier expansion}), bounding the Fourier denominator for $n_{j+1}\times n_j$ matrix coordinates by
\[
 F_j=\operatorname{Nm}_{E/\mathbb Q}((N))^{n_{j+1}n_j}
     =N^{2f n_{j+1}n_j}.
\]
The actual Fourier denominator is the cardinality of the residue matrix space. Since the modulus ideal contains $(N)$, the cardinality divides $F_j$. Multiplication by $F_j$ makes the datum an integral combination of a vacuum projection above $N$ followed by finitely many Heisenberg translations. Since $\bar\alpha$ clears the Fourier denominators on $\Sigma_N$, the local denominators of the translations can be cleared by $\alpha$. The vacuum projection itself sends a translated vacuum to zero or to a translate integral above $N$. Thus, after multiplication by $F_j$, the data belong to the standard-data module with translation denominators cleared by $\alpha$. Since the extra kernel scaling is a $N$-adic integer and $|\bar\alpha d|_\iota=1$, the standard kernel estimate gives a factor at most $1$.\leanfn{See \href{https://github.com/lsernau/grh-lean/blob/72fead102f663d5c7a407964a13dc0cf9d977207/grh/src/New/UpperBound/Divisibility/ValuationAtLevel.lean\#L126}{Vacuum projection and Fourier expansion of the direct datum}, \href{https://github.com/lsernau/grh-lean/blob/72fead102f663d5c7a407964a13dc0cf9d977207/grh/src/New/UpperBound/Kernels/ThetaLiftGauge.lean\#L288}{Kernel gauge for standard data}.}

Dividing by $F_j$ gives a factor at most $N^{2fn_{j+1}n_j}$ in the standard bases. Together with $n_j,n_{j+1}\le n_m$ and the admissible scales $2n_m^2\le b$, this is at most $N^{fb}$. Finally, apply Weil translation by $(\delta_j,\delta_{j+1})$ to the data. Kernel equivariance carries the standard-basis gauge estimate to $g\delta_j$ and $g'\delta_{j+1}$ with the same constant.\leanfn{See \href{https://github.com/lsernau/grh-lean/blob/72fead102f663d5c7a407964a13dc0cf9d977207/grh/src/New/UpperBound/Divisibility/ValuationAtLevel.lean\#L197}{Scale bound for the single-step loss}, \href{https://github.com/lsernau/grh-lean/blob/72fead102f663d5c7a407964a13dc0cf9d977207/grh/src/New/UpperBound/Divisibility/ValuationAtLevel.lean\#L488}{Changing to the $L_j$ lattice bases}.}
\end{proof}

\thmsource{Lean source: \leansourcelink{src/New/UpperBound/Divisibility/ValuationAtLevel.lean}{504}.}
\begin{lemma}[The gauge of the composite lift at the level prime]
\label{lem:a-level-composite-gauge}
\emitthmsource
If a ramified source form $\phi\in U$ is in $\mathcal A'$, then for every twist $\vartheta$ and its chain data $\Psi$ the composite lift $\operatorname{Comp}_{\vartheta,\Psi}(\phi)$ has $N$-adic Weyl gauge at most $N^{2mfb}$ at every point for the standard lattice.
\end{lemma}
\begin{proof}
Set $\psi=\mathsf U_N^{m/2}\phi$. By definition of $\mathcal A'$, the source form $\psi$ is integral at $\delta_0$. The denominators of the chain data come only from $M$ and the twist modulus, both coprime to $N$. Each vacuum translate is integral above $N$ and hence fixed by the vacuum projection at $N$. Apply \statementref{Lemma}{lem:a-level-direct-chain} to each translation and sum with the coefficients of $\Psi$. This gives the sum of the direct lifts of $\psi$ multiplied by a common constant $c_\vartheta$ of gauge at most $N^{mfb}$.\leanfn{See \href{https://github.com/lsernau/grh-lean/blob/72fead102f663d5c7a407964a13dc0cf9d977207/grh/src/New/UpperBound/Divisibility/ChainDatumExpansion.lean\#L52}{Projection of admissible vacuum translates at $N$}, \href{https://github.com/lsernau/grh-lean/blob/72fead102f663d5c7a407964a13dc0cf9d977207/grh/src/New/UpperBound/Divisibility/ValuationAtLevel.lean\#L516}{Changing from the glued chain to the direct chain}.}

In the $L_j$ basis at $\Sigma_N$ the upper-chain levels are integral invertible matrices whose last column reduces to the last basis vector modulo $N$. Acting on the matrix argument of the direct datum by levels $j$ and $j+1$ preserves integrality and this last-column condition. The direct-data operator therefore commutes with the Weil actions of those levels. By \statementref{Lemma}{lem:a-weak-chain-vacuum-projection} the original unglued tensor is jointly invariant, so applying the direct-data operators at each step preserves joint invariance. At each point the sum of direct lifts then expands as a finite class sum. Each term is an integral coefficient $\Psi(x)$ times $m$ kernels applied successively to the source value. By \statementref{Lemma}{lem:a-level-direct-gauge} each kernel contributes at most $N^{fb}$, so the gauge of a term at $\delta_m$ is at most $N^{mfb}$. The ultrametric inequality gives this for the class sum. Multiplication by $c_\vartheta$ bounds the gauge of the upper-chain result at $\delta_m$ by $N^{2mfb}$.\leanfn{See \href{https://github.com/lsernau/grh-lean/blob/72fead102f663d5c7a407964a13dc0cf9d977207/grh/src/New/UpperBound/Divisibility/ValuationAtLevel.lean\#L449}{Direct data commute with the level actions}, \href{https://github.com/lsernau/grh-lean/blob/72fead102f663d5c7a407964a13dc0cf9d977207/grh/src/New/UpperBound/Divisibility/ChainDatumExpansion.lean\#L415}{Joint invariance of the unglued tensor}.}

The gauge estimate for the product of kernels is \statementref{Lemma}{lem:a-weak-chain-gauge}.

Restrict the form on $X_m$ to $Y_m$. The integral branching matrix does not increase the gauge. Restriction also introduces a scalar $\prod_\tau\tau(\bar\pi_N)^{-\beta_\tau}$ with branching exponent $\beta_\tau$. Under a weight embedding, the last coordinate of $\delta_m^{-1}$ is $N^{-1}$ or $1$, with the first corresponding to $\bar w_\tau\mid\pi_N$. Branching raises the action of that coordinate to power $\beta_\tau$. Since $|\tau(\bar\pi_N)|_\iota$ equals $N^{-1}$ or $1$ by the same condition, the absolute value of this branching factor cancels that of the scalar in restriction. Thus the standard-lattice gauge on $Y_m$ is controlled by the gauge of $X_m$ at $\delta_m$, still at most $N^{2mfb}$.\leanfn{See \href{https://github.com/lsernau/grh-lean/blob/72fead102f663d5c7a407964a13dc0cf9d977207/grh/src/New/UpperBound/Divisibility/ChainValuation.lean\#L286}{The branching matrix and scalar in restriction}, \href{https://github.com/lsernau/grh-lean/blob/72fead102f663d5c7a407964a13dc0cf9d977207/grh/src/New/UpperBound/Divisibility/ValuationAtLevel.lean\#L541}{The final gauge of the composite lift}.}
\end{proof}

\thmsource{Lean source: \leansourcelink{src/New/UpperBound/Divisibility/LocalDivisibility.lean}{482}.}
\begin{lemma}[Uniform local losses for the composite lift]
\label{lem:a-uniform-divisibility}
\emitthmsource
There is a family $e_p$ of nonnegative reals with $e_p=0$ whenever $p\ne N$ and either $\chi_K(p)=1$ or $x_1\le p$ holds. For every finite index set $J$ we have $\sum_{p\in J}e_p\log p\le B$. Fix a prime $p$ and a $p$-adic place above it. Let $M$ be the integral ramified source module there, with $u_0=\dim_{\mathbb C}U_0$. There exist $j\in\mathbb N$ and a submodule $M'\subseteq M$ with $j\le e_pu_0$ such that every finite family in $M$ can be sent into $M'$ by an integral matrix with determinant dividing $p^j$. For each source form in $M'$, every twist with level coprime to $p$, and all its chain data, the composite lift has $p$-adic Weyl gauge at most $p^{e_p-\gamma_p}$ at each finite adelic point.
\end{lemma}
\begin{proof}
Defining the allowable losses
\[
 e_p=
 \begin{cases}
 2d_p, &p\in Q,\\
 2mfb,&p=N,\\
 0,&\text{otherwise},
 \end{cases}
 \qquad
 d_p=f\ell^2\frac ap
       \begin{cases}m,&p\equiv\pm1\pmod m,\\ 1,&\text{otherwise}.
       \end{cases}
\]
Here $d_p$ is used only when $p\in Q$, and $N\ge x_1$ implies $Q$ and $\{N\}$ are disjoint. Nonnegativity and the support condition follow from the definition. Write $\chi_K$ for the characteristic of $K=\mathbb Q(\sqrt{-d})$. In $Q$ we require $\chi_K(p)\ne1$.\leanfn{See \href{https://github.com/lsernau/grh-lean/blob/72fead102f663d5c7a407964a13dc0cf9d977207/grh/src/New/UpperBound/Size/Estimates.lean\#L314}{Definition of allowable losses}.}

Write $\lambda=\log x_1$ and $T_0=4fa\ell^2\lambda$. If $p\in Q$ and $p\not\equiv\pm1\pmod m$, take $c=p$ to get $e_p\log p\le T_0/c$. If $p\in Q$ and $p\equiv\pm1\pmod m$, write $p=mc\pm1$ with $1\le c\le p$ and $mc\le p+1\le2p$ to get $e_p\log p\le T_0/c$ again. The level prime $N$ corresponds to $c=1$: by $\log N\le30m\ell$, $m^2\le4x$, and $a=xb$,
\[
 e_N\log N\le60m^2fb\ell\le240fa\ell\le T_0 .
\]
The last step uses $\ell\lambda\ge60$ from the scales. For fixed $c$ this bound controls at most the four indices $N,mc-1,c,mc+1$, so for every finite set $J$,
\[
 \sum_{p\in J} e_p\log p\le
 4T_0\sum_{c=1}^{\lfloor x_1\rfloor}\frac1c
 \le16fa\ell^2\lambda(1+\lambda)\le fa\ell^3=B.
\]
The last step uses $\ell\ge\lambda^6$ and $\lambda\ge3$, for instance $16(1+\lambda)\le(64/3)\lambda$ and $\lambda^4\ge81$.\leanfn{See \href{https://github.com/lsernau/grh-lean/blob/72fead102f663d5c7a407964a13dc0cf9d977207/grh/src/New/UpperBound/Size/Estimates.lean\#L326}{A bound for harmonic partial sums}, \href{https://github.com/lsernau/grh-lean/blob/72fead102f663d5c7a407964a13dc0cf9d977207/grh/src/New/Reduction/Data/LevelAndTwists.lean\#L228}{A size bound for the level prime}.}

Now fix a prime $p$ and a $p$-adic place above it, and let $M_0$ be the full integral source module there. First choose $S_0$ so any finite family in $M_0$ is sent into $S_0$ by a matrix of determinant dividing $p^j$. Then verify the local lift estimate for ramified vectors in $S_0$. We construct $S_0$ in four cases.
\begin{enumerate}
\item If $p\notin Q\cup\{N\}$, then $d<x_1$ and $\chi_K(d)=0$ force $p\ne d$, and $\gamma_p=0=e_p$. \statementref{Lemma}{lem:a-weak-chain-replacement} gives weak chain data for the same lift, and \statementref{Lemma}{lem:a-weak-chain-integral-away} gives integrality. Take $S_0=M_0$ and $j=0$.
\item If $p\in Q$ and $p$ is nonexceptional, still take $S_0=M_0$ and $j=0$. \statementref{Lemma}{lem:a-weak-chain-unmarked-valuation} gives the weak-chain local bound of exponent $mC_p(\varnothing)+d_p-\gamma_p$. By $mC_p(\varnothing)\le d_p$ of \statementref{Lemma}{lem:a-nonexceptional-collision} the exponent is at most $e_p-\gamma_p$.
\item If $p\in Q$ and $p$ is exceptional, take a full marking $\Lambda$ from \statementref{Lemma}{lem:a-full-mark}. Let $Y_\Lambda$ be the product of scaled Hecke operators at the marked places, put $j_p=\lceil 2mfa/p\rceil$, and take $S_0=Y_\Lambda^m M_0+p^{j_p}M_0$. \statementref{Lemma}{lem:a-exceptional-hecke-colength} gives $j\le d_pu_0$ with correction determinant dividing $p^j$. By \statementref{Lemma}{lem:a-full-mark-collision} the unmarked places have no collisions, so \statementref{Lemma}{lem:a-full-mark-valuation} gives the bound $p^{2d_p-\gamma_p}=p^{e_p-\gamma_p}$.
\item If $p=N$, take $S_0$ to be the preimage in the source space of the integral module for the raising operator in \statementref{Lemma}{lem:a-level-up-colength}. That lemma gives $j\le mfbu_0\le e_pu_0$, and \statementref{Lemma}{lem:a-level-composite-gauge} gives the lift bound $N^{2mfb}$. By $\gamma_N=0$ this equals $p^{e_p-\gamma_p}$.
\end{enumerate}
The local bounds hold simultaneously for all the specified twists and chain data. Finally put $M'=M\cap S_0$. A finite family from $M$ is also one in $M_0$, so an appropriate matrix sends it into $S_0$. Since the matrix is integral and $M$ is a module, the corrected family remains in $M$ and thus lies in $M'$ with the same determinant bound. The local lift estimate then applies to its vectors.\leanfn{See \href{https://github.com/lsernau/grh-lean/blob/72fead102f663d5c7a407964a13dc0cf9d977207/grh/src/New/UpperBound/Divisibility/LocalDivisibility.lean\#L402}{The four local cases and the uniform bound}.}
\end{proof}

\section{The seesaw identity}
\label{sec:a-seesaw}

The classical seesaw identity interchanges the pairings along two theta-lifting paths~\cite{Kudla1984}. For unfolding of nilpotent orbit periods under a single theta lift see Wang~\cite[Theorem~4.5]{Wang2025}. To compare the two chains in~\eqref{eq:lift-comparison}, we unfold the constant term of the Bessel period. The unipotent character determines the rational string classes that survive. Absolute integrability and lattice support then express the expansion as a restriction of the upper-chain lift.

\thmsource{Lean source: \leansourcelink{src/New/UpperBound/Seesaw/IntegralAsCosetSum.lean}{353}.}
\begin{lemma}[A weighted section integral for the Bessel period]
\label{lem:a-seesaw-cut-integral}
\emitthmsource
Let $G=G_\infty\times G_f$ be a lower-chain unitary group, $R\subset G$ the corresponding \emph{Bessel subgroup}, $R_f$ its finite part, and $\Gamma_R$ its rational point group. Let $A,J$ be the finite level and $J^\circ$ the standard integral level. Take a finite subset $\Omega$ of $G_f/J$ stable under $R_f\cap A$. For $\omega\in G_f/J$ let $\mu(\omega)$ count the classes in $R_f/(R_f\cap A)$ of elements $u\in R_f$ satisfying $u^{-1}\omega\in\Omega$, with the count defined to be zero if infinite.

Let $H$ be continuous, left $\Gamma_R$-invariant and integrable on $\Gamma_R\backslash G$, vanishing off points $t$ satisfying $\mu(\bar t_f)\ne0$. Suppose when $\bar t_f\in\Omega$ its \emph{Bessel period} satisfies
\[
 \int_{\Gamma_R\backslash R}H(rt)\,d[r]
     =[J^\circ:J]\, a(t_\infty)b(\bar t_f),
\]
Then there is a constant $\kappa$ depending only on $R,A,a$ and the measure such that
\[
 \int_{\Gamma_R\backslash G} H(t)\,d[t]
   =\kappa\sum_{\omega\in\Omega}\frac{b(\omega)}{\mu(\omega)} ,
\]
where the reciprocal is taken to be zero when the multiplicity vanishes.
\end{lemma}
\begin{proof}
Normalize Haar measure $\eta$ on $R_f$ so that $R_f\cap A$ has volume $1$. Put $\psi=\mathbf 1_\Omega/\mu$ and regard it as a function on $G_f/J$. Fix $\omega$. Since $\Omega$ is stable under $R_f\cap A$, the elements $u$ satisfying $u^{-1}\omega\in\Omega$ form a union of $R_f\cap A$ cosets. Left translation by $R_f$ identifies the sets being counted, so $\mu$ is constant on each $R_f$ orbit. Thus when $\mu(\omega)>0$,
\[
 \int_{R_f}\psi(u^{-1}\omega)\,d\eta(u)
   =\frac{1}{\mu(\omega)}\eta\{u:u^{-1}\omega\in\Omega\}
   =1.
\]
When $\mu(\omega)=0$ the integrand vanishes identically. The orbit integral of the finite section weight is thus the indicator of the covered set.\leanfn{See \href{https://github.com/lsernau/grh-lean/blob/72fead102f663d5c7a407964a13dc0cf9d977207/grh/src/New/UpperBound/Seesaw/IntegralAsCosetSum.lean\#L200}{Orbit integral of the finite section weight}.}

Choose a nonnegative measurable cutoff $\beta$ on $G$ with $\int_R\beta(rt)\,dr=1$ for each $t$. Such a cutoff exists by properness of left translation by a closed subgroup. Put
\[
 \alpha(v)=\int_{R_f}\beta(v,u)\,d\eta(u),\qquad
 w(t)=\alpha(t_\infty)\psi(\bar t_f).
\]
To compute the orbit integral of $w$, use Tonelli to exchange the integrals over $R$ and $R_f$. With $u$ fixed, left translate by $r\mapsto(1,u^{-1})r$ in $R$. Exchange the integrals again, and with $r$ fixed, left translate by $u\mapsto r_fu$ in $R_f$. These substitutions give\leanfn{See \href{https://github.com/lsernau/grh-lean/blob/72fead102f663d5c7a407964a13dc0cf9d977207/grh/src/New/UpperBound/Seesaw/IntegralAsCosetSum.lean\#L278}{From the cutoff to the product weight}.}
\[
 \int_R w(rt)\,dr
   =\int_{R_f}\psi(u^{-1}\bar t_f)
                 \left(\int_R\beta(r(t_\infty,u))\,dr\right)d\eta(u)
   =\mathbf 1_{\{\mu(\bar t_f)\ne0\}}.
\]

Since $H$ vanishes outside the covered set, insert the orbit normalization and expand the quotient integral in stages along $\Gamma_R\subset R\subset G$ to get
\[
 \int_{\Gamma_R\backslash G}H(t)\,d[t]
   =\int_G w(t)\left(\int_{\Gamma_R\backslash R}H(rt)\,d[r]\right)dt .
\]
The map $(t,r)\mapsto rt$ pushes $w(t)\,dt\,dr$ forward to the restriction of Haar measure to the covered set. Left $\Gamma_R$-invariance makes the decomposition descend to the rational-point quotient, and integrability justifies the integral identity there. The right-hand side samples only at $\bar t_f\in\Omega$, so the period formula applies.\leanfn{See \href{https://github.com/lsernau/grh-lean/blob/72fead102f663d5c7a407964a13dc0cf9d977207/grh/src/New/UpperBound/Seesaw/IntegralAsCosetSum.lean\#L318}{Successive unfolding on the rational-point quotient}.}

The product weight $w$ factors into archimedean and finite parts. By the normalization of Haar measure on $G_f$, the volume of $J^\circ$ is $1$. Each $J$ coset therefore has volume $[J^\circ:J]^{-1}$, so
\[
 \int_G w(t)[J^\circ:J]a(t_\infty)b(\bar t_f)\,dt
   =\left(\int_{G_\infty}\alpha(v)a(v)\,dv\right)
      \sum_{\omega\in\Omega}\frac{b(\omega)}{\mu(\omega)} .
\]
The integral in parentheses depends only on the fixed $R,A,a$ and the measure. Take it as $\kappa$.\leanfn{See \href{https://github.com/lsernau/grh-lean/blob/72fead102f663d5c7a407964a13dc0cf9d977207/grh/src/New/UpperBound/Seesaw/IntegralAsCosetSum.lean\#L367}{Cancellation of level volumes and the archimedean constant}.}
\end{proof}

\thmsource{Lean source: \leansourcelink{src/New/UpperBound/Seesaw/GoodDatum.lean}{526}.}
\begin{lemma}[Vacuum expansion of the finite evaluation tensor]
\label{lem:a-seesaw-vacuum-expansion}
\emitthmsource
Write $R=M^{f n_m^2}$. There exists a common constant $C\in\mathbb C$ with the following property for every tame character $\vartheta$. Write $q$ for the modulus norm of twist $\vartheta$. If the lower-chain finite variables $y_j$ are integral away from $q$, then the pure tensor $E_\vartheta(\bar y)$ obtained by evaluating the double steps has a finite expansion
\[
 E_\vartheta(\bar y)
  =C\sum_x r(x)\operatorname{Vac}_{\rm gl}(x),
\]
Here $x$ is an admissible translation and $R^{2m}r(x)$ is integral at every finite place away from $q$. If the support of the finite conductor is empty, then for each $\sigma\in\operatorname{Aut}(\mathbb C)$ we have $\sum_x|\sigma(r(x))|\le1$.
\end{lemma}
\begin{proof}
Let $\beta=\pi_M$ be the modulus element for the principal level in $E$, with $A=\operatorname{Nm}((\beta))=M^f$. Fix double step $Y_{2i}\to Y_{2i+1}\to Y_{2i+2}$ of index $i$. For a single edge $Y_j\to Y_{j+1}$, the finite Schwartz datum translated at $y_j,y_{j+1}$ is obtained by applying the finite Weil action to the initial vacuum expansion and multiplying by the character value of the determinant of the point. Its absolute value is $1$ at finite places away from $q$ and at all $\operatorname{Aut}(\mathbb C)$-conjugates. Since $y_j,y_{j+1}$ is integral away from $q$, the Weil action sends each vacuum translate to a finite linear combination on the same admissible set with coefficients integral away from $q$. The averaging denominator $A^{\dim Y_j\dim Y_{j+1}}$ of the initial datum is also a denominator of the translated expansion. With empty modulus support the sum of absolute values of the coefficients at each conjugate is at most $1$.\leanfn{See \href{https://github.com/lsernau/grh-lean/blob/72fead102f663d5c7a407964a13dc0cf9d977207/grh/src/New/UpperBound/Seesaw/GoodDatum.lean\#L359}{Finite Weil translation of edge data and the coefficients}.}

Take the arrow product of the edge data. It is bilinear, and the arrow product of two vacuum translates is another vacuum translate times a root of unity. The new translation belongs to the finite adelic coordinate space
\[
 \operatorname{Hom}\bigl(Y_{2i+1},\,Y_{2i}^{-}\oplus Y_{2i+2}\bigr)
\]
and becomes integral away from $q$ upon multiplying the two blocks by $\beta$ and $\bar\beta$, respectively. By bilinearity and the unit absolute values of roots of unity, use the product of the two edge denominators for the denominator. The bound on the sum of conjugate absolute values is still $1$.\leanfn{See \href{https://github.com/lsernau/grh-lean/blob/72fead102f663d5c7a407964a13dc0cf9d977207/grh/src/New/UpperBound/Seesaw/GoodDatum.lean\#L446}{Admissible expansion of the arrow product}.}

Apply double-step evaluation to this expansion. Its local evaluation theorem applies to each vacuum translate and writes it as a finite expansion of glued double-step vacuum translates with a common constant $C_i$. The two translations in the output have allowed denominator $\bar\beta$ away from $q$. The evaluation denominator counts a class of matrices modulo $\beta$ with at most $2n_m$ rows and $n_m$ columns. Writing $A^{d_i}$ for the denominator, $d_i\le2n_m^2$. Combining with the denominators of the two edges gives divisibility of the full denominator by
\[
 A^{n_m^2+n_m^2+2n_m^2}=R^4 .
\]
With empty modulus support, multiply the absolute value of each original coefficient by the absolute coefficient sum in the corresponding local expansion and add. The result for each conjugate is still at most $1$.\leanfn{See \href{https://github.com/lsernau/grh-lean/blob/72fead102f663d5c7a407964a13dc0cf9d977207/grh/src/New/UpperBound/Seesaw/GoodDatum.lean\#L418}{Counting residue matrices and bounding the denominator}, \href{https://github.com/lsernau/grh-lean/blob/72fead102f663d5c7a407964a13dc0cf9d977207/grh/src/New/UpperBound/Seesaw/GoodDatum.lean\#L481}{Local expansion of double-step evaluation}, \href{https://github.com/lsernau/grh-lean/blob/72fead102f663d5c7a407964a13dc0cf9d977207/grh/src/New/UpperBound/Seesaw/GoodDatum.lean\#L155}{The local evaluation theorem}.}

To pass to the integral vacua of the chain data, let $\kappa_j$ be the nonzero multiple giving the integral vacuum at edge $j$ from the ordinary one. Linearity of translation and gluing changes the common factor $C_i$ for double step $i$ to $c_i=C_i/(\kappa_{2i}\kappa_{2i+1})$ when using integral vacua. Take the tensor product of the expansions for all $m/2$ double steps. Each term is an admissible tuple of translations and its coefficient is the product of the double-step coefficients. The common constant is $C=\prod_i c_i$ and the total denominator divides $(R^4)^{m/2}=R^{2m}$. If the finite modulus support is empty, the sum of the absolute values under each $\sigma$ is at most the product of the separate sums, hence at most $1$. Combining repeated tuples preserves the bounds and gives the expansion.\leanfn{See \href{https://github.com/lsernau/grh-lean/blob/72fead102f663d5c7a407964a13dc0cf9d977207/grh/src/New/UpperBound/Seesaw/GoodDatum.lean\#L531}{Integral vacuum normalization and the double-step tensor product}.}
\end{proof}

\thmsource{Lean source: \leansourcelink{src/Lib/Weil/Periods/BesselIntegral.lean}{200}.}
\begin{lemma}[Eichler transformations in a graded unipotent group]
\label{lem:a-bessel-eichler}
\emitthmsource
Let $\mathcal Y$ be a string grading of the Hermitian space $V$, with $g(a)$ the degree of basis index $a$, $\sigma$ the string reversal, and $J$ the Gram matrix in the rational basis with coefficients viewed as adeles. The only nonzero pairing of $J$ at index $a$ is with $\sigma(a)$, and $g(\sigma(a))=-g(a)$. Given an integer $k$, let $v$ have support in degrees strictly greater than $\max(0,k)$ and $x$ have support in degree $-k$. Put $c=x^\ast Jx/2$ with $^\ast$ the conjugate transpose. Then the matrix
\[
 U=1+v x^\ast J-(x+c v)v^\ast J
\]
is the matrix in rational coordinates of an element of the unipotent group of $\mathcal Y$.
\end{lemma}
\begin{proof}
The degree pairing gives $v^\ast Jv=0$. Since $x$ can pair only with degree $k$, we also have $x^\ast Jv=0$. Put $B=U-1$. By $J^\ast=J$, the element $c$ is conjugation-invariant. Direct multiplication gives
\[
 B^\ast J+JB=-2c\,Jv v^\ast J,\qquad
 B^\ast JB=(x^\ast Jx)\,Jv v^\ast J ,
\]
In the second formula the terms with $v^\ast Jv$ or $x^\ast Jv$ vanish. This gives $U^\ast JU=J$, so $U$ is unitary.

Now consider degrees. The term $v x^\ast J$ sends source degree $k$ to strictly higher degree. The term $x v^\ast J$ has nonzero source degrees strictly below $-k$ and target degree $-k$. Finally $v v^\ast J$ sends negative degrees to positive ones. Thus $U-1$ strictly raises degree, the unipotent condition for the string grading.\leanfn{See \href{https://github.com/lsernau/grh-lean/blob/72fead102f663d5c7a407964a13dc0cf9d977207/grh/src/Lib/Adelic/Witt/WittTheorem.lean\#L736}{The Eichler matrix}, \href{https://github.com/lsernau/grh-lean/blob/72fead102f663d5c7a407964a13dc0cf9d977207/grh/src/Lib/Weil/Periods/BesselIntegral.lean\#L204}{The Eichler transformation is unipotent}.}
\end{proof}

\thmsource{Lean source: \leansourcelink{src/Lib/Weil/Periods/BesselIntegral.lean}{541}.}
\begin{lemma}[Equality of Gram entries from equality of unipotent phases]
\label{lem:a-bessel-eichler-gram}
\emitthmsource
Let $\mathcal Y,\mathcal Y'$ be string gradings of $V,V'$ of lengths $s,\max(s-1,0)$, with string map $T$. Keep $J,g,\sigma$, write $J'$ for the target Gram matrix, and $^\dagger$ for the Hermitian adjoint. Fix integer weights $d'(b)$ on the target basis. Suppose $w$ passes the following phase test. For every source unipotent $U$, whenever $U-1$ has degree at least $2$ and $w(U-1)$ has degree at least $1$ relative to the target weights $d'$, we have
\[
 \mathfrak b(TU,T)=\mathfrak b(wU,w),
\]
where $\mathfrak b$ is the Weil bicharacter on the Hom space. Fix a source index $k$ and a source vector $v$ satisfying
\[
 \begin{aligned}
 g(a)<\max(1,2+g(k))&\ \Longrightarrow\ v_a=0,\\
 d'(b)\le\max(-1,g(k))&\ \Longrightarrow\ (wv)_b=0,\\
 d'(b)+g(a)\le0&\ \Longrightarrow\ 
                        w_{b,\sigma(k)}=0\ \text{or}\ v_a=0 .
 \end{aligned}
\]
Then $(T^\dagger T v)_k=(w^\dagger wv)_k$.
\end{lemma}
\begin{proof}
Take any adele $t$ and let $x$ be supported only at index $\sigma(k)$, with coordinate
\[
  \overline{\jmath{}^{-1}(J^{-1})_{k,\sigma(k)}\,t},
\]
Here $\jmath{}$ is the nonzero purely imaginary parameter of the Weil data, so $\bar{\jmath{}}=-\jmath{}$. The first support condition allows \statementref{Lemma}{lem:a-bessel-eichler} for $v,x$, and the resulting $U-1$ has degree at least $2$. In fact, $v x^\ast J$ sends degree $g(k)$ to at least $g(k)+2$, $x v^\ast J$ sends degrees at most $-g(k)-2$ to $-g(k)$, and $v v^\ast J$ sends degrees at most $-1$ to at least $1$.

For the degree condition on $w(U-1)$, check its three terms. For $wv x^\ast J$ the source degree is $g(k)$ and the second condition gives $d'$ at least $g(k)+1$ on every nonzero target row. For $wx v^\ast J$, if $v_a$ contributes, the source degree is $-g(a)$ and the third condition gives $d'$ at least $1-g(a)$ on the nonzero rows. For $wv v^\ast J$ the second condition gives $d'$ at least $0$ on each nonzero row and the degree of a nonzero source column is at most $-1$. The phase test applies.\leanfn{See \href{https://github.com/lsernau/grh-lean/blob/72fead102f663d5c7a407964a13dc0cf9d977207/grh/src/Lib/Weil/Periods/BesselIntegral.lean\#L550}{An Eichler transformation for the phase test}.}

For a matrix $y$ put $C_y=y^\ast J'y$. Substitute $U$ into the trace formula for the Weil bicharacter to get
\[
 \operatorname{tr}(J^{-1}C_yU)
 =\operatorname{tr}(J^{-1}C_y)
      +x^\ast C_yv-v^\ast C_yx-c\,v^\ast C_yv .
\]
The first and last terms are conjugation-invariant and the two middle terms are opposite conjugates. The bicharacter evaluates the standard additive character at the trace times $\jmath{}/2$. This character is invariant under replacing its argument with its average with its conjugate, so
\[
 \mathfrak b(yU,y)=\psi_E(\jmath{} x^\ast C_yv).
\]
Row $k$ of the inverse source Gram matrix $J^{-1}$ is nonzero only at $\sigma(k)$. The choice of $x$ gives $\jmath{} x^\ast C_yv=t(J^{-1}C_yv)_k=t(y^\dagger y v)_k$. Taking $y=T,w$ in turn, the phase test gives for each $t$
\[
  \psi_E\!\left(t\bigl((T^\dagger Tv)_k-(w^\dagger wv)_k\bigr)\right)=1 .
\]
Nondegeneracy of the standard additive pairing makes the difference zero.\leanfn{See \href{https://github.com/lsernau/grh-lean/blob/72fead102f663d5c7a407964a13dc0cf9d977207/grh/src/Lib/Weil/Periods/BesselIntegral.lean\#L572}{The phase trace and the additive pairing}.}
\end{proof}

\thmsource{Lean source: \leansourcelink{src/Lib/Weil/Periods/BesselIntegral.lean}{765}.}
\begin{lemma}[Rational string representatives from the unipotent test]
\label{lem:a-bessel-rational-chord-representative}
\emitthmsource
Keep the grading and notation of \statementref{Lemma}{lem:a-bessel-eichler-gram}. Let $w$ be a rational map from $V$ to $V'$ satisfying its phase test with target weight $d'=0$. There is a rational unitary transformation $\delta$ of $V'$ such that $\delta w$ agrees with $T$ on each positive-degree source basis vector. In particular,
\[
 (\delta w)_{ba}=T_{ba}
 \quad\text{whenever}\quad
        g(a)>0,\quad g'(b)<g(a),
\]
where $g'$ is the target grading.
\end{lemma}
\begin{proof}
Let $V_{>0}$ be spanned by the positive-degree basis vectors. First prove that $w|_{V_{>0}}$ is injective. If there were a nonzero $v\in V_{>0}$ with $wv=0$, take a string index $a$ of minimal degree among the nonzero coordinates of $v$. Positive degrees on the string decrease by $2$ each step. The next string index $k=a+1$ therefore exists and satisfies $g(k)=g(a)-2$, and every nonzero coordinate has degree at least $\max(1,g(k)+2)$. Since $wv=0$ and $v$ vanishes in nonpositive degrees, the three conditions of \statementref{Lemma}{lem:a-bessel-eichler-gram} hold, giving $(T^\dagger T v)_k=(w^\dagger wv)_k=0$.

But $T^\dagger T$ is the shift on the source string. It sends each nonterminal string vector to a nonzero multiple of the next and kills the definite block and the terminal vector. On positive-degree columns the multiple is $-Nc_{\mathcal Y}^2$ when the degree is $1$ and $c_{\mathcal Y}$ in the other cases, with $N,c_{\mathcal Y}$ nonzero. Thus $(T^\dagger Tv)_k$ is a nonzero multiple of $v_a$, a contradiction. Hence $w|_{V_{>0}}$ is injective. Also $T$ carries the positive-degree source vectors to distinct nonzero target string vectors with the same indices, so $T$ is injective on the positive-degree part as well.\leanfn{See \href{https://github.com/lsernau/grh-lean/blob/72fead102f663d5c7a407964a13dc0cf9d977207/grh/src/Lib/Weil/Periods/Grading.lean\#L875}{The adjoint product of the string map is a shift}, \href{https://github.com/lsernau/grh-lean/blob/72fead102f663d5c7a407964a13dc0cf9d977207/grh/src/Lib/Weil/Periods/BesselIntegral.lean\#L722}{Independence of the positive-degree columns}.}

Compare the Gram matrices of the two restrictions. For positive-degree index $a,b$ take $v=e_a$ and $k=\sigma(b)$ in \statementref{Lemma}{lem:a-bessel-eichler-gram}. Since $g(a),g(b)\ge1$, the first support condition holds. By $g(k)=-g(b)\le-1$ the second condition with target weights zero is vacuous. The third says exactly that $e_a$ vanishes in nonpositive degree. Thus $(T^\dagger T e_a)_{\sigma(b)}=(w^\dagger w e_a)_{\sigma(b)}$. Multiply by entry $(b,\sigma(b))$ of the source Gram matrix to obtain at $(b,a)$
\[
 (T^\ast J'T)_{ba}=(w^\ast J'w)_{ba}.
\]
Here $y^\ast J'y=Jy^\dagger y$ is used along with the fact that row $b$ of $J$ is nonzero only at $\sigma(b)$. The number field embeds injectively into its adele ring, so this is an identity of rational-coordinate matrices.

Thus $w v\mapsto Tv$ from $w(V_{>0})$ to $T(V_{>0})$ is an isometric isomorphism of rational subspaces. By Witt extension it extends to a rational unitary transformation $\delta$ of $V'$. By construction $\delta w$ and $T$ agree in each positive-degree source column, giving the required string agreement.\leanfn{See \href{https://github.com/lsernau/grh-lean/blob/72fead102f663d5c7a407964a13dc0cf9d977207/grh/src/Lib/Adelic/Unitary/RationalPoints.lean\#L173}{Witt extension for rational unitary spaces}, \href{https://github.com/lsernau/grh-lean/blob/72fead102f663d5c7a407964a13dc0cf9d977207/grh/src/Lib/Weil/Periods/BesselIntegral.lean\#L774}{A rational representative agreeing with the string map}.}
\end{proof}

\thmsource{Lean source: \leansourcelink{src/Lib/Weil/Periods/BesselIntegral.lean}{642}.}
\begin{lemma}[Target transformations preserving string agreement]
\label{lem:a-bessel-chord-stabilizer}
\emitthmsource
Suppose $w$ and the string map $T$ agree in negative-degree entries of every positive-degree source column, meaning $w_{ba}=T_{ba}$ when $g(a)>0$ and $g'(b)<g(a)$. If $\gamma$ is in the source Bessel group and $\delta$ in the target unitary group, then $\delta w\gamma$ still satisfies the string agreement if and only if $\delta$ belongs to the target Bessel group. Matrices use rational coordinates.
\end{lemma}
\begin{proof}
If $\delta$ is in the target Bessel group, the matrices $\delta,\gamma$ do not lower degree. For a positive-degree source column, each component of $(w-T)\gamma$ has degree at least the degree of that column, which remains true after left multiplication by $\delta$. Also,
\[
 \delta T\gamma-T=(\delta-1)T\gamma+T(\gamma-1)
\]
has nonnegative degree. Indeed $T$ joins corresponding string coordinates and lowers degree by one, while $\delta-1$ and $\gamma-1$ strictly raise degree on the relevant string coordinates. Thus $\delta w\gamma-T$ vanishes in negative-degree entries of positive-degree source columns.\leanfn{See \href{https://github.com/lsernau/grh-lean/blob/72fead102f663d5c7a407964a13dc0cf9d977207/grh/src/Lib/Weil/Periods/Grading.lean\#L817}{Left and right Bessel actions preserve negative-degree terms of the string map}.}

Conversely, suppose $\delta w\gamma$ satisfies string agreement. The preceding argument with $\delta=1$ shows that $w\gamma$ also satisfies it. To show that $\delta$ is in the target Bessel group, its column criterion asks for $(\delta-1)_{bk}=0$ whenever column $k$ is a string index with $g'(k)\ge0$ and row $b$ satisfies $g'(b)\le g'(k)$.

Fix $b$ and induct over the indicated columns in decreasing degree. For $k$ take the matching source string index $a$. Then $g(a)=g'(k)+1>0$ and $T_{ka}\ne0$, and the column of $T$ is nonzero only at $k$. For columns $i$ with $g'(i)>g'(k)$ the induction hypothesis gives $(\delta-1)_{bi}=0$. For columns $i$ with $g'(i)\le g'(k)$, use $g'(i)<g(a)$ and string agreement to get $(w\gamma)_{ia}=T_{ia}$. Hence
\[
 ((\delta-1)w\gamma)_{ba}=(\delta-1)_{bk}T_{ka}.
\]
Here $\delta w\gamma$ and $w\gamma$ both equal $T$, so the left-hand side is zero. Cancel the nonzero rational factor $T_{ka}$ to get $(\delta-1)_{bk}=0$. Induction and the column criterion imply that $\delta$ is in the target Bessel group.\leanfn{See \href{https://github.com/lsernau/grh-lean/blob/72fead102f663d5c7a407964a13dc0cf9d977207/grh/src/Lib/Weil/Periods/BesselIntegral.lean\#L617}{The unique nonzero entry of the matching column}, \href{https://github.com/lsernau/grh-lean/blob/72fead102f663d5c7a407964a13dc0cf9d977207/grh/src/Lib/Weil/Periods/BesselIntegral.lean\#L645}{The group preserving string agreement}.}
\end{proof}

\thmsource{Lean source: \leansourcelink{src/Lib/Weil/Periods/BesselIntegral.lean}{672}.}
\begin{lemma}[The class passing the unipotent test in the target grading]
\label{lem:a-bessel-surviving-chord-class}
\emitthmsource
Keep the two string gradings. Suppose $w$ and $T$ agree on the negative-degree entries of positive-degree source columns. If $w$ also passes the phase test of \statementref{Lemma}{lem:a-bessel-eichler-gram} with target grading $d'=g'$, then $w-T$ has nonnegative degree: $w_{ba}=T_{ba}$ for all index pairs satisfying $g'(b)<g(a)$.
\end{lemma}
\begin{proof}
Induct on the target row degree $g'(b)$ in increasing order. For $g(a)>0$ use string agreement, leaving the case $g'(b)<g(a)\le0$. Let $b^\vee=\sigma'(b)$ denote the target dual index, of degree $-g'(b)>0$, so $b^\vee$ lies on the target string. Take the matching source string index $i$. Then
\[
 g(i)=1-g'(b),\qquad T_{b^\vee i}\ne0,\qquad
       T_{u i}=0\quad(u\ne b^\vee).
\]
Put $k=\sigma(a)$ and $v=e_i$ and check the three support conditions of \statementref{Lemma}{lem:a-bessel-eichler-gram}.

First $g(i)\ge2-g(a)=2+g(k)$ and $g(i)>0$. Next, $g'(u)\le\max(-1,g(k))=-g(a)$ gives $g'(u)<-g'(b)=g'(b^\vee)$. With $g(i)>0$ and string agreement this gives $w_{ui}=T_{ui}=0$. For the third condition, consider a source index $i$. If $g'(u)+g(i)\le0$ then $g'(u)\le g'(b)-1<g'(b)$, so row induction gives $w_{ua}=T_{ua}$. By $g'(u)<g'(b)<g(a)$ we get $g'(u)\le g(a)-2$, and since $T$ lowers degree by just one, $w_{ua}=0$.

This gives
\[
 (w^\dagger w)_{ki}=(T^\dagger T)_{ki}.
\]
For $u$ with $g'(u)\le g'(b^\vee)$, use $g(i)=g'(b^\vee)+1$ and string agreement to get $w_{ui}=T_{ui}$. For $u$ with $g'(u)>g'(b^\vee)$, the dual row $\sigma'(u)$ has degree strictly below $g'(b)$. The induction hypothesis and the fact that $T$ lowers degree by one give $w_{\sigma'(u),a}=0$. The entry formula for the adjoint gives $(w^\dagger)_{k u}=0$. The left-hand side now has only $(w^\dagger)_{k b^\vee}T_{b^\vee i}$ and the right only $(T^\dagger)_{k b^\vee}T_{b^\vee i}$. Cancel the nonzero factor $T_{b^\vee i}$ and use the adjoint entry formula again to get $w_{ba}=T_{ba}$, completing row induction.\leanfn{See \href{https://github.com/lsernau/grh-lean/blob/72fead102f663d5c7a407964a13dc0cf9d977207/grh/src/Lib/Weil/Periods/Grading.lean\#L453}{Dual entries of the adjoint matrix}, \href{https://github.com/lsernau/grh-lean/blob/72fead102f663d5c7a407964a13dc0cf9d977207/grh/src/Lib/Weil/Periods/BesselIntegral.lean\#L680}{The second orbit selection leaves only the string class}.}
\end{proof}

\thmsource{Lean source: \leansourcelink{src/Lib/Weil/Periods/Unfolding.lean}{432}.}
\begin{lemma}[An orbit expansion of the Bessel period by constant terms]
\label{lem:a-bessel-orbit-unfold}
\emitthmsource
Let $V,V'$ be two Hermitian spaces with gradings $\mathcal Y,\mathcal Y'$ of string lengths $s,\max(s-1,0)$, respectively. Let $R,R'$ be their Bessel groups, $\Gamma'\subset \Ugrp(V')(\mathbb A)$ their rational point groups, and $R'_{\mathbb Q}=\Gamma'\cap R'$. Suppose $W$ is continuous and left rationally invariant on $R$, and equivariant for the graded unipotent character on elements raising degree by at least $2$. Let $T$ be the string map from $\mathcal Y$ to $\mathcal Y'$, $\Phi$ Weil data on $\operatorname{Hom}(V,V')$, and $\omega$ the fixed twisted Weil action. For $h\in \Ugrp(V')(\mathbb A)$ define
\[
 A(h)=\int_{R(\mathbb Q)\backslash R}W(r)
        \operatorname{CT}_{T}\bigl(\omega(r,h)\Phi\bigr)\,dr ,
\]
where the \emph{constant term} averages over positive-degree maps for the grading of $\mathcal Y,\mathcal Y'$. Then for every $h\in \Ugrp(V')(\mathbb A)$,
\[
 \int_{R(\mathbb Q)\backslash R}
       W(r)\theta\bigl(\omega(r,h)\Phi\bigr)\,dr
   =\sum_{\gamma\in R'_{\mathbb Q}\backslash\Gamma'} A(\gamma h)
\]
and the right-hand side converges absolutely.
\end{lemma}
\begin{proof}
Set $\deg_d(b,a)=d(b)-\deg_{\mathcal Y}(a)$, first with target weight $d=0$, and take rational maps modulo the rational subspace of maps satisfying $\deg_d\ge0$. The theta function is the sum of the constant terms over these classes. Indeed, the positive-degree coordinates give an isotropic subspace. Parametrize negative-degree coordinates by the dual frame. The constant term for each class is the theta value of the mixed-model slice at that rational parameter. The Weil data norm bounds the sum of the slice norms, giving absolute convergence. On a compact covering set of representatives of $R(\mathbb Q)\backslash R$, the product of $W(r)$ and the norm of $\omega(r,1)\Phi$ is bounded. This permits termwise integration.\leanfn{See \href{https://github.com/lsernau/grh-lean/blob/72fead102f663d5c7a407964a13dc0cf9d977207/grh/src/Lib/Weil/Periods/ConstantTerm.lean\#L492}{Rational classes and mixed-model slices}, \href{https://github.com/lsernau/grh-lean/blob/72fead102f663d5c7a407964a13dc0cf9d977207/grh/src/Lib/Weil/Periods/ConstantTerm.lean\#L549}{Absolute summation of the constant terms}, \href{https://github.com/lsernau/grh-lean/blob/72fead102f663d5c7a407964a13dc0cf9d977207/grh/src/Lib/Weil/Periods/Unfolding.lean\#L215}{Termwise integration}.}

A representative $w$ passes the unipotent character test if for each source unipotent $u$, whenever $u-1$ has degree at least $2$ and $w(u-1)$ degree at least $1$, the phase of the constant term of $w$ under $u$ cancels the unipotent phase of the source function. If the test fails, some such $u$ gives the weighted term $F_w(r)$ with $F_w(ur)=cF_w(r)$, where $|c|=1$ and $c\ne1$. In any family of classes stable under the rational Bessel action, first group the sum by rational orbits. Unfold each orbit contribution as an integral of one term on the stabilizer quotient. The term is continuous and left invariant under the discrete stabilizer. The vanishing theorem for a nontrivial phase on the quotient makes the integral zero.\leanfn{See \href{https://github.com/lsernau/grh-lean/blob/72fead102f663d5c7a407964a13dc0cf9d977207/grh/src/Lib/Weil/Periods/Unfolding.lean\#L228}{The phase when the test fails}, \href{https://github.com/lsernau/grh-lean/blob/72fead102f663d5c7a407964a13dc0cf9d977207/grh/src/Lib/Weil/Periods/Unfolding.lean\#L250}{An orbit integral on the stabilizer quotient}, \href{https://github.com/lsernau/grh-lean/blob/72fead102f663d5c7a407964a13dc0cf9d977207/grh/src/Lib/Adelic/QuotientMeasure/UnfoldingInStages.lean\#L365}{Vanishing by a nontrivial phase on the quotient}.}

Write $\mathcal C_d$ for the classes satisfying string agreement, where $w$ and $T$ agree in the entries $(b,a)$ satisfying $\max(0,\deg_{\mathcal Y'}(b))<\deg_{\mathcal Y}(a)$. When $d=0$, \statementref{Lemma}{lem:a-bessel-rational-chord-representative} moves every passing $w$ into $\mathcal C_0$ by a rational target point. Thus only the rational target translates of $\mathcal C_0$ contribute to the integral. The rational source Bessel group preserves $\mathcal C_0$, and \statementref{Lemma}{lem:a-bessel-chord-stabilizer} shows that a rational target point maps classes of $\mathcal C_0$ back into $\mathcal C_0$ exactly when it belongs to $R'_{\mathbb Q}$. Thus $R'_{\mathbb Q}\backslash\Gamma'$ times $\mathcal C_0$ parametrizes the remaining classes. A rational change of variable expresses each contribution as an integral of the sum over $\mathcal C_0$ evaluated at $\gamma h$.\leanfn{See \href{https://github.com/lsernau/grh-lean/blob/72fead102f663d5c7a407964a13dc0cf9d977207/grh/src/Lib/Weil/Periods/BesselIntegral.lean\#L774}{Surviving classes move to classes satisfying string agreement}, \href{https://github.com/lsernau/grh-lean/blob/72fead102f663d5c7a407964a13dc0cf9d977207/grh/src/Lib/Weil/Periods/Unfolding.lean\#L346}{Removing the other orbits}, \href{https://github.com/lsernau/grh-lean/blob/72fead102f663d5c7a407964a13dc0cf9d977207/grh/src/Lib/Weil/Periods/Unfolding.lean\#L404}{Rational cosets and classes with string agreement}.}

Finally replace target weights $0$ by $\deg_{\mathcal Y'}$ in each block. The sum of string-agreement classes is the same Fourier coefficient for either choice. Let $v$ range over adelic maps from $V'$ to $V$ modulo rational maps, and let $v_+$ keep only entries $(a,b)$ with $\max(0,\deg_{\mathcal Y'}(b))<\deg_{\mathcal Y}(a)$. For Weil datum $\Psi$ this coefficient is
\[
 \operatorname*{mean}_{v}\,
   \frac{\bigl(H(v_+^\dagger)\Psi\bigr)(0)}
        {\psi\bigl(\jmath{}\,\operatorname{tr}(v_+T)\bigr)} ,
\]
Here $\psi$ is the standard additive character, $\jmath{}$ the nonzero imaginary parameter of the Weil data, and $H$ Heisenberg translation. For both choices of target weights, $v_+^\dagger$ has positive degree. Expand the numerator in rational classes and interchange the two translations. Class $w$ acquires character $\psi\bigl(\jmath{}\,\operatorname{tr}(v_+(w-T))\bigr)$. Averaging over $v$ keeps exactly the trivial character, hence $\mathcal C_d$. The block integral is therefore independent of the choice of target weights. With $\deg_{\mathcal Y'}$, \statementref{Lemma}{lem:a-bessel-surviving-chord-class} shows that any passing representative satisfying string agreement differs from $T$ by a map of nonnegative degree. It must belong to the class of $T$. This class is fixed by the source rational Bessel group. Phase vanishing for the other classes leaves the constant term of $T$ alone, equal to $A(\gamma h)$. The original integrable absolute sum bounds the sum of the block norms, so the resulting coset sum converges absolutely.\leanfn{See \href{https://github.com/lsernau/grh-lean/blob/72fead102f663d5c7a407964a13dc0cf9d977207/grh/src/Lib/Weil/Periods/ConstantTerm.lean\#L592}{The Fourier coefficient from the sum over string-agreement classes}, \href{https://github.com/lsernau/grh-lean/blob/72fead102f663d5c7a407964a13dc0cf9d977207/grh/src/Lib/Weil/Periods/BesselIntegral.lean\#L680}{Classes satisfying string agreement and passing the test}, \href{https://github.com/lsernau/grh-lean/blob/72fead102f663d5c7a407964a13dc0cf9d977207/grh/src/Lib/Weil/Periods/Unfolding.lean\#L370}{The retained string constant term in each block}, \href{https://github.com/lsernau/grh-lean/blob/72fead102f663d5c7a407964a13dc0cf9d977207/grh/src/Lib/Weil/Periods/Unfolding.lean\#L435}{The absolutely convergent period expansion}.}
\end{proof}

\thmsource{Lean source: \leansourcelink{src/Lib/Weil/Periods/Unfolding.lean}{512}.}
\begin{lemma}[Bessel unfolding and integrability for one step]
\label{lem:a-bessel-unfold-integral}
\emitthmsource
Keep the data of \statementref{Lemma}{lem:a-bessel-orbit-unfold}. If $z$ is a rapidly decreasing automorphic function on $\Ugrp(V')(\mathbb A)$, then $A(h)z(h)$ is absolutely integrable on $R'_{\mathbb Q}\backslash \Ugrp(V')(\mathbb A)$ and
\[
 \int_{R(\mathbb Q)\backslash R} W(r)
   \int_{\Gamma'\backslash \Ugrp(V')(\mathbb A)}
     \theta\bigl(\omega(r,h)\Phi\bigr)z(h)\,dh\,dr
  =
 \int_{R'_{\mathbb Q}\backslash \Ugrp(V')(\mathbb A)}
     A(h)z(h)\,dh .
\]
\end{lemma}
\begin{proof}
Write $a_w(r,h)$ for the weighted term obtained from the decomposition using $d=0$. There exists a constant $C$ such that for every $r,h$,
\[
 \sum_w|a_w(r,h)z(h)|\le C.
\]
First use the rational source Bessel group to move $r$ to a compact covering set of representatives and take a uniform norm bound for $W(r)\omega(r,1)\Phi$. Next use $\Gamma'$ to reduce $h$. The rational action permutes $w$. After reduction, the operator norm of $\omega(1,h)$ is bounded by a power of the target height. Automorphic invariance and rapid decay of $z$ bound the corresponding product of $h$ at the original point. The sum of constant-term norms then gives $C$. Both original quotient measures are finite, so the product integrand is absolutely integrable.\leanfn{See \href{https://github.com/lsernau/grh-lean/blob/72fead102f663d5c7a407964a13dc0cf9d977207/grh/src/Lib/Weil/Periods/Unfolding.lean\#L459}{Combining reduction and rapid decay}, \href{https://github.com/lsernau/grh-lean/blob/72fead102f663d5c7a407964a13dc0cf9d977207/grh/src/Lib/Weil/Periods/Unfolding.lean\#L562}{Integrability on the product space}.}

\statementref{Lemma}{lem:a-bessel-orbit-unfold} and its blockwise proof give
\[
 \sum_{\gamma\in R'_{\mathbb Q}\backslash\Gamma'}
      |A(\gamma h)z(h)|\le C\,\operatorname{vol}\bigl(R(\mathbb Q)\backslash R\bigr).
\]
To prove this bound, take absolute values and expand each block in rational classes. The block index sets are disjoint, so the preceding sum over all classes bounds the result. Thus the absolute coset sum is integrable on $\Gamma'\backslash \Ugrp(V')(\mathbb A)$. Since $A$ is left $R'_{\mathbb Q}$-invariant and $z$ left $\Gamma'$-invariant, the unfolding formula for the quotient integral gives
\[
 \int_{\Gamma'\backslash \Ugrp(V')(\mathbb A)}
      \sum_{\gamma\in R'_{\mathbb Q}\backslash\Gamma'}A(\gamma h)z(h)\,dh
 = \int_{R'_{\mathbb Q}\backslash \Ugrp(V')(\mathbb A)}A(h)z(h)\,dh .
\]
Applying the same formula to absolute values proves absolute integrability on the right. Apply Fubini to the original double integral and replace the inner Bessel period by \statementref{Lemma}{lem:a-bessel-orbit-unfold} to get the identity.\leanfn{See \href{https://github.com/lsernau/grh-lean/blob/72fead102f663d5c7a407964a13dc0cf9d977207/grh/src/Lib/Weil/Periods/Unfolding.lean\#L483}{A bound for the absolute coset sum}, \href{https://github.com/lsernau/grh-lean/blob/72fead102f663d5c7a407964a13dc0cf9d977207/grh/src/Lib/Weil/Periods/Unfolding.lean\#L518}{Quotient unfolding and Fubini}.}
\end{proof}

\thmsource{Lean source: \leansourcelink{src/New/UpperBound/Seesaw/UnfoldingStages.lean}{678}.}
\begin{lemma}[Single-step unfolding along the lower chain]
\label{lem:a-seesaw-single-stage}
\emitthmsource
Fix a twist $\vartheta$, source form $\phi$, linear functional $t$ on the source Weyl module, and lower-chain points $g_k=(\xi_k,y_k)$. Let $k\le m-1$ and let $\mathcal Y'$ be a grading of string length $(m-2-k)_+$ on $Y_{k+1}$, with $(j)_+=\max(j,0)$. Let $R_k$ be the Bessel group of the standard grading of $Y_k$, and $R'_{\mathbb Q}$ the rational points of the Bessel group for $\mathcal Y'$. Write $W_k$ for the function on $R_k$ entering step $k+1$ and $\mathcal W_{k+1}$ for the result of unfolding with $\mathcal Y'$. Then for every rapidly decreasing automorphic function $z$ on $Y_{k+1}$,
\[
 \int_{R_k(\mathbb Q)\backslash R_k} W_k(r)
 \int_{[\Ugrp(Y_{k+1})]}
  \theta\bigl(\omega_k(rg_k,h)\Phi_k\bigr)z(h)\,dh\,dr
 =
 \int_{R'_{\mathbb Q}\backslash \Ugrp(Y_{k+1})(\mathbb A)}
    \mathcal W_{k+1}(h)z(h)\,dh ,
\]
and the right-hand side is absolutely integrable. Here $\Phi_k$ is the lower-chain kernel datum at step $k+1$ and $\omega_k$ uses that step's fixed weight and twist.
\end{lemma}
\begin{proof}
Put $\Phi_k'=\omega_k(g_k,1)\Phi_k$. The representation law for the Weil action gives
\[
 \omega_k(r,h)\Phi_k'=\omega_k(rg_k,h)\Phi_k .
\]
By definition $W_0$ is the source Bessel function of $t,\phi$. For the next standard grading the recursion is
\[
 W_{k+1}(r')=\mathcal W_{k+1}(r'g_{k+1}),
 \qquad
 \mathcal W_{k+1}(h)=
 \int_{R_k(\mathbb Q)\backslash R_k}
    W_k(r)\operatorname{CT}_{T_k}\bigl(\omega_k(r,h)\Phi_k'\bigr)\,dr ,
\]
Here $T_k$ is the string map between the two stages. The initial $W_0$ satisfies source-grading equivariance, and the generic characters of adjacent standard gradings use the same constant. A target unipotent change of variables in one unfolding sends the string map $T_k$ to $uT_k$. Since $u-1$ raises degree by at least $2$ and $T_k$ lowers it by $1$, the two maps differ by positive degree. The phase of the constant term under this change is determined by $T_kT_k^\dagger$, precisely the shift in the target character. With continuity and left rational invariance, the recursion gives the equivariance of each $W_k$.\leanfn{See \href{https://github.com/lsernau/grh-lean/blob/72fead102f663d5c7a407964a13dc0cf9d977207/grh/src/New/UpperBound/Seesaw/UnfoldingStages.lean\#L589}{Moving $g_k$ into the kernel data}, \href{https://github.com/lsernau/grh-lean/blob/72fead102f663d5c7a407964a13dc0cf9d977207/grh/src/New/UpperBound/Seesaw/UnfoldingStages.lean\#L610}{Recursion for the unfolded functions}, \href{https://github.com/lsernau/grh-lean/blob/72fead102f663d5c7a407964a13dc0cf9d977207/grh/src/Lib/Weil/Periods/ConstantTerm.lean\#L887}{One unfolding step preserves equivariance}, \href{https://github.com/lsernau/grh-lean/blob/72fead102f663d5c7a407964a13dc0cf9d977207/grh/src/New/UpperBound/Seesaw/UnfoldingStages.lean\#L657}{Equivariance at each stage}.}

Apply \statementref{Lemma}{lem:a-bessel-unfold-integral} to $W_k$, $\Phi_k'$, $\mathcal Y'$, and $z$. Its constant-term integral is by definition $\mathcal W_{k+1}$. Reexpress the Weil action using the original data at $(rg_k,h)$ to obtain the identity and absolute integrability.\leanfn{See \href{https://github.com/lsernau/grh-lean/blob/72fead102f663d5c7a407964a13dc0cf9d977207/grh/src/New/UpperBound/Seesaw/UnfoldingStages.lean\#L688}{Unfolding step $k+1$ along the lower chain}.}
\end{proof}

\thmsource{Lean source: \leansourcelink{src/New/UpperBound/Seesaw/UnfoldingStages.lean}{828}.}
\begin{lemma}[The matrix coefficient after unfolding a double step]
\label{lem:a-seesaw-double-coefficient}
\emitthmsource
Fix the functional $t$, the earlier infinite variables $\xi$, and the current terminal infinite variable $v_\infty$. For $0\le i<m/2$ consider the double step formed from the gradings of $Y_{2i},Y_{2i+1},Y_{2i+2}$. The two edge evaluation data determine the embedding of the terminal definite group. Evaluating $\mathcal W_{2i+2}$ at the product of the embedded $k$ and $(v_\infty,y_{2i+2})$ for a terminal point $k$ gives a matrix coefficient of the first $i+1$ upper-chain double lifts on the $X_{2i+2}$ side of the embedding of $k$. The functional defining the coefficient is independent of $\vartheta,\phi,y$.
\end{lemma}
\begin{proof}
Start from $\mathcal W_0$. On the stage-zero Bessel group it is the source function determined by $\phi$ and $t$, completely specified by its matrix coefficient on the definite part $\Ugrp(X_0)$ and equivariance for the unipotent character. Assume the function at even stage $2i$ is still of this form, with definite-side form $\psi$ and functional $\lambda$. We compute the functional after the double step.

Write the next two unfoldings as integrals over the two Bessel groups. Each edge contributes a theta constant term translated along the string map. Write $\Phi_i,\Phi'_i$ for the Schwartz data of the two edges after translation by the outer variables. The product of their constant terms can be computed using the arrow product of $\Phi_i,\Phi'_i$. This datum lies in
\[
 \operatorname{Hom}\bigl(Y_{2i+1},Y_{2i}^{-}\oplus Y_{2i+2}\bigr).
\]
On the first unipotent group the characters of the source function and the first-edge kernel cancel. On the second unipotent group the integral restricts the arrow product to the degree-zero term, obtained by Heisenberg translation along the combined string map followed by the degree-zero projection. The two averages give a constant $z_i$ depending only on the gradings and measures, and a double theta integral on $\Ugrp(X_{2i})$ and $\Ugrp(X_{2i+1})$ with that datum as kernel, at the definite embedding of the terminal variable in $k$.\leanfn{See \href{https://github.com/lsernau/grh-lean/blob/72fead102f663d5c7a407964a13dc0cf9d977207/grh/src/New/UpperBound/Seesaw/UnfoldingStages.lean\#L428}{Constant terms and unipotent averages for a double step}.}

The degree-zero datum factors as an infinite Fock vector times the evaluated finite tensor. First separate the infinite and finite parts of the two Schwartz data. Weil translation, the arrow product, Heisenberg translation, and the degree-zero projection all preserve the product structure. The infinite Weil action depends only on the character weights, fixed for admissible twists. Thus the infinite vector $u_i$ depends only on the fixed string maps and infinite variables, independently of $\vartheta,y$. The finite factor is obtained by double-step evaluation of the two translated data. Absorb $z_i$ into the fixed infinite vector and use $z_i u_i$ for the infinite datum.\leanfn{See \href{https://github.com/lsernau/grh-lean/blob/72fead102f663d5c7a407964a13dc0cf9d977207/grh/src/New/UpperBound/Seesaw/UnfoldingStages.lean\#L705}{Factoring the real and finite data}.}

For any fixed infinite Fock vector, the double theta integral on the definite spaces gives a common terminal functional. First use a product of edge polynomials for the infinite vector and a pure tensor of edge data for the finite part. The theta kernel of the arrow product is the product of the edge theta kernels times a fixed constant. Apply the matrix coefficient formula for algebraic theta lifting to $\Ugrp(X_{2i})$ with the input functional $\lambda$ to obtain the intermediate functional. The formula for $\Ugrp(X_{2i+1})$ then gives the terminal functional, so the double integral agrees with the matrix coefficient of the two successive lifts. Linearity extends this to any finite tensor, and products of the two sets of edge polynomials span all polynomial vectors. Finally the coefficient families from terminal functionals form a finite-dimensional linear space, closed in the pointwise topology indexed by twist, finite data, input form, and terminal point. The double integrals depend continuously on the infinite vector and polynomial vectors are dense in Fock space. The formula extends to all infinite vectors with a terminal functional valid for the entire family.\leanfn{See \href{https://github.com/lsernau/grh-lean/blob/72fead102f663d5c7a407964a13dc0cf9d977207/grh/src/New/UpperBound/Seesaw/UnfoldingStages.lean\#L116}{The definite double theta integral and closedness of the coefficient space}.}

Take $z_i u_i$ as infinite vector. The argument above turns the double integral into a matrix coefficient of the double lift with the upper-chain evaluation tensor. Equivariance for the unipotent character is preserved on the terminal Bessel group, so this coefficient determines the source function there. The induction starts with $\lambda=t$ and $\psi=\phi$ and supplies a common functional at every even stage. Evaluate on the indicated terminal embedding.\leanfn{See \href{https://github.com/lsernau/grh-lean/blob/72fead102f663d5c7a407964a13dc0cf9d977207/grh/src/New/UpperBound/Seesaw/UnfoldingStages.lean\#L749}{Double-step induction and evaluation at the terminal stage}.}
\end{proof}

\thmsource{Lean source: \leansourcelink{src/New/UpperBound/Seesaw/EdgePairEvaluation.lean}{199}.}
\begin{lemma}[Lattice support of the finite degree-zero term]
\label{lem:a-finite-zero-term-support}
\emitthmsource
In the finite Weil model, let $H(z)$ be the Heisenberg operator for displacement $z\in\mathbb A_{E,f}^{n}$. Let $C$ be the isotropic frame of a Witt decomposition in standard lattice coordinates, and $Z$ the linear map taking the finite degree-zero term along it. Write $\widehat{\mathcal O}$ for the ring of finite adelic integers and $^\ast$ for conjugate transpose. Suppose $M,M'$ satisfies $MM'^\ast=1$ and $H(Mx)F=F$ for every $x\in\widehat{\mathcal O}^{n}$. If $Z(H(y)F)\ne0$, there exists $l\in\widehat{\mathcal O}^{n}$ such that
\[
 C^\ast(y-M'l)=0.
\]
\end{lemma}
\begin{proof}
Since $C$ is a frame of the decomposition and $MM'^\ast=1$, the matrix $A=M'^\ast C$ has a finite adelic left inverse. We can further choose an invertible square matrix $B$ so that $AB$ is integral and has integral left inverse $L$. At each finite place $w$ the module $\operatorname{im}(A_w)\cap\mathcal O_w^n$ is saturated, so over the discrete valuation ring it is a direct summand. A basis for the intersection lattice is $A_wB_w$ and gives an integral left inverse $L_w$. At almost all places $A_w$ and its original left inverse are integral, so take $B_w=1$. Assembling the local matrices gives $B,L$.\leanfn{See \href{https://github.com/lsernau/grh-lean/blob/72fead102f663d5c7a407964a13dc0cf9d977207/grh/src/New/UpperBound/Seesaw/EdgePairEvaluation.lean\#L204}{Finite support from an integral frame}.}

For an integral vector $a$ put $z=CBa$. By $z=M(ABa)$ we have $H(z)F=F$. Since $z$ is in the span of the isotropic frame, the averaging definition of the degree-zero term gives $Z(H(z)F')=Z(F')$ for every datum $F'$. The Heisenberg commutation relation yields
\[
 \psi_{\mathrm{sd}}(z^\ast y)\,Z(H(y)F)=Z(H(y)F),
\]
Here $\psi_{\mathrm{sd}}$ is the additive character for which the standard lattice is self-dual. By nonvanishing it always takes value $1$, so the pairing of $v=B^\ast C^\ast y$ with each integral $a$ lies in its kernel. Self-duality gives integrality of $v$.\leanfn{See \href{https://github.com/lsernau/grh-lean/blob/72fead102f663d5c7a407964a13dc0cf9d977207/grh/src/New/UpperBound/Seesaw/EdgePairEvaluation.lean\#L186}{Isotropic translation and the degree-zero term}, \href{https://github.com/lsernau/grh-lean/blob/72fead102f663d5c7a407964a13dc0cf9d977207/grh/src/New/UpperBound/Seesaw/EdgePairEvaluation.lean\#L217}{The commutation relation and the dual lattice condition}.}

Take an integral vector $l=L^\ast v$. Conjugate transposition of $LAB=1$ gives $B^\ast C^\ast M' L^\ast=1$ and hence $B^\ast C^\ast(y-M'l)=0$. Cancel the invertible matrix $B^\ast$.
\end{proof}

\thmsource{Lean source: \leansourcelink{src/New/UpperBound/Seesaw/Frequency.lean}{222}.}
\begin{lemma}[A lattice map from a nonzero string constant term]
\label{lem:a-frequency-negative-degree}
\emitthmsource
Keep the two gradings $\mathcal Y,\mathcal Y'$, with string map $T$ and finite part $T_f$. Let $\alpha$ be the finite idele defining the self-dual pairing of the standard lattice. Suppose $\Phi$ factors into an infinite Fock vector and finite datum $F$, with $F$ invariant under the Heisenberg translations given by $b\widehat{\mathcal O}^{\,\dim\operatorname{Hom}(V,V')}$ for a finite idele $b$. If the constant term at $T$ of the twisted Weil translate of $\Phi$ at the adelic point $(x,x')$ is nonzero, there is a finite adelic matrix $W$ such that
\[
  W\Lambda(x_f)\subseteq\Lambda'(x_f'),
  \qquad
  W-\frac{\bar b}{\alpha}T_f\ \text{has degree at least }0 ,
\]
where lattices and matrices are written in rational coordinates of $V,V'$.
\end{lemma}
\begin{proof}
In rational coordinates of $\operatorname{Hom}(V,V')$ take the isotropic frame $C_0$ of the positive-degree coordinates. Source and target reversal both reverse the weights, so Gram matrix $J$ pairs each degree with its negative. Thus
\[
 P_{<0}=J^{-1}C_0C_0^\ast J
\]
is the diagonal projection to negative-degree coordinates. Taking the degree-zero term along $C_0$ and evaluating theta on the complement gives the constant term along the positive-degree coordinates. Factor the Weil translation and the Heisenberg translation by $T$ into infinite and finite parts. The degree-zero term on a pure tensor also factors, and the twist contributes just a scalar. Thus nonvanishing of the original constant term implies that the finite degree-zero term of the moved datum translated along $T_f$ is nonzero.\leanfn{See \href{https://github.com/lsernau/grh-lean/blob/72fead102f663d5c7a407964a13dc0cf9d977207/grh/src/Lib/Weil/Periods/ConstantTerm.lean\#L403}{The graded Witt decomposition and the constant term}, \href{https://github.com/lsernau/grh-lean/blob/72fead102f663d5c7a407964a13dc0cf9d977207/grh/src/New/UpperBound/Seesaw/Frequency.lean\#L172}{Nonvanishing of the finite degree-zero term}.}

Let $g$ be the unitary matrix induced by $(x_f,x_f')$ in the lattice coordinates of $\operatorname{Hom}(V,V')$ and $F_g$ the finite Weil translate of $F$. Commutation of the Weil and Heisenberg actions makes $F_g$ invariant under translations in $bg\widehat{\mathcal O}^{\,\dim\operatorname{Hom}(V,V')}$. Let $z$ be the finite displacement coordinate of $T$ in this model. In \statementref{Lemma}{lem:a-finite-zero-term-support} use $M=bg$ and $M'=\bar b^{-1}g$, satisfying $MM'^\ast=1$. Since the finite degree-zero term is nonzero, there is a vectorization $l$ of an integral matrix such that
\[
 C^\ast\bigl(z-\bar b^{-1} g l\bigr)=0,
\]
where $C$ is the frame of $C_0$ in finite Weil lattice coordinates.

Let $A_x,A_{x'}$ be the coordinate matrices from the chosen standard lattices to $\Lambda(x_f),\Lambda'(x_f')$ at the two stages. Let $L_0$ be the integral matrix corresponding to $l$, and put $W=A_{x'}L_0A_x^{-1}$. Then $W\Lambda(x_f)\subseteq\Lambda'(x_f')$. If $A_H$ changes $\operatorname{Hom}(V,V')$ from finite Weil lattice coordinates to rational coordinates, then
\[
 A_H\bigl(z-\bar b^{-1}g l\bigr)
     =\operatorname{vec}\bigl(\alpha^{-1}T_f-\bar b^{-1}W\bigr).
\]
The displacement coordinate of $T$ includes the self-dual scaling $\alpha^{-1}$, and the unitary action on Hom uses left multiplication by the target matrix and right multiplication by the source inverse. Also $C^\ast=C_0^\ast J A_H$. Multiplying the frame identity on the left by $J^{-1}C_0$ gives
\[
 P_{<0}\operatorname{vec}
     \bigl(\alpha^{-1}T_f-\bar b^{-1}W\bigr)=0.
\]
This says every negative-degree entry of $W-(\bar b/\alpha)T_f$ vanishes.\leanfn{See \href{https://github.com/lsernau/grh-lean/blob/72fead102f663d5c7a407964a13dc0cf9d977207/grh/src/New/UpperBound/Seesaw/Frequency.lean\#L55}{Conjugate transpose of the frame}, \href{https://github.com/lsernau/grh-lean/blob/72fead102f663d5c7a407964a13dc0cf9d977207/grh/src/New/UpperBound/Seesaw/Frequency.lean\#L207}{Change of basis for the finite displacement coordinate}, \href{https://github.com/lsernau/grh-lean/blob/72fead102f663d5c7a407964a13dc0cf9d977207/grh/src/New/UpperBound/Seesaw/Frequency.lean\#L229}{The lattice map in negative degree}.}
\end{proof}

\thmsource{Lean source: \leansourcelink{src/New/UpperBound/Seesaw/Frequency.lean}{289}.}
\begin{lemma}[String lattice inclusion for graded source and target lattices]
\label{lem:a-frequency-graded-lattices}
\emitthmsource
Under \statementref{Lemma}{lem:a-frequency-negative-degree} take finite Bessel points $r,r'$ such that $\Lambda(rx_f)$ and $\Lambda'(r'x_f')$ are graded lattices, meaning they contain every homogeneous component of their vectors. Write $\Lambda(\,\cdot\,)^+$ for lattice projection to the source string and $\Lambda'(\,\cdot\,)^-$ for intersection with the target string. Then
\[
  \frac{\bar b}{\alpha}\,
     T_f\bigl(\Lambda(rx_f)^+\bigr)
  \ \subseteq\ \Lambda'(r'x_f')^- .
\]
\end{lemma}
\begin{proof}
Let $R,R'$ be the matrix of $r,r'$ in rational coordinates. Put $c=\bar b/\alpha$ and take $W$ from \statementref{Lemma}{lem:a-frequency-negative-degree}. Then
\[
 W_0=R'W R^{-1}
\]
sends $\Lambda(rx_f)$ into $\Lambda'(r'x_f')$. The Bessel matrices and their inverses do not lower degree, and
\[
 R'T_f R^{-1}-T_f
   =(R'-1)T_f R^{-1}+T_f(R^{-1}-1)
\]
has nonnegative degree. Indeed, $T_f$ changes degree $d$ to $d-1$ between corresponding string coordinates only. On the relevant entries the Bessel conditions make the columns of $R'-1$ and rows of $R^{-1}-1$ strictly raise degree. Thus both terms have nonnegative degree. From
\[
 W_0-cT_f=R'(W-cT_f)R^{-1}+c(R'T_fR^{-1}-T_f)
\]
we see $W_0-cT_f$ also has nonnegative degree. Since $T_f$ is homogeneous of degree $-1$, the degree-$-1$ component of $W_0$ equals $cT_f$.\leanfn{See \href{https://github.com/lsernau/grh-lean/blob/72fead102f663d5c7a407964a13dc0cf9d977207/grh/src/Lib/Weil/Periods/Grading.lean\#L817}{Bessel translations preserve the negative-degree part of the string map}, \href{https://github.com/lsernau/grh-lean/blob/72fead102f663d5c7a407964a13dc0cf9d977207/grh/src/New/UpperBound/Seesaw/Frequency.lean\#L136}{Matrix comparison under the two Bessel translations}.}

For $\lambda\in\Lambda(rx_f)$ each source homogeneous component $\operatorname{pr}_d\lambda$ is still in the source lattice. Mapping it by $W_0$ into the target lattice and taking $\operatorname{pr}'_{d-1}$ keeps it in that lattice. Sum over $d$ to obtain
\[
 cT_f\lambda=\sum_d\operatorname{pr}'_{d-1}
                         W_0\operatorname{pr}_d\lambda
            \in\Lambda'(r'x_f').
\]
Finally $T_f$ kills the definite source block and takes values in the target string space. Thus the image of the source lattice by $cT_f$ is the image of its projection to the string, and it lies in the intersection of the target lattice with the string space.\leanfn{See \href{https://github.com/lsernau/grh-lean/blob/72fead102f663d5c7a407964a13dc0cf9d977207/grh/src/New/UpperBound/Seesaw/Frequency.lean\#L74}{Homogeneous parts preserve graded lattices}, \href{https://github.com/lsernau/grh-lean/blob/72fead102f663d5c7a407964a13dc0cf9d977207/grh/src/New/UpperBound/Seesaw/Frequency.lean\#L108}{String projection and the string map}, \href{https://github.com/lsernau/grh-lean/blob/72fead102f663d5c7a407964a13dc0cf9d977207/grh/src/New/UpperBound/Seesaw/Frequency.lean\#L307}{String lattice inclusion from a nonzero constant term}.}
\end{proof}

\thmsource{Lean source: \leansourcelink{src/New/UpperBound/Seesaw/UnfoldingStages.lean}{1093}.}
\begin{lemma}[Lattice support of the unfolded functions]
\label{lem:a-seesaw-lattice-support}
\emitthmsource
Suppose $\phi\in U$ and the first finite variable $y_0$ belongs to the chosen principal level. At intermediate stage $1\le j<m$, suppose the unfolded function $\mathcal W_j$ is nonzero at some infinite point and finite point $x$. Then every finite Bessel translation $r'$ sending $x$ to a graded lattice satisfies
\[
 \bar{\mathfrak q}^{\,j}\Lambda_j^{-}\ \subseteq\ \Lambda_j'{}^{-},
 \qquad
 {\mathfrak q}^{\,j}\Lambda_j'{}^{-}\ \subseteq\ \Lambda_j^{-},
\]
Here $\mathfrak q$ is the finite idele generating the twist modulus. The lattices $\Lambda_j,\Lambda'_j$ correspond to the identity and $r'x$, respectively, and the superscript $-$ denotes intersection with the string space.
\end{lemma}
\begin{proof}
Induct on $j=k+1$. The function $\mathcal W_{k+1}$ integrates the preceding Bessel function times the single-edge theta constant term. A nonzero integral supplies a point where both factors are nonzero. Write $x_k$ for its finite variable at the preceding stage. If $k=0$, then $x_0$ is a Bessel translate of $y_0$ and can be returned to the standard graded lattice. If $k>0$, the nonzero value of the preceding function comes from $\mathcal W_k$. The induction hypothesis and existence of graded representatives give a graded lattice $\Lambda'_k$ for $x_k$ satisfying $\mathfrak q^{\,k}\Lambda_k'{}^{-}\subseteq\Lambda_k^{-}$.\leanfn{See \href{https://github.com/lsernau/grh-lean/blob/72fead102f663d5c7a407964a13dc0cf9d977207/grh/src/New/UpperBound/Seesaw/UnfoldingStages.lean\#L1092}{Choosing a lattice at the previous stage from a nonzero integral}.}

Let $\Lambda_k^+$ and $\Lambda_k'{}^+$ be the projections of the two lattices to the string spaces. Lattice pairing duality turns the preceding inclusion into $\bar{\mathfrak q}^{\,k}\Lambda_k^+\subseteq\Lambda_k'{}^+$. Write $\mathrm{sh}_k$ for the edge string map and $\mathfrak a$ for the fixed finite idele of the data. The nonzero constant term condition of \statementref{Lemma}{lem:a-frequency-graded-lattices} gives
\[
 \frac{\bar{\mathfrak q}}{\mathfrak a}\,
       \mathrm{sh}_k(\Lambda_k'{}^+)\ \subseteq\ \Lambda'_{k+1}{}^{-}.
\]
The standard lattice also satisfies $\mathrm{sh}_k(\Lambda_k^+)=\mathfrak a\Lambda_{k+1}^{-}$. Substitute the dual inclusion in the frequency condition to obtain
\[
 \bar{\mathfrak q}^{\,k+1}\Lambda_{k+1}^{-}
   =\frac{\bar{\mathfrak q}}{\mathfrak a}\,
           \mathrm{sh}_k(\bar{\mathfrak q}^{\,k}\Lambda_k^+)
   \ \subseteq\ \Lambda'_{k+1}{}^{-}.
\]
The argument works for every graded lattice in the terminal orbit.\leanfn{See \href{https://github.com/lsernau/grh-lean/blob/72fead102f663d5c7a407964a13dc0cf9d977207/grh/src/New/UpperBound/Seesaw/Sandwich.lean\#L136}{Dual lattices and transfer of the lower bound}.}

It remains to prove the reverse inclusion. Use the lower bound just proved in the comparison lemma for standard string lattices. There are two possibilities. The first gives $\mathfrak q^{\,k+1}\Lambda'_{k+1}{}^{-}\subseteq\Lambda_{k+1}^{-}$. The second is possible only when $k+1$ is even and the terminal lattice splits along the definite part at some $w\mid\pi_N$. The twist modulus is coprime to $N$, so $\mathfrak q$ is a unit at $w$ and its conjugate. In the second case the splitting makes the scalar operator $z_w$ fix the terminal finite evaluation data. Equivariance along double lifts transfers this invariance to the source and gives the average over its $z_w$ powers. Ramification of the source form makes that average zero. Thus the even-stage unfolded function vanishes on the Bessel orbit, contradicting nonvanishing. Only the first possibility remains.\leanfn{See \href{https://github.com/lsernau/grh-lean/blob/72fead102f663d5c7a407964a13dc0cf9d977207/grh/src/New/UpperBound/Seesaw/UnfoldingStages.lean\#L1053}{Vanishing in the split case by ramification}, \href{https://github.com/lsernau/grh-lean/blob/72fead102f663d5c7a407964a13dc0cf9d977207/grh/src/New/UpperBound/Seesaw/UnfoldingStages.lean\#L1141}{The upper lattice bound and exclusion of the exceptional case}.}
\end{proof}

\thmsource{Lean source: \leansourcelink{src/New/UpperBound/Seesaw/SeesawIdentity.lean}{535}.}
\begin{lemma}[Seesaw unfolding and the section tensor]
\label{lem:a-seesaw-unfolding}
\emitthmsource
Fix a complex linear functional $t$ on the Weyl module of $X_0$ and an archimedean adelic point $g_\infty$ of $Y_0$. Write $\mathcal P=\mathcal P_{g_\infty,t}$ for the corresponding Poincaré map. For tame character $\vartheta$ and lower-chain finite variables $y=(y_j)_{j=0}^m$, write $E_\vartheta(\bar y)$ for the pure tensor of the double-step finite data evaluated at $y$, where $\bar y_j$ is the class of $y_j$ modulo its level. Let $\Omega_j$ be the finite section classes at step $j$ and $\mu_j(\omega)$ their covering multiplicity. The terminal finite variable is the identity. Define
\[
 F_\vartheta
 =\sum_{\omega\in\Omega_0\times\cdots\times\Omega_{m-1}}
  \left(\prod_{j=0}^{m-1}\mu_j(\omega_j)^{-1}\right)E_\vartheta(\omega).
\]
There exists a constant $c_0\in\mathbb C$ depending only on $D,t,g_\infty$ such that for every $\vartheta$ and $\phi\in U$,
\[
 T_\vartheta(\mathcal P\phi)
   = c_0\,\operatorname{Res}_{X_m\to Y_m}
            \mathcal T_\vartheta(\phi;F_\vartheta),
\]
where $T_\vartheta$ is the full lower-chain lift and $\mathcal T_\vartheta(\phi;F)$ is the form on $X_m$ given by $m$ lifts of $\phi$ using upper-chain finite tensor $F$.
\end{lemma}
\begin{proof}
Let $G_j=\Ugrp(Y_j)$ be the lower-chain group at stage $j$, with Bessel subgroup $R_j$. Write $J_j$ for the level on $G_j(\mathbb A_f)$ with twist $\vartheta$, $J_j^\circ$ for the standard lattice stabilizer, $J_M$ for the chosen principal level at stage zero, and $[\cdot]$ for the corresponding rational left quotient. Write $L_j$ for the lower lift on edge $j$, including factor $[J_j^\circ:J_j]$. Interchanging the integration variables of its two-sided theta kernel gives the transpose lift $L_j^{\mathrm t}$ with the same factor. For a stage-$j$ form $v$ and a conjugate-type stage-$j+1$ form $z$,
\[
 \int_{[G_{j+1}]} z\,L_jv
   =\int_{[G_j]}(L_j^{\mathrm t}z)\,v .
\]
The arrow property of the lift and transpose lift for Schwartz data gives absolute convergence and permits exchanging the integrals. As the theta data for $L_j^{\mathrm t}$ are fixed by $J_j$, $L_j^{\mathrm t}z$ is invariant under $J_j$.\leanfn{See \href{https://github.com/lsernau/grh-lean/blob/72fead102f663d5c7a407964a13dc0cf9d977207/grh/src/New/UpperBound/Seesaw/SeesawIdentity.lean\#L168}{Transpose lifts and level invariance}.}

\emph{Recursion for the finite section sums.}
Write $T_\vartheta^{[j]}$ for the first $j$ lifts of the lower chain, and $\xi=(\xi_i)$ and $y=(y_i)$ for the archimedean and finite coordinates of the outer variables. Let $\mathcal W_j(\xi,y;h)$ be the function remaining on $G_j$ after $j$ unfoldings and put
\[
 \Xi_j(z;\xi,y)=\int_{R_j(\mathbb Q)\backslash G_j(\mathbb A)}
                    \mathcal W_j(\xi,y;h)z(h)\,dh ,
\]
where $R_j(\mathbb Q)$ is the corresponding rational Bessel subgroup. Single-step unfolding expresses the Bessel period of the stage-$j$ remaining function and the transpose lift as a period at the next stage. Fix $(\xi,y)$ and put
$H_j(h)=\mathcal W_j(\xi,y;h)(L_j^{\mathrm t}z)(h)$. Updating outer variable $j$ to $h=(\alpha,v)$ gives
\[
 \int_{[R_j]}H_j(rh)\,dr
    =[J_j^\circ:J_j]\,
      \Xi_{j+1}\bigl(z;\xi[j\leftarrow\alpha],y[j\leftarrow v]\bigr).
\]
\statementref{Lemma}{lem:a-seesaw-single-stage} also gives integrability of the product of $\mathcal W_j$ and the conjugate-type form on the required quotient.\leanfn{See \href{https://github.com/lsernau/grh-lean/blob/72fead102f663d5c7a407964a13dc0cf9d977207/grh/src/New/UpperBound/Seesaw/SeesawIdentity.lean\#L267}{The single-step period recursion}.}

Suppose at an endpoint $\Xi_j$ factors as $a(\xi)S(\bar y)$, with the factorization needed only when $y_0\in J_M$ and $y_m=1$. Induct on $j$ to show that $\int_{[G_j]}z\,T_\vartheta^{[j]}(\mathcal P\phi)$ equals a constant independent of $\vartheta,\phi,z,S$ times
\[
 \sum_{\omega\in\Omega_0\times\cdots\times\Omega_{j-1}}
   \left(\prod_{i=0}^{j-1}\mu_i(\omega_i)^{-1}\right)S(\omega),
\]
where the coordinates not summed over are the identity classes.

For $j=1$ first transpose $L_0$ onto $z$. Expanding the Poincaré map averages the Bessel period over representatives of $J_M/J_0$. This quotient corresponds to $\Omega_0$ and $J_0$-invariance of the transpose lift makes the average independent of the representatives. Unfold the period for each representative. The numerator has factor $[J_0^\circ:J_0]$ and the averaging denominator is $[J_M:J_0]$, leaving $[J_0^\circ:J_M]$. Since $\mu_0=1$ and the first archimedean variable is fixed at $g_\infty$, the constant is $[J_0^\circ:J_M]a(\xi^{(0)})$, with first coordinate $g_\infty$ of $\xi^{(0)}$ and all others the identity.\leanfn{See \href{https://github.com/lsernau/grh-lean/blob/72fead102f663d5c7a407964a13dc0cf9d977207/grh/src/New/UpperBound/Seesaw/SeesawIdentity.lean\#L327}{The first period and level indices}.}

In passing from stage $j+1$ back to stage $j$, first transpose $L_j$ onto $z$ and substitute the assumed factorization in the Bessel period of $H_j$. By $1\le j<m$, updating coordinate $j$ preserves $y_0\in J_M$ and $y_m=1$. \statementref{Lemma}{lem:a-seesaw-cut-integral} turns $\Xi_j(L_j^{\mathrm t}z;\xi,y)$ into
\[
 a'(\xi)\sum_{\omega\in\Omega_j}
    \mu_j(\omega)^{-1} S(\bar y[j\leftarrow\omega]).
\]
The finite part of this formula cancels the factor $[J_j^\circ:J_j]$ in the stage-$j$ period. Here $a'(\xi)$ is the integral over $\alpha$ of $a(\xi[j\leftarrow\alpha])$ for a fixed archimedean measure depending only on $R_j$ and standard level $J_j^\circ$, so $a'$ is independent of $\vartheta$. The hypotheses follow from continuity and rational left Bessel invariance of the unfolded functions, the integrability just proved, and the lattice support bound of \statementref{Lemma}{lem:a-seesaw-lattice-support}. That bound ensures that $\Omega_j$ covers the orbit of the finite variables whenever $\mathcal W_j$ is nonzero. Apply the stage-$j$ induction hypothesis to the new factorization and combine the new step-$j$ sum with the previous $j$-fold finite sum. This gives stage $j+1$.\leanfn{See \href{https://github.com/lsernau/grh-lean/blob/72fead102f663d5c7a407964a13dc0cf9d977207/grh/src/New/UpperBound/Seesaw/IntegralAsCosetSum.lean\#L352}{The section integral formula for the Bessel period}, \href{https://github.com/lsernau/grh-lean/blob/72fead102f663d5c7a407964a13dc0cf9d977207/grh/src/New/UpperBound/Seesaw/SeesawIdentity.lean\#L413}{Induction on the section sum and orbit coverage}.}

\emph{Factoring the terminal period.}
Take $j=m$ and put $\psi=\mathcal T_\vartheta(\phi;E_\vartheta(\bar y))$. The terminal space $Y_m$ is definite and embeds in $X_m$ by $k\mapsto\operatorname{diag}(k,1)$. When $y_m=1$, \statementref{Lemma}{lem:a-seesaw-double-coefficient} gives a functional $\lambda_\xi$ such that $\mathcal W_m(\xi,y;k)$ is the matrix coefficient of $\psi$ for $\lambda_\xi$ at this embedding. The functional $\lambda_\xi$ depends only on $t,\xi,D$ and is independent of $\vartheta,\phi,y$.\leanfn{See \href{https://github.com/lsernau/grh-lean/blob/72fead102f663d5c7a407964a13dc0cf9d977207/grh/src/New/UpperBound/Seesaw/UnfoldingStages.lean\#L828}{The matrix coefficient formula for double-step unfolding}.}

For a form $v$ on $Y_m$ take $z=\overline v$. Write $A(v,\psi)$ for the class integral over finite points whose integrand is the Gram pairing at $h$ of the algebraic value of $v$ with the Weyl vector obtained by branching from $\psi(\operatorname{diag}(h,1))$. The archimedean matrix coefficient integral gives $\Xi_m(\overline v;\xi,y)=C(\lambda_\xi)A(v,\psi)$, with coefficient $C(\lambda_\xi)$ depending only on $\lambda_\xi$ and the fixed weight and measures. Also let $d_m$ be the Weyl dimension of $Y_m$, and let $c_N=\prod_{\tau\in\Phi}\tau(\overline{\pi_N})^{|\mu_{X_m,\tau}|-|\mu_{Y_m,\tau}|}$ with $|\mu_{X_m,\tau}|$ and $|\mu_{Y_m,\tau}|$ the box counts of the source and target Young diagrams, respectively. Restriction includes the constant $c_N^{-1}$ and the target inner product is the class integral divided by $d_m$. Hence
\[
 \big\langle v,\operatorname{Res}_{X_m\to Y_m}\psi\big\rangle
  =d_m^{-1}c_N^{-1}A(v,\psi),
\]
Here $d_m,c_N$ are nonzero. Combining the two identities shows that $\Xi_m(\overline v;\xi,y)$ has the required factorization. Take\leanfn{See \href{https://github.com/lsernau/grh-lean/blob/72fead102f663d5c7a407964a13dc0cf9d977207/grh/src/New/UpperBound/Seesaw/SeesawIdentity.lean\#L82}{The terminal matrix coefficient and the inner product after restriction}.}
\[
 a(\xi)=C(\lambda_\xi)d_mc_N,\qquad
 S(\bar y)=
 \big\langle v,\operatorname{Res}_{X_m\to Y_m}
                 \mathcal T_\vartheta(\phi;E_\vartheta(\bar y))\big\rangle .
\]

Substitute in the section-sum formula from the induction. The inner product is linear in the second argument, and lifting and restriction are linear in the data. Thus the weighted sum moves inside the tensor argument, giving $F_\vartheta$. For every $v$,
\[
 \big\langle v,T_\vartheta(\mathcal P\phi)\big\rangle
 =c_0\big\langle v,\operatorname{Res}_{X_m\to Y_m}
                          \mathcal T_\vartheta(\phi;F_\vartheta)\big\rangle.
\]
The constant $c_0$ from the induction depends only on $a,D,g_\infty$, and the terminal $a$ depends only on $t,D$ and the outer archimedean variables. Thus $c_0$ is common to $\vartheta$ and $\phi$. Nondegeneracy of the inner product gives the identity of forms.\leanfn{See \href{https://github.com/lsernau/grh-lean/blob/72fead102f663d5c7a407964a13dc0cf9d977207/grh/src/New/UpperBound/Seesaw/SeesawIdentity.lean\#L535}{The section tensor formula for the full lower chain}.}
\end{proof}

\thmsource{Lean source: \leansourcelink{src/New/UpperBound/Seesaw/SeesawIdentity.lean}{624}. See also \href{https://github.com/lsernau/grh-lean/blob/72fead102f663d5c7a407964a13dc0cf9d977207/grh/src/New/UpperBound/Seesaw/SeesawIdentity.lean\#L592}{An admissible expansion of the section tensor}.}
\begin{theorem}[The seesaw identity]
\label{thm:a-seesaw}
\emitthmsource
The \emph{seesaw identity} here compares the following lifts. Fix an archimedean adelic point $g_\infty$ of $Y_0$ and a complex functional $t$ on the Weyl module of $X_0$. Let $\mathcal P=\mathcal P_{g_\infty,t}\colon U\to\operatorname{Forms}(Y_0)$ be the associated Poincaré map. There is a constant $c_{\mathrm{ss}}\in\mathbb C^\times$ independent of twist and source form such that for every twist $\vartheta$ there are chain data $\Psi_\vartheta$ satisfying
\begin{equation}
 T_\vartheta(\mathcal P\phi)
   =c_{\mathrm{ss}}\operatorname{Comp}_{\vartheta,\Psi_\vartheta}(\phi)
 \qquad(\phi\in U),
\label{eq:a-seesaw}
\end{equation}
where $T_\vartheta$ is the lift along lower chain $Y_0\to\cdots\to Y_m$ and $\operatorname{Comp}_{\vartheta,\Psi_\vartheta}$ lifts along upper chain $X_0\to\cdots\to X_m$ using glued data and then restricts to $Y_m$.
\end{theorem}
\begin{proof}
Let $F_\vartheta$ be the finite tensor of \statementref{Lemma}{lem:a-seesaw-unfolding}. The lemma also gives a constant $c_0\in\mathbb C$ such that $T_\vartheta(\mathcal P\phi)$ equals the upper-chain lift for $c_0F_\vartheta$ followed by restriction to $Y_m$. We construct chain data expressing a nonzero multiple of $c_0F_\vartheta$ independent of $\vartheta$.

Let $Q_\vartheta$ be the support of the finite twist modulus and $q_\vartheta=\operatorname{Nm}(Q_\vartheta)$ the product of the norms of those places. Put $R=M^{f n_m^2}$ and $\Delta=R^{2m}$. Write the section sum for the finite tensor as
\[
 F_\vartheta=\sum_{\omega\in\Omega_0\times\cdots\times\Omega_{m-1}}
             \left(\prod_{j=0}^{m-1}\mu_j(\omega_j)^{-1}\right)
             E_\vartheta(\omega),
\]
Here $\Omega_j$ are the section classes for lower-chain finite variables, $\mu_j$ their covering multiplicities, and $E_\vartheta(\omega)$ the pure tensor from evaluation on the classes. Representatives are integral away from $q_\vartheta$. Since $\mu_j(\omega_j)$ divides a power of $q_\vartheta$, the rational weight $\prod_j\mu_j(\omega_j)^{-1}$ is integral at every prime outside $q_\vartheta$. If $Q_\vartheta=\varnothing$, then $\#\Omega_j\le R$ and each weight has real absolute value at most $1$, so their sum is at most $R^m$.\leanfn{See \href{https://github.com/lsernau/grh-lean/blob/72fead102f663d5c7a407964a13dc0cf9d977207/grh/src/New/UpperBound/Seesaw/Cut.lean\#L201}{Section multiplicities and class counts}, \href{https://github.com/lsernau/grh-lean/blob/72fead102f663d5c7a407964a13dc0cf9d977207/grh/src/New/UpperBound/Seesaw/Cut.lean\#L321}{Integrality of the weighted section sum}.}

\statementref{Lemma}{lem:a-seesaw-vacuum-expansion} gives a common constant $C\in\mathbb C$ such that for each $\omega$,
\[
 E_\vartheta(\omega)=\sum_x r_{\vartheta,\omega}(x)\,
                         C\,\operatorname{Vac}_{\rm gl}(x),
\]
where $x$ is a tuple of adelic translations at the upper stages and $\operatorname{Vac}_{\rm gl}(x)$ the tensor obtained by translating their integral vacua and then gluing in pairs. Every $x$ in the expansion is admissible, becoming integral away from $q_\vartheta$ after multiplication by $\overline{\pi_M}$. Each $\Delta r_{\vartheta,\omega}(x)$ is integral at the primes away from $q_\vartheta$, and if $Q_\vartheta=\varnothing$, then for every $\sigma\in\operatorname{Aut}(\mathbb C)$,
$\sum_x|\sigma(r_{\vartheta,\omega}(x))|\le1$.
Substituting into the section sum and collecting the coefficients of each $x$ gives
\[
 F_\vartheta=C\sum_x r_\vartheta(x)\operatorname{Vac}_{\rm gl}(x).
\]
The valuation rings are closed under addition and multiplication, so $\Delta r_\vartheta(x)$ is integral at the primes away from $q_\vartheta$. When $Q_\vartheta=\varnothing$, take the sum of absolute values under $\sigma\in\operatorname{Aut}(\mathbb C)$ and use the triangle inequality to get $\sum_x|\sigma(r_\vartheta(x))|\le R^m$.\leanfn{See \href{https://github.com/lsernau/grh-lean/blob/72fead102f663d5c7a407964a13dc0cf9d977207/grh/src/New/UpperBound/Seesaw/GoodDatum.lean\#L526}{The finite pure tensor evaluation theorem}, \href{https://github.com/lsernau/grh-lean/blob/72fead102f663d5c7a407964a13dc0cf9d977207/grh/src/New/UpperBound/Seesaw/SeesawIdentity.lean\#L581}{Finite expansion of the section sum}.}

Choose a common scalar $u\in\mathbb C^\times$: take $u=\Delta/(c_0C)$ if $c_0C\ne0$ and $u=1$ if $c_0C=0$. In the former case the vacuum expansion of $u c_0F_\vartheta$ has coefficients $\Delta r_\vartheta(x)$. In the latter $u c_0F_\vartheta=0$ and use zero coefficients. In both cases the coefficients are integral away from $q_\vartheta$. If $Q_\vartheta=\varnothing$, their sum of absolute values under any conjugate is at most
\[
 \Delta R^m=R^{3m}=M^{3fm n_m^2}.
\]

Set the coordinates of each $x$ to zero at finite places not dividing $M q_\vartheta$, giving $x'$. The element $\overline{\pi_M}$ is a unit at those places, so the removed components are integral and act trivially on the vacuum. The supports of the retained and removed parts are disjoint and stable under complex conjugation, so the Heisenberg product adds no phase. Thus $\operatorname{Vac}_{\rm gl}(x')=\operatorname{Vac}_{\rm gl}(x)$. Add coefficients for the same $x'$ and discard zeros. This gives a finitely supported function $\Psi_\vartheta$ with admissible support at the places dividing $M q_\vartheta$, preserving integrality and the absolute coefficient sum bound. The glued tensor is\leanfn{See \href{https://github.com/lsernau/grh-lean/blob/72fead102f663d5c7a407964a13dc0cf9d977207/grh/src/New/UpperBound/Divisibility/ChainValuation.lean\#L107}{Reducing the support of the translations}.}
\[
 \operatorname{Glue}(\Psi_\vartheta)
 =\sum_{x'}\Psi_\vartheta(x')\operatorname{Vac}_{\rm gl}(x')
 =u c_0F_\vartheta.
\]

Finally check joint invariance. Let $K_j$ be the level of $X_j$ and take $b_j\in K_j$ with $b_m=1$. For $j<m$ let $h_j$ be the finite embedding of $b_j$ in the lower-chain group and put $h_m=1$. Each $h_j$ permutes the section classes $\Omega_j$ and preserves the multiplicities $\mu_j$. By equivariance of finite evaluation, simultaneous left translation of all lower-chain variables by $h_j$ before evaluating the double steps gives the joint Weil action of $b_j$ on the evaluated upper-chain tensor. The terminal embedding in the last double step uses the identity, so $b_m=1$ preserves this equivariance. Changing variables by $\omega_j\mapsto h_j\omega_j$ in the section sum shows that the joint Weil action fixes $F_\vartheta$ and hence $u c_0F_\vartheta$. Thus $\Psi_\vartheta$ satisfies every condition for chain data.\leanfn{See \href{https://github.com/lsernau/grh-lean/blob/72fead102f663d5c7a407964a13dc0cf9d977207/grh/src/New/UpperBound/Seesaw/Cut.lean\#L184}{Translation of sections and preservation of multiplicities}, \href{https://github.com/lsernau/grh-lean/blob/72fead102f663d5c7a407964a13dc0cf9d977207/grh/src/New/UpperBound/Seesaw/SeesawIdentity.lean\#L503}{Joint invariance of the section tensor}.}

Upper-chain lifting and restriction are linear in the finite tensor. Combining \statementref{Lemma}{lem:a-seesaw-unfolding} with $\operatorname{Glue}(\Psi_\vartheta)=u c_0F_\vartheta$ gives
\[
 u\,T_\vartheta(\mathcal P\phi)=\operatorname{Comp}_{\vartheta,\Psi_\vartheta}(\phi).
\]
The identification of automorphic and algebraic forms gives the same identity in both models. Take $c_{\mathrm{ss}}=u^{-1}$.\leanfn{See \href{https://github.com/lsernau/grh-lean/blob/72fead102f663d5c7a407964a13dc0cf9d977207/grh/src/New/UpperBound/Seesaw/SeesawIdentity.lean\#L612}{Common normalization of the chain data}.}
\end{proof}

\section{Norm bounds for short forms and bounds for local lattice indices}
\label{sec:a-conclusions}

Equation~\eqref{eq:a-short} covers all translates and all complex automorphisms. Equation~\eqref{eq:a-index} covers any finite family drawn from finitely many translated images, with the local correction required to work simultaneously for all twists admissible at the given prime.

\thmsource{Lean source: \leansourcelink{src/New/Reduction/ShortAndEnough.lean}{254}. See also \href{https://github.com/lsernau/grh-lean/blob/72fead102f663d5c7a407964a13dc0cf9d977207/grh/src/New/Reduction/ShortAndEnough.lean\#L231}{Rational normalization and conjugate squared norms}.}
\begin{theorem}[A(a): short integral forms span the source space]
\label{thm:short-forms}
\emitthmsource
Put $h=H-mS+L$. The forms
\begin{equation}
  \phi\in I,\qquad
  \|\Theta_1P_g\phi\|_\sigma^2\le e^h
  \quad (g\in G,\ \sigma\in\operatorname{Aut}(\mathbb C))
\label{eq:a-short}
\end{equation}
$\phi$ satisfying this bound span $U$.
\end{theorem}
\begin{proof}
Let $U_0$ be the full source space and $N$ the level prime. A source form is $M$-short if it is integral at every $p$-adic place and its conjugate value at every finite adelic point under every field automorphism of the complex numbers has squared Weyl norm at most $M$. Continue to use $\|\cdot\|_{\mathrm{Weyl}}$ for the Weyl norm on the value space.

\emph{Choosing short integral forms spanning the ramified source space.}
By \statementref{Lemma}{lem:a-short-source-basis}, $U_0$ is spanned by $e^B/4^f$-short forms.

Let $z_w$ be the finite-place scalar translation for each $w\mid\pi_N$. The proof of \statementref{Lemma}{lem:a-ramified-dimension} shows that the ramification operator $\prod_{w\mid\pi_N}(1-z_w)$ has image $U$. Each $z_w$ comes from the standard integral level, so right translation preserves integrality. At any conjugate and any point, a squared norm bound $M$ for the form gives bound $(\sqrt M+\sqrt M)^2=4M$ for the difference from its translate. After all $f$ factors the forms remain integral and their squared norms are at most $4^f(e^B/4^f)=e^B$. A linear map sends a spanning family to a spanning family of its image. If $\mathcal S$ is the family of images, then $\operatorname{span}_{\mathbb C}\mathcal S=U$ and $\mathcal S\subseteq I$.\leanfn{See \href{https://github.com/lsernau/grh-lean/blob/72fead102f663d5c7a407964a13dc0cf9d977207/grh/src/New/UpperBound/Size/ShortColumns.lean\#L962}{Ramification and spanning by short forms}.}

\emph{Estimating the lifts before normalization.}
Write $\mathcal P$ for the Poincaré map, $c_{\mathrm{ss}}\ne0$ for the seesaw constant, and $e_0$ for the projection to the target. For $\phi\in\mathcal S$ put $z=\mathcal P(c_{\mathrm{ss}}^{-1}\phi)$ and let $T_\vartheta$ be the lower-chain lift for twist $\vartheta$. For a fixed twist, \statementref{Theorem}{thm:a-seesaw} gives chain data $\Psi$ such that
\[
 T_\vartheta(\mathcal P\phi)=
 c_{\mathrm{ss}}\operatorname{Comp}_{\vartheta,\Psi}(\phi)
 \qquad (\phi\in U),
\]
where $\operatorname{Comp}_{\vartheta,\Psi}$ is the restriction of the upper-chain composite lift to the terminal target. Linearity of $T_\vartheta\circ\mathcal P$ gives
$T_\vartheta z=\operatorname{Comp}_{\vartheta,\Psi}(\phi)$. The projected translate is defined by $P_g^{\mathrm{proj}}\phi=|G|e_0gz$ and its normalization is $P_g\phi=c_B P_g^{\mathrm{proj}}\phi$.\leanfn{See \href{https://github.com/lsernau/grh-lean/blob/72fead102f663d5c7a407964a13dc0cf9d977207/grh/src/New/Reduction/ShortAndEnough.lean\#L168}{Normalization and the seesaw identity}.}

Use the set of primes $Q$ and local gains $\gamma_p$ of \statementref{Lemma}{lem:a-composite-size} and set
\[
 G_0=\sum_{p\in Q}\gamma_p\log p,\qquad
 \beta=H-mS+2G_0+B .
\]
\statementref{Lemma}{lem:a-composite-size} uses chain data for the trivial twist and a bound on source values. Given $\sigma$, if the conjugates of the point values of the source form have squared Weyl norms at most $e^u$, then those of the composite lift are at most $e^{u+\beta}$. Take $u=B$ to obtain
\[
 \|\sigma((T_1 z)(v))\|_{\mathrm{Weyl}}^2\le e^{B+\beta}
\]
at each finite adelic point $v$.\leanfn{See \href{https://github.com/lsernau/grh-lean/blob/72fead102f663d5c7a407964a13dc0cf9d977207/grh/src/New/UpperBound/Size/ArchimedeanSize.lean\#L426}{The size theorem for the composite lift}.}

After projection and translation, the point values have the expansion
\[
 \bigl(T_1(|G|e_0gz)\bigr)(v)
   = \sum_{t\in G}|G| e_0(t)\,(T_1z)(v\,a_1(tg)),
\]
Here $a_1(tg)$ is the first-block lift of group element $tg$ and the coefficients satisfy $|\sigma(e_0(t))|\le1$. The triangle inequality under $\sigma$ bounds the Weyl norm of each term by $|G|e^{(B+\beta)/2}$ and the norm of the sum by $|G|^2e^{(B+\beta)/2}$. The target squared norm is the class sum of squared pointwise norms, with each stabilizer weight at most $1$. Use the class number bound $\exp(x^{30})$ and the scale estimate $\exp(x^{30})|G|^4\le\exp(x^{99})$ to get\leanfn{See \href{https://github.com/lsernau/grh-lean/blob/72fead102f663d5c7a407964a13dc0cf9d977207/grh/src/New/UpperBound/Size/ArchimedeanSize.lean\#L475}{The projected translate formula}, \href{https://github.com/lsernau/grh-lean/blob/72fead102f663d5c7a407964a13dc0cf9d977207/grh/src/New/UpperBound/Size/ArchimedeanSize.lean\#L519}{A size bound for the projected translate}.}
\[
 \|\Theta_1 P_g^{\mathrm{proj}}\phi\|_\sigma^2
 \le \exp(x^{30})|G|^4 e^{B+\beta}
 \le \exp(x^{99}+B+\beta).
\]

\emph{Normalization and cancellation of the local gains.}
Fix the normalization witnesses $(e_p)_p$ given by \statementref{Lemma}{lem:a-uniform-divisibility}, common to every $P_g$. The errors satisfy
\[
 \sum_{p\in J}e_p\log p\le B
\]
for any finite set $J$ of primes. Put $Q'=Q\cup\{N\}$ and define
\[
 g_p=\lfloor\gamma_p-e_p\rfloor,\qquad c_B=\prod_{p\in Q'}p^{-g_p}.
\]
The number $c_B$ is positive rational. Since the level prime has $N\ge x_1$, we have $\gamma_N=0$ and hence $\sum_{p\in Q'}\gamma_p\log p=G_0$. The rounding inequality gives
\[
 \log c_B
 =-\sum_{p\in Q'}g_p\log p
 \le -G_0 +\sum_{p\in Q'}e_p\log p+\sum_{p\in Q'}\log p
 \le -G_0+B+b .
\]
The rounding error in the last step is bounded by
$\sum_{p\in Q'}\log p\le\lceil x_1\rceil x_1+\log N\le b$. The first term uses $|Q|\le\lceil x_1\rceil$ and $p<x_1$, and the last inequality is the scale bound for the level prime.\leanfn{See \href{https://github.com/lsernau/grh-lean/blob/72fead102f663d5c7a407964a13dc0cf9d977207/grh/src/New/Reduction/ShortAndEnough.lean\#L109}{Rounding error and the normalization bound}.}

Every automorphism $\sigma$ fixes the rational $c_B$, so multiplying by $c_B$ multiplies every conjugate squared norm by $c_B^2$. Combining the estimates gives
\[
 \|\Theta_1P_g\phi\|_\sigma^2
 \le \exp(x^{99}+B+\beta+2\log c_B)
 \le \exp(H-mS+x^{99}+4B+2b).
\]
Here $2G_0$ cancels the term $-2G_0$ in $2\log c_B$. The scales give $\ell\ge128$, $x\ge1$, and $x^{99}\le b$. Together with $a=xb$ and $f\ge1$ this yields $B=fa\ell^3\ge b$ and hence
\[
 x^{99}+4B+2b\le7B\le\ell B=L .
\]
Thus $\|\Theta_1P_g\phi\|_\sigma^2\le e^{H-mS+L}$ for each $\phi$ in $\mathcal S$ and every $g,\sigma$. Since $\mathcal S$ spans $U$, the proof is complete.
\end{proof}

\thmsource{Lean source: \leansourcelink{src/New/Reduction/ShortAndEnough.lean}{373}. See also \href{https://github.com/lsernau/grh-lean/blob/72fead102f663d5c7a407964a13dc0cf9d977207/grh/src/New/Reduction/ShortAndEnough.lean\#L342}{Local lattice index of one translated image}.}
\begin{theorem}[A(b): local lattice indices of normalized images]
\label{thm:index-bound}
\emitthmsource
Let $T\subset G$ be finite. If each member of a finite family $y$ belongs to $\bigcup_{g\in T}P_g(I)$, then $y$ satisfies
\begin{equation}
  \mathcal I_{2|T|\dim_{\mathbb C}(U)B}(y).
\label{eq:a-index}
\end{equation}
\end{theorem}
\begin{proof}
Let $U_0$ be the full source space, with $u_0=\dim_{\mathbb C}U_0$ and $u=\dim_{\mathbb C}U$, and keep $Q'$, $\gamma_p$, $e_p$, $g_p$, $c_B$, and $P_g^{\mathrm{proj}}$ from \statementlink{thm:short-forms}{Theorem~A(a)}. For a $p$-adic place $\iota$ write $\mathcal O_\iota$ for the valuation ring it induces on $\mathbb C$ and $W_{\iota}^{\mathrm{int}}\subseteq W$ for the integral target lattice. By definition, $y\in W_{\iota}^{\mathrm{int}}$ is equivalent to integrality of $\Theta_\vartheta y$ at $\iota$ for every $m$-fold twist $\vartheta$ with level norm coprime to $p$. The underlying twist of an $m$-fold twist has the same modulus norm, so the local divisibility lemma applies.\leanfn{See \href{https://github.com/lsernau/grh-lean/blob/72fead102f663d5c7a407964a13dc0cf9d977207/grh/src/New/Reduction/IntegralTargetForms.lean\#L65}{Definition of the integral target lattice}.}

\emph{A common correction matrix from local divisibility.}
Fix a $p$-adic place $\iota$. With the chosen loss family, \statementref{Lemma}{lem:a-uniform-divisibility} supplies a nonnegative integer $j_\iota\le e_pu_0$ and a local submodule $M_\iota$ of the integral ramified source space with two properties.
\begin{enumerate}
\item For any finite family $(\phi_i)_i$ from $I$ there is a matrix $A\in\operatorname{Mat}(\mathcal O_\iota)$ such that
$\det A\mid p^{j_\iota}$ and each $\sum_i A_{ki}\phi_i$ lies in $M_\iota$.
\item For $\phi\in M_\iota$, every admissible twist, and any chain data for it, the composite lift has $p$-adic Weyl gauge at most $p^{e_p-\gamma_p}$ at every finite adelic point.
\end{enumerate}
The same submodule $M_\iota$ works for every twist and set of chain data in the second property.\leanfn{See \href{https://github.com/lsernau/grh-lean/blob/72fead102f663d5c7a407964a13dc0cf9d977207/grh/src/New/UpperBound/Divisibility/LocalDivisibility.lean\#L482}{Local divisibility for the selected loss family}.}

\emph{Mapping the corrected source forms into the integral target lattice.}
When $p\in Q'$, the other prime factors of $c_B$ are units at $\iota$, so $|\iota(c_B)|_p=p^{g_p}$. When $p\notin Q'$, both gain and error are zero and hence $g_p=0$, so the same absolute-value formula holds. Combining this with $g_p\le\gamma_p-e_p$ shows that, for $\phi\in M_\iota$, the normalized composite lift has $p$-adic gauge at most
\[
 |\iota(c_B)|_p p^{e_p-\gamma_p}=p^{g_p+e_p-\gamma_p}\le1 .
\]
The seesaw identity transfers this bound to the normalized lower-chain lift. The projected translate formula in \statementlink{thm:short-forms}{Theorem~A(a)} also holds for arbitrary twists. Its coefficients $|G|e_0(t)$ are algebraic integers and the right translations are by elements of the standard integral level, so the local gauge of each term is at most $1$. The ultrametric inequality bounds the sum by $1$. Thus for every admissible $\vartheta$, the lift $\Theta_\vartheta P_g\phi$ is integral at $\iota$, which proves\leanfn{See \href{https://github.com/lsernau/grh-lean/blob/72fead102f663d5c7a407964a13dc0cf9d977207/grh/src/New/Reduction/ShortAndEnough.lean\#L151}{Local absolute value of the normalizing factor}, \href{https://github.com/lsernau/grh-lean/blob/72fead102f663d5c7a407964a13dc0cf9d977207/grh/src/New/UpperBound/Size/ArchimedeanSize.lean\#L490}{The projection preserves integrality}.}
\[
 P_g(M_\iota)\subseteq W_\iota^{\mathrm{int}}
 \qquad (g\in G).
\]

\emph{Bounding the logarithmic index for one translate.}
Fix $g$. For each $y_i\in P_g(I)$ in the family choose $\phi_i\in I$ with $P_g\phi_i=y_i$ and take the matrix $A$ from local divisibility. Linearity of $P_g$ gives, for each row index $k$,
\[
 \sum_i A_{ki}y_i=P_g\Bigl(\sum_i A_{ki}\phi_i\Bigr)\in W_\iota^{\mathrm{int}}.
\]
Divisibility of the determinant gives $v_p(\det A)\le j_\iota\le e_pu_0$. Thus the places above $p$ can use the same loss budget $e_pu_0$ even though the correction matrix may depend on the place. At each fixed place the matrix sends the vectors into $W_\iota^{\mathrm{int}}$ and works for every twist admissible there.

\statementref{Lemmas}{lem:a-ramified-dimension} and~\ref{lem:a-ramified-loss-scale} give $u\ge(1-f/(N-1))u_0$ and $f/(N-1)\le b^{-8}\le1/2$, hence $u_0\le2u$. For every finite set $J$ of primes,
\[
 \sum_{p\in J}(u_0e_p)\log p
 \le u_0B\le2uB.
\]
By the definition of $\mathcal I_\delta$, any finite family in a single translated image has index budget $2uB$.\leanfn{See \href{https://github.com/lsernau/grh-lean/blob/72fead102f663d5c7a407964a13dc0cf9d977207/grh/src/New/UpperBound/Size/ShortColumns.lean\#L942}{The codimension bound for ramification}, \href{https://github.com/lsernau/grh-lean/blob/72fead102f663d5c7a407964a13dc0cf9d977207/grh/src/New/Reduction/ShortAndEnough.lean\#L338}{The index estimate for one translate}.}

\emph{Combining finitely many translates.}
For each $y_i$ choose $g\in T$ with $y_i\in P_g(I)$ and group the indices by the choices. Fix $\iota$ and combine the correction matrices of the groups into a block diagonal matrix, restoring the index order if needed. All corrected vectors lie in $W_\iota^{\mathrm{int}}$. The determinant is the product of the block determinants, so the valuation is at most the sum of the group budgets. Summing the local losses over the same finite prime set $J$ gives total budget at most $|T|(2uB)$. For the empty family take the empty matrix with determinant $1$ and zero loss. This is the required $\mathcal I_{2|T|\dim_{\mathbb C}(U)B}(y)$.\leanfn{See \href{https://github.com/lsernau/grh-lean/blob/72fead102f663d5c7a407964a13dc0cf9d977207/grh/src/New/Reduction/Lattices/TwistedLattice.lean\#L205}{Combining local indices block by block}.}
\end{proof}
